% La protección del medioambiente y el cambio climático — Espagnol 1ère A — Cours signé OnBuch+
% Programme officiel MINESEC. Compiler avec : tectonic EA05-medioambiente-cambio-climatico.tex
\newcommand{\DOCMATIERE}{Espagnol}
\newcommand{\DOCNIVEAU}{1ère A}
\newcommand{\DOCMODULE}{Módulo 2 : Environnement, santé et bien-être}
\newcommand{\DOCLECON}{EA05}
\newcommand{\DOCTITRE}{La protección del medioambiente y el cambio climático}
% OnBuch+ — préambule « cours signés » ENRICHI (compilé avec tectonic / XeLaTeX)
% Système visuel « Pop » d'OnBuch+ : fond crème (jamais blanc), bordures encre
% épaisses, ombres « dures » décalées, titres Archivo Black en capitales.
% Version MATHÉMATIQUES : graphes (pgfplots/tikz), boîtes pédagogiques
% (définition, propriété, méthode, attention, le savais-tu, exercice résolu,
% exercice, TP, bilan), page de garde et signature OnBuch+ ancrée en bas.
%
% Macros attendues dans main.tex AVANT \input{preamble} :
%   \DOCMATIERE  \DOCNIVEAU  \DOCMODULE  \DOCLECON (numéro)  \DOCTITRE
\documentclass[11pt]{article}

\usepackage{fontspec}
\usepackage[spanish,french]{babel} % français = langue principale ; l'espagnol via \esp{..} / espagnol
\usepackage[a4paper,margin=2cm,top=3.4cm,bottom=2.8cm,headheight=1.4cm,footskip=1.3cm]{geometry}
\usepackage{amsmath,amssymb,amsfonts}
\usepackage{mathtools} % \xmapsto, \coloneqq... souvent utilisés par les modèles
\usepackage{cancel} % simplifications barrées (bilans de forces)
% \abs{z} (module, valeur absolue) et \norm{u} (norme), utilisés par les modèles
\providecommand{\abs}{}\let\abs\relax\DeclarePairedDelimiter{\abs}{\lvert}{\rvert}
\providecommand{\norm}{}\let\norm\relax\DeclarePairedDelimiter{\norm}{\lVert}{\rVert}
\usepackage[table]{xcolor} % [table] : fournit aussi \rowcolors (alternance de lignes)
\usepackage{pagecolor}
\usepackage{tikz}
\usetikzlibrary{arrows.meta,positioning,calc,shapes.geometric,shapes.misc,shapes.arrows,decorations.pathmorphing,decorations.pathreplacing,decorations.markings,patterns,shadows,mindmap,trees,fit,backgrounds,matrix,intersections,angles,quotes}
\usepackage{pgfplots}
\usepackage[version=4]{mhchem} % équations nucléaires \ce{^{235}_{92}U}
\usepackage[european,siunitx]{circuitikz} % schémas électriques
\pgfplotsset{compat=1.18}
\usepgfplotslibrary{fillbetween,groupplots} % aires entre courbes (calcul intégral)
\usepackage{enumitem}
\usepackage{fancyhdr}
% [most] charge toutes les bibliothèques tcolorbox (listings, minted, raster,
% documentation...) et gonfle inutilement la mémoire TeX sur un document long
% (~300 boîtes) : on ne charge que ce qui est réellement utilisé.
\usepackage[skins,breakable]{tcolorbox}
\usepackage{graphicx}
\usepackage{colortbl}
\usepackage{booktabs}
\usepackage{tabularx}
\usepackage{diagbox} % case d'en-tête diagonale des tableaux à double entrée
% Colonnes extensibles centrée / gauche / droite (C, L, R), souvent utilisées par les modèles
\newcolumntype{C}{>{\centering\arraybackslash}X}
\newcolumntype{L}{>{\raggedright\arraybackslash}X}
\newcolumntype{R}{>{\raggedleft\arraybackslash}X}
\usepackage{array}
\usepackage{multicol}
\usepackage{siunitx}
% parse-numbers=false : accepte aussi \SI{2,5 \times 10^{-3}}{...}
\sisetup{locale=FR,output-decimal-marker={,},group-minimum-digits=4,parse-numbers=false}
\DeclareSIUnit\ohm{\ensuremath{\symup{\Omega}}} % U+2126 absent de la police
\DeclareSIUnit\division{div} % réglages d'oscilloscope (ms/div, V/div)
\DeclareSIUnit\tr{tr} % tours (vitesse de rotation, ex. \SI{1500}{\tr\per\minute})
\DeclareSIUnit\molar{M} % concentration molaire (ex. \SI{3}{\molar}, \SI{1}{\micro\molar})
\DeclareSIUnit\an{an} % année (le modèle confond parfois avec \year, un compteur TeX) — ex. \SI{5}{\kilo\meter\per\an}
\DeclareSIUnit\FCFA{FCFA} % franc CFA (ex. \SI{500}{\FCFA\per\kilo\gram})
\DeclareSIUnit\degreeBrix{\textdegree Brix} % teneur en sucre d'un sirop/jus (réfractométrie)
\DeclareSIUnit\month{mois} % le modèle utilise parfois \month, un compteur TeX, comme unité de durée
\usepackage{qrcode}
% \degree : commande fréquemment utilisée par les modèles pour un angle ou une
% température en dehors de \SI{}{} (non fournie par les paquets chargés).
\providecommand{\degree}{^{\circ}}
% \qed (fin de démonstration), utilisé par les modèles sans amsthm
\providecommand{\qed}{\hfill$\square$}
% \barre : double barre des valeurs interdites dans un tableau de variations (tkz-tab)
\providecommand{\barre}{\ensuremath{\|}}
% environnement proof (amsthm non chargé) utilisé par les modèles
\ifdefined\proof\else\newenvironment{proof}[1][Démonstration]{\par\noindent\textit{#1.}\ }{\qed\par}\fi
\usepackage{lastpage}
\usepackage{hyperref}
\hypersetup{hidelinks}

% ============================================================
% Palette « Pop » OnBuch+ — mêmes hex que la classe P (plus_ui.dart)
% ============================================================
\definecolor{poporange}{HTML}{F59321}
\definecolor{popdark}{HTML}{E07A0C}
\definecolor{popcream}{HTML}{FFF6E8}
\definecolor{popink}{HTML}{1C1714}
\definecolor{poppurple}{HTML}{7A5AE0}
\definecolor{popgreen}{HTML}{2E9E6A}
\definecolor{poppink}{HTML}{E0457B}
\definecolor{popblue}{HTML}{2F7DE1}
\definecolor{popgold}{HTML}{C98A12}
\definecolor{popmuted}{HTML}{8A7F76}
\definecolor{popline}{HTML}{EADFD2}
\definecolor{popgray}{HTML}{8A7F76}
\colorlet{popgrey}{popgray}
% Alias tolérants (le modèle nomme parfois les couleurs différemment)
\colorlet{popwhite}{white}
\colorlet{popblack}{popink}
\colorlet{popred}{poppink}
\colorlet{popyellow}{popgold}
% Teintes claires (fonds de tableaux/encadrés) — le modèle nomme parfois
% les couleurs avec un suffixe L (ex. popblueL) sans les avoir définies.
\colorlet{poporangeL}{poporange!20}
\colorlet{popdarkL}{popdark!20}
\colorlet{poppurpleL}{poppurple!20}
\colorlet{popgreenL}{popgreen!20}
\colorlet{poppinkL}{poppink!20}
\colorlet{popblueL}{popblue!20}
\colorlet{popgoldL}{popgold!20}
\colorlet{popredL}{popred!20}
\colorlet{popgrayL}{popgray!20}
% Teintes claires pour fonds de tableaux / zones
\colorlet{popblueL}{popblue!10!white}
\colorlet{poporangeL}{poporange!14!white}
\colorlet{popgreenL}{popgreen!12!white}
\colorlet{poppurpleL}{poppurple!12!white}
\colorlet{poppinkL}{poppink!12!white}
\colorlet{popgoldL}{popgold!14!white}

\pagecolor{popcream}

% ============================================================
% Typographie
% ============================================================
\defaultfontfeatures{Ligatures=TeX}
\setmainfont{PlusJakartaSans-Medium.ttf}[Path=../../../fonts/,BoldFont=PlusJakartaSans-Bold.ttf]
% Maths en sans-serif (Fira Math) avec chiffres et lettres droites de Plus
% Jakarta, pour que les formules se fondent dans le texte.
\usepackage[math-style=ISO,bold-style=ISO]{unicode-math}
\setmathfont{FiraMath-Regular.otf}[Path=../../../fonts/]
\setmathfont{PlusJakartaSans-Medium.ttf}[Path=../../../fonts/,range={up/num,up/{Latin,latin},\mathup}]
\usepackage{icomma} % virgule décimale sans espace en mode math
% Caractères mathématiques Unicode absents des polices de texte (Plus Jakarta,
% Archivo Black) : rendus par leur équivalent mathématique, en texte comme en
% formule — sinon ils s'affichent en carré vide (ex. « Arithmétique dans ℤ »).
\usepackage{newunicodechar}
\newunicodechar{ℝ}{\ensuremath{\mathbb{R}}}
\newunicodechar{ℤ}{\ensuremath{\mathbb{Z}}}
\newunicodechar{ℂ}{\ensuremath{\mathbb{C}}}
\newunicodechar{ℕ}{\ensuremath{\mathbb{N}}}
\newunicodechar{ℚ}{\ensuremath{\mathbb{Q}}}
\newunicodechar{𝔻}{\ensuremath{\mathbb{D}}}
\newunicodechar{α}{\ensuremath{\alpha}}
\newunicodechar{β}{\ensuremath{\beta}}
\newunicodechar{γ}{\ensuremath{\gamma}}
\newunicodechar{δ}{\ensuremath{\delta}}
\newunicodechar{ε}{\ensuremath{\varepsilon}}
\newunicodechar{θ}{\ensuremath{\theta}}
\newunicodechar{λ}{\ensuremath{\lambda}}
\newunicodechar{μ}{\ensuremath{\mu}}
\newunicodechar{σ}{\ensuremath{\sigma}}
\newunicodechar{τ}{\ensuremath{\tau}}
\newunicodechar{φ}{\ensuremath{\varphi}}
\newunicodechar{ψ}{\ensuremath{\psi}}
\newunicodechar{ω}{\ensuremath{\omega}}
\newunicodechar{Σ}{\ensuremath{\Sigma}}
% majuscules produites par \MakeUppercase dans les titres (α-aminés -> Α)
\newunicodechar{Α}{\ensuremath{\alpha}}
\newunicodechar{Β}{\ensuremath{\beta}}
\newunicodechar{Γ}{\ensuremath{\Gamma}}
\newunicodechar{Δ}{\ensuremath{\Delta}}
\newunicodechar{Ε}{\ensuremath{\varepsilon}}
\newunicodechar{Θ}{\ensuremath{\Theta}}
\newunicodechar{Λ}{\ensuremath{\Lambda}}
\newunicodechar{Μ}{\ensuremath{\mu}}
\newunicodechar{Τ}{\ensuremath{\tau}}
\newunicodechar{Φ}{\ensuremath{\Phi}}
\newunicodechar{Ψ}{\ensuremath{\Psi}}
\newunicodechar{Ω}{\ensuremath{\Omega}}
\newunicodechar{↦}{\ensuremath{\mapsto}}
\newunicodechar{∈}{\ensuremath{\in}}
\newunicodechar{∉}{\ensuremath{\notin}}
\newunicodechar{∧}{\ensuremath{\wedge}}
\newunicodechar{⇔}{\ensuremath{\Leftrightarrow}}
\newunicodechar{⟺}{\ensuremath{\Longleftrightarrow}}
\newunicodechar{⇒}{\ensuremath{\Rightarrow}}
\newunicodechar{⟹}{\ensuremath{\Longrightarrow}}
\newunicodechar{∘}{\ensuremath{\circ}}
\newunicodechar{∪}{\ensuremath{\cup}}
\newunicodechar{∩}{\ensuremath{\cap}}
\newunicodechar{≡}{\ensuremath{\equiv}}
\newunicodechar{⊂}{\ensuremath{\subset}}
\newunicodechar{⊥}{\ensuremath{\perp}}
\newunicodechar{∃}{\ensuremath{\exists}}
\newunicodechar{∀}{\ensuremath{\forall}}
\newunicodechar{∛}{\ensuremath{\sqrt[3]{\,}}}
\newunicodechar{∜}{\ensuremath{\sqrt[4]{\,}}}
\newunicodechar{⃗}{\ensuremath{{}^{\to}}}
\newunicodechar{ⁿ}{\ifmmode^{n}\else\textsuperscript{\ensuremath{n}}\fi}
\newunicodechar{ˣ}{\ifmmode^{x}\else\textsuperscript{\ensuremath{x}}\fi}
\newunicodechar{ᵇ}{\ifmmode^{b}\else\textsuperscript{\ensuremath{b}}\fi}
\newunicodechar{ʳ}{\ifmmode^{r}\else\textsuperscript{\ensuremath{r}}\fi}
\newunicodechar{ᵘ}{\ifmmode^{u}\else\textsuperscript{\ensuremath{u}}\fi}
\newunicodechar{ʸ}{\ifmmode^{y}\else\textsuperscript{\ensuremath{y}}\fi}
\newunicodechar{ᵃ}{\ifmmode^{a}\else\textsuperscript{\ensuremath{a}}\fi}
\newunicodechar{ᵉ}{\ifmmode^{e}\else\textsuperscript{\ensuremath{e}}\fi}
\newunicodechar{ᵏ}{\ifmmode^{k}\else\textsuperscript{\ensuremath{k}}\fi}
\newunicodechar{ᵖ}{\ifmmode^{p}\else\textsuperscript{\ensuremath{p}}\fi}
\newunicodechar{ᴺ}{\ifmmode^{N}\else\textsuperscript{\ensuremath{N}}\fi}
\newunicodechar{ᶜ}{\ifmmode^{c}\else\textsuperscript{\ensuremath{c}}\fi}
\newunicodechar{ʰ}{\ifmmode^{h}\else\textsuperscript{\ensuremath{h}}\fi}
\newunicodechar{ᵠ}{\ifmmode^{q}\else\textsuperscript{\ensuremath{q}}\fi}
\newunicodechar{ₙ}{\ifmmode_{n}\else\textsubscript{\ensuremath{n}}\fi}
\newunicodechar{ₐ}{\ifmmode_{a}\else\textsubscript{\ensuremath{a}}\fi}
\newunicodechar{ᵢ}{\ifmmode_{i}\else\textsubscript{\ensuremath{i}}\fi}
\newunicodechar{ᵣ}{\ifmmode_{r}\else\textsubscript{\ensuremath{r}}\fi}
\newunicodechar{ₖ}{\ifmmode_{k}\else\textsubscript{\ensuremath{k}}\fi}
\newunicodechar{ᵦ}{\ifmmode_{\beta}\else\textsubscript{\ensuremath{\beta}}\fi}
\newunicodechar{ₓ}{\ifmmode_{x}\else\textsubscript{\ensuremath{x}}\fi}
\newfontfamily\popdisplay{ArchivoBlack-Regular.ttf}[Path=../../../fonts/]
\newfontfamily\poplogo{SpaceGrotesk-Bold.ttf}[Path=../../../fonts/]
\newfontfamily\popsemi{PlusJakartaSans-SemiBold.ttf}[Path=../../../fonts/]
\newfontfamily\popxbold{PlusJakartaSans-ExtraBold.ttf}[Path=../../../fonts/]
\newfontfamily\popxboldls{PlusJakartaSans-ExtraBold.ttf}[Path=../../../fonts/,LetterSpace=3.0]

\setlist[enumerate]{leftmargin=1.7em,itemsep=5pt,topsep=4pt}
\setlist[itemize]{leftmargin=1.7em,itemsep=4pt,topsep=4pt,label={\color{poporange}\textbullet}}
\setlength{\parindent}{0pt}
\setlength{\parskip}{5pt}
\renewcommand{\arraystretch}{1.25}

% pgfplots : style Pop par défaut
\pgfplotsset{
  every axis/.append style={axis line style={popink,line width=1pt},tick style={popink},
    grid style={popline,line width=0.5pt},label style={font=\small},tick label style={font=\footnotesize}},
  popcurve/.style={poporange,line width=1.8pt,no markers,smooth},
}
% Même style disponible dans un \draw TikZ simple (hors pgfplots)
\tikzset{popcurve/.style={poporange,line width=1.8pt}}
\tikzset{popgrid/.style={popline,very thin}}
\colorlet{popcurve}{poporange} % le nom du style est parfois utilisé comme couleur

% ============================================================
% Pill / squiggle
% ============================================================
\newcommand{\poppill}[3]{%
  \tikz[baseline=-0.55ex]\node[draw=popink,line width=1.2pt,rounded corners=2.6pt,%
    fill=#2,inner xsep=6.5pt,inner ysep=3pt]{%
    \popxboldls\fontsize{7}{7}\selectfont\color{#3}\MakeUppercase{#1}};%
}
\newcommand{\popsquiggle}[1]{%
  \tikz[baseline=-0.4ex]{
    \draw[#1,line width=2.6pt,line cap=round]
      plot [smooth] coordinates {(0,0.09) (0.34,0.01) (0.68,0.21) (1.02,0.01) (1.36,0.10)};
  }%
}
% Mot-clé surligné (marqueur orange sous le texte)
\newcommand{\cle}[1]{{\popxbold\color{popdark}#1}}

% ============================================================
% En-tête / pied
% ============================================================
\pagestyle{fancy}
\fancyhf{}
\renewcommand{\headrulewidth}{0pt}
\renewcommand{\footrulewidth}{0pt}
\lhead{\raisebox{-0.15ex}{\poplogo\fontsize{13}{13}\selectfont\color{popink}OnBuch\color{poporange}+}}
\rhead{\poppill{\DOCMATIERE\ \textbullet\ \DOCNIVEAU\ \textbullet\ Leçon \DOCLECON}{white}{popink}}
\lfoot{\popxbold\fontsize{7.6}{7.6}\selectfont\color{popmuted}\MakeUppercase{Cours signé OnBuch+}\ \textbullet\ {\popsemi\DOCTITRE}}
\rfoot{\tikz[baseline=-0.6ex]\node[rounded corners=8.5pt,fill=popink,minimum height=17pt,minimum width=17pt,inner xsep=5pt,inner sep=0]{\popxbold\fontsize{8}{8}\selectfont\color{popcream}\thepage\,/\,\pageref*{LastPage}};}

% ============================================================
% Titres
% ============================================================
\newcommand{\chaptitre}[2]{%
  \begin{tcolorbox}[enhanced,colback=poporange,colframe=popink,boxrule=2.6pt,arc=11pt,
      drop shadow=popink,left=15pt,right=15pt,top=12pt,bottom=11pt,before skip=3pt,after skip=18pt]
    \poppill{#2}{popink}{popcream}\par\vspace{9pt}
    {\raggedright\hyphenpenalty=10000\exhyphenpenalty=10000\popdisplay\fontsize{22}{24}\selectfont\color{popcream}\MakeUppercase{#1}\par}
  \end{tcolorbox}
}
% Section numérotée : pastille orange avec le numéro + titre Archivo
\newcounter{popsec}
\newcommand{\coursec}[1]{%
  \stepcounter{popsec}\setcounter{popsub}{0}%
  \par\vspace{16pt}\noindent
  \begin{minipage}{\linewidth}
  \tikz[baseline=-0.7ex]\node[draw=popink,line width=1.6pt,fill=poporange,rounded corners=4pt,minimum size=22pt,inner sep=0]{\popdisplay\fontsize{12}{12}\selectfont\color{popcream}\thepopsec};%
  \hspace{8pt}{\popdisplay\fontsize{15}{17}\selectfont\color{popink}\MakeUppercase{#1}}\par
  \vspace{4pt}\noindent{\color{poporange}\rule{36pt}{3.6pt}}
  \end{minipage}\par\nopagebreak\vspace{8pt}
}
\newcounter{popsub}
\newcommand{\courssub}[1]{%
  \stepcounter{popsub}%
  \par\vspace{9pt}\noindent{\popxbold\fontsize{11.5}{13}\selectfont\color{popink}{\color{poporange}\thepopsec.\thepopsub}\hspace{6pt}#1}\par\nopagebreak\vspace{4pt}
}

% ============================================================
% Boîtes pédagogiques — même recette : fond blanc, bordure épaisse colorée,
% ombre dure encre, badge plein. Titre optionnel [#1].
% ============================================================
\tcbset{popbase/.style={breakable,enhanced,colback=white,boxrule=2.3pt,arc=9pt,
  shadow={3pt}{-3pt}{0pt}{popink},left=14pt,right=14pt,top=11pt,bottom=11pt,before skip=11pt,after skip=15pt,
  fontupper=\popsemi\fontsize{10.6}{14.5}\selectfont\color{popink},
  fontlower=\popsemi\fontsize{10.6}{14.5}\selectfont\color{popink}}}
% Coche dessinée (le glyphe ✓ n'existe pas dans les polices du document)
\newcommand{\popcheck}[1]{\tikz[baseline=-0.5ex]\draw[#1,line width=1.6pt,line cap=round,line join=round](-0.13,0.0)--(-0.04,-0.09)--(0.13,0.1);}
\providecommand{\coche}{\popcheck{popgreen}}
\providecommand{\resultat}[1]{\par\smallskip\noindent\fcolorbox{popgreen}{popgreenL}{\parbox{\dimexpr\linewidth-2\fboxsep-2\fboxrule\relax}{\textbf{Résultat.} #1}}\par\smallskip}
\newcommand{\popboxhead}[3]{\poppill{#1}{#2}{white}\ifx\relax#3\relax\else\hspace{6pt}{\popxbold\fontsize{10.6}{12}\selectfont\color{popink}#3}\fi\par\vspace{7pt}}

\newtcolorbox{aretenir}[1][]{popbase,colframe=popgreen,before upper={\popboxhead{À retenir}{popgreen}{#1}}}
% Texte en espagnol (document de lecture, dialogue, poème, extrait) : cadre
% d'étude. Un titre optionnel (source, auteur) s'affiche dans l'en-tête.
\newtcolorbox{textebox}[1][]{popbase,colframe=popblue,before upper={\popboxhead{Texto}{popblue}{#1}\begin{otherlanguage}{spanish}},after upper={\end{otherlanguage}}}
% Marques ✓ / ✗ (forme correcte / fautive) souvent utilisées par les modèles
\providecommand{\cross}{\textcolor{poppink}{\ensuremath{\times}}}
\providecommand{\xmark}{\cross}\providecommand{\crossmark}{\cross}\providecommand{\cmark}{\ensuremath{\checkmark}}
% Ligne de réponse / justification à compléter (inventée par les modèles)
\providecommand{\justification}[1]{\par\noindent\textit{Justification~:} \dotfill\par}
% \textipa (package tipa non chargé) : la police ne couvre pas l'API -> simple passage
\providecommand{\textipa}[1]{#1}
% Soulignement (ulem non chargé) : \uline, \ul
\providecommand{\uline}[1]{\underline{#1}}\providecommand{\ul}[1]{\underline{#1}}
% Mot ou phrase en espagnol à l'intérieur d'un paragraphe français
\newcommand{\esp}[1]{\ifmmode\text{\textit{#1}}\else\begingroup\selectlanguage{spanish}\itshape #1\endgroup\fi}
\newtcolorbox{exemplebox}[1][]{popbase,colframe=poporange,before upper={\popboxhead{Exemple}{poporange}{#1}}}
\newtcolorbox{definition}[1][]{popbase,colframe=popblue,before upper={\popboxhead{Définition}{popblue}{#1}}}
\newtcolorbox{propriete}[1][]{popbase,colframe=poppurple,before upper={\popboxhead{Propriété}{poppurple}{#1}}}
\newtcolorbox{methode}[1][]{popbase,colframe=popgold,before upper={\popboxhead{Méthode}{popgold}{#1}}}
\newtcolorbox{attention}[1][]{popbase,colframe=poppink,before upper={\popboxhead{Attention}{poppink}{#1}}}
\newtcolorbox{experience}[1][]{popbase,colframe=popblue,colback=popblueL,before upper={\popboxhead{Expérience}{popblue}{#1}}}
% « Le savais-tu ? » : carte violette pleine (contexte, culture, Cameroun)
\newtcolorbox{savaistu}[1][]{popbase,colback=poppurple,colframe=popink,
  fontupper=\popsemi\fontsize{10.4}{14.2}\selectfont\color{white},
  before upper={\poppill{Le savais-tu\,?}{white}{poppurple}\ifx\relax#1\relax\else\hspace{6pt}{\popxbold\color{white}#1}\fi\par\vspace{7pt}}}
% Exercice résolu : énoncé en haut, \tcblower puis la solution sur fond crème
\newcounter{popexo}\newcounter{popexr}
\newtcolorbox{exoresolu}[1][]{popbase,colframe=poporange,colbacklower=popcream,
  segmentation style={popink,line width=1.2pt,dashed},
  before upper={\stepcounter{popexr}\popboxhead{Exercice résolu \thepopexr}{poporange}{#1}},
  before lower={\poppill{Solution}{popink}{popcream}\par\vspace{6pt}}}
% Exercice (non résolu en place) — la correction est dans la partie Corrigés
\newtcolorbox{exercice}[1][]{popbase,colframe=popink,
  before upper={\stepcounter{popexo}\popboxhead{Exercice \thepopexo}{popink}{#1}}}
% Corrigé
\newtcolorbox{corrige}[1][]{popbase,colframe=popgreen,colback=popgreenL,
  before upper={\popboxhead{Corrigé}{popgreen}{#1}}}
% Objectifs / prérequis / bilan
\newtcolorbox{objectifs}{popbase,colframe=popink,colback=poporangeL,
  before upper={\popboxhead{Objectifs de la leçon}{popink}{}}}
\newtcolorbox{prerequis}{popbase,colframe=popink,colback=popblueL,
  before upper={\popboxhead{Prérequis}{popblue}{}}}
\newtcolorbox{bilan}{popbase,colframe=popink,colback=white,boxrule=2.8pt,
  before upper={\popboxhead{Fiche bilan}{poporange}{L'essentiel à connaître pour l'examen}}}
% Figure encadrée avec légende
\newtcolorbox{popfigure}[1][]{enhanced,breakable=false,colback=white,colframe=popline,boxrule=1.4pt,arc=8pt,
  left=10pt,right=10pt,top=10pt,bottom=8pt,before skip=10pt,after skip=12pt,halign=center,
  lower separated=false,halign lower=center,
  fontlower=\popsemi\fontsize{9}{11.5}\selectfont\color{popmuted},
  #1}
\newcommand{\legende}[1]{\par\vspace{6pt}{\popsemi\fontsize{9}{11.5}\selectfont\color{popmuted}#1}}

% ============================================================
% Signature OnBuch+
% ============================================================
\newcommand{\poplogoinline}[1]{{\poplogo\fontsize{#1}{#1}\selectfont\color{popink}OnBuch\color{poporange}+}}
\newcommand{\signatureonbuch}{%
  \begin{tcolorbox}[enhanced,colback=popink,colframe=popink,boxrule=2.2pt,arc=10pt,
      drop shadow=poporange,left=14pt,right=14pt,top=10pt,bottom=10pt,before skip=0pt,after skip=0pt]
    \tikz[baseline=-0.7ex]\node[circle,fill=popgreen,draw=popcream,line width=1.3pt,minimum size=22pt,inner sep=0]{\popcheck{white}};%
    \hspace{9pt}\begin{minipage}[c]{0.62\linewidth}
      {\popxbold\fontsize{12}{13}\selectfont\color{popcream}Signé\ \textbullet\ {\poplogo OnBuch\color{poporange}+}}\par\vspace{2pt}
      {\raggedright\popsemi\fontsize{8}{10.5}\selectfont\color{popline}\MakeUppercase{Équipe pédagogique OnBuch+}\\\MakeUppercase{Conforme au programme officiel MINESEC}\par}
    \end{minipage}\hfill
    {\poplogo\fontsize{22}{22}\selectfont\color{popcream}OnBuch\color{poporange}+}
  \end{tcolorbox}%
}

% ============================================================
% Page de garde
% #1 titre · #2 matière · #3 niveau · #4 module · #5 n° leçon · #6 durée ·
% #7 description / objectifs (texte ou itemize)
% ============================================================
\newcommand{\pagegarde}[7]{%
  \thispagestyle{empty}
  \newgeometry{a4paper,margin=1.7cm,top=1.6cm,bottom=1.4cm}%
  \begin{tikzpicture}[remember picture,overlay]
    \fill[poporange!18] ([xshift=-3.2cm,yshift=-3.6cm]current page.north east) circle (4.6cm);
    \fill[poppurple!14] ([xshift=2.4cm,yshift=7.6cm]current page.south west) circle (3.8cm);
  \end{tikzpicture}%
  {\poplogo\fontsize{30}{30}\selectfont\color{popink}OnBuch\color{poporange}+}\hfill
  \raisebox{0.6ex}{\poppill{Cours signé \textbullet\ Premium}{popink}{popcream}}\par
  \vspace{1pt}\popsquiggle{poppurple}\par
  \vspace{8pt}
  \begin{tcolorbox}[enhanced,colback=poporange,colframe=popink,boxrule=2.8pt,arc=13pt,
      drop shadow=popink,left=18pt,right=18pt,top=12pt,bottom=13pt,after skip=10pt]
    \poppill{#2 \textbullet\ #3}{popink}{popcream}\hspace{5pt}\poppill{#4}{white}{popink}\par\vspace{8pt}
    {\popdisplay\fontsize{14}{14}\selectfont\color{popink}LEÇON #5}\par\vspace{4pt}
    {\raggedright\hyphenpenalty=10000\exhyphenpenalty=10000\popdisplay\fontsize{25}{27}\selectfont\color{popcream}\MakeUppercase{#1}\par}
    \vspace{8pt}
    {\popxbold\fontsize{9.5}{9.5}\selectfont\color{popink}\MakeUppercase{Durée conseillée : #6}}
  \end{tcolorbox}
  \begin{tcolorbox}[enhanced,colback=white,colframe=popink,boxrule=2.2pt,arc=10pt,
      drop shadow=popink,left=16pt,right=16pt,top=10pt,bottom=10pt,after skip=8pt]
    \poppill{Dans cette leçon}{poppurple}{white}\par\vspace{5pt}
    \popsemi\fontsize{9.3}{12.6}\selectfont\color{popink}#7
  \end{tcolorbox}
  \noindent\begin{minipage}[t]{0.487\linewidth}
    \begin{tcolorbox}[enhanced,colback=popgreen,colframe=popink,boxrule=2.2pt,arc=10pt,
        drop shadow=popink,left=11pt,right=11pt,top=10pt,bottom=10pt,height=3.1cm,valign=center]
      \tcbox[colback=white,colframe=popink,boxrule=1.3pt,arc=5pt,boxsep=0pt,nobeforeafter,
        left=3pt,right=3pt,top=3pt,bottom=3pt,valign=center]{\qrcode[height=1.9cm]{https://play.google.com/store/apps/details?id=com.luvvix.onbuch&hl=fr}}%
      \hspace{8pt}\begin{minipage}[c]{0.5\linewidth}
        {\popdisplay\fontsize{11}{12.5}\selectfont\color{white}\MakeUppercase{Télécharge OnBuch}}\par\vspace{4pt}
        {\raggedright\popsemi\fontsize{8.4}{11}\selectfont\color{white}Gratuit \textbullet\ Google Play\\Tuteur IA, annales, concours, résultats\par}
      \end{minipage}
    \end{tcolorbox}
  \end{minipage}\hfill
  \begin{minipage}[t]{0.487\linewidth}
    \begin{tcolorbox}[enhanced,colback=poppurple,colframe=popink,boxrule=2.2pt,arc=10pt,
        drop shadow=popink,left=11pt,right=11pt,top=10pt,bottom=10pt,height=3.1cm,valign=center]
      \tikz[baseline=-0.7ex]\node[circle,fill=white,draw=popink,line width=1.3pt,minimum size=1.6cm,inner sep=0]{\popdisplay\fontsize{24}{24}\selectfont\color{poppurple}+};%
      \hspace{8pt}\begin{minipage}[c]{0.6\linewidth}
        {\popdisplay\fontsize{11}{12.5}\selectfont\color{white}\MakeUppercase{Passe à OnBuch+}}\par\vspace{4pt}
        {\raggedright\popsemi\fontsize{8.4}{11}\selectfont\color{white}Tous les cours signés, Tuteur IA illimité, fiches PDF — onglet OnBuch+ de l'app\par}
      \end{minipage}
    \end{tcolorbox}
  \end{minipage}
  \vspace{8pt}
  \popcontenu
  \vfill
  \signatureonbuch
  \restoregeometry
  \newpage
}

% Rangée « Ce cours contient » (page de garde)
\newcommand{\poptile}[3]{% #1 couleur #2 gros texte #3 légende
  \begin{tcolorbox}[enhanced,colback=#1,colframe=popink,boxrule=2pt,arc=9pt,shadow={2.5pt}{-2.5pt}{0pt}{popink},
      left=6pt,right=6pt,top=9pt,bottom=9pt,halign=center,valign=center,height=1.75cm,width=\linewidth,nobeforeafter]
    {\popdisplay\fontsize{12.5}{13}\selectfont\color{white}\MakeUppercase{#2}}\par\vspace{3pt}
    {\centering\popsemi\fontsize{7.8}{9.5}\selectfont\color{white}#3\par}
  \end{tcolorbox}}
\newcommand{\popcontenu}{%
  \par\noindent{\popxboldls\fontsize{7.5}{7.5}\selectfont\color{popmuted}CE COURS SIGNÉ CONTIENT}\par\vspace{7pt}
  \noindent\begin{minipage}[t]{0.235\linewidth}\poptile{popblue}{Cours}{complet, illustré, pas à pas}\end{minipage}\hfill
  \begin{minipage}[t]{0.235\linewidth}\poptile{poporange}{Méthodes}{et exercices résolus}\end{minipage}\hfill
  \begin{minipage}[t]{0.235\linewidth}\poptile{poppink}{Exercices}{type examen + corrigés}\end{minipage}\hfill
  \begin{minipage}[t]{0.235\linewidth}\poptile{popgreen}{Fiche bilan}{l'essentiel à retenir}\end{minipage}\par}

% Sommaire
\newtcolorbox{sommaire}{enhanced,colback=white,colframe=popink,boxrule=2.2pt,arc=10pt,
  drop shadow=popink,left=18pt,right=18pt,top=13pt,bottom=13pt,before skip=6pt,after skip=20pt,
  before upper={\poppill{Sommaire}{poppurple}{white}\par\vspace{9pt}\popsemi\fontsize{10.6}{14.5}\selectfont\color{popink}}}

% Signature de fin de cours — poussée tout en bas de la dernière page
\newcommand{\findecours}{%
  \par\vfill\nopagebreak
  \begin{center}\popsquiggle{poporange}\end{center}\vspace{4pt}
  \signatureonbuch
}
\AtBeginDocument{\def\llbracket{\mathopen{[\mkern-3.5mu[}}\def\rrbracket{\mathclose{]\mkern-3.5mu]}}} % intervalles d'entiers (glyphe ⟦ absent de la police)
% Glyphes absents de Fira Math : redéfinis à partir de glyphes disponibles
\AtBeginDocument{%
  \def\setminus{\mathbin{\backslash}}\let\smallsetminus\setminus
  \def\vdots{\mathord{\vbox{\baselineskip3.2pt\lineskiplimit0pt\kern4pt\hbox{.}\hbox{.}\hbox{.}}}}%
  \def\ddots{\mathinner{\mkern1mu\raise6.5pt\hbox{.}\mkern2mu\raise3.6pt\hbox{.}\mkern2mu\raise0.7pt\hbox{.}\mkern1mu}}%
  \def\triangle{\mathord{\scalebox{0.9}{$\symup{\Delta}$}}}%
  \def\nabla{\mathord{\raisebox{\depth}{\scalebox{1}[-1]{$\symup{\Delta}$}}}}%
  \def\checkmark{\popcheck{popdark}}%
  \def\star{\mathbin{\ast}}\def\bigstar{\mathord{\ast}}%
  \def\diamond{\mathbin{\scalebox{0.6}{$\Diamond$}}}%
  \def\bigcup{\mathop{\mathchoice{\raisebox{-0.3em}{\scalebox{1.7}{$\cup$}}}{\scalebox{1.25}{$\cup$}}{\cup}{\cup}}\displaylimits}%
  \def\bigcap{\mathop{\mathchoice{\raisebox{-0.3em}{\scalebox{1.7}{$\cap$}}}{\scalebox{1.25}{$\cap$}}{\cap}{\cap}}\displaylimits}%
  \def\lesseqqgtr{\mathrel{\lessgtr}}\def\bigoplus{\mathop{\oplus}\displaylimits}%
}
\providecommand{\ssi}{si et seulement si} % abréviation fréquente
\AtBeginDocument{\providecommand{\wideparen}{\overparen}} % arc de cercle

%% FIGFIX-BEGIN v5 — correctifs globaux des figures (docs/onbuch-generation/tools/figfix)
\usepackage{adjustbox}\usepackage{etoolbox}\tcbuselibrary{raster}
% toute figure TikZ/pgfplots plus large que la ligne est réduite à la largeur disponible
\BeforeBeginEnvironment{tikzpicture}{\begin{adjustbox}{max width=\linewidth}}
\AfterEndEnvironment{tikzpicture}{\end{adjustbox}}
% légendes pgfplots lisibles par-dessus les courbes
\pgfplotsset{every axis/.append style={legend style={fill=white,fill opacity=0.92,text opacity=1,draw=black!15,inner sep=3pt}}}
% étiquettes d'arêtes (syntaxe edge["..."]) avec fond blanc
\tikzset{every edge quotes/.append style={fill=white,inner sep=1.2pt}}
%% FIGFIX-END
\begin{document}

\pagegarde{La protección del medioambiente y el cambio climático}{Espagnol}{1ère A}{Módulo 2 : Environnement, santé et bien-être}{EA05}{2 h}{Cette leçon propose la lecture et le commentaire d'un texte sur le changement climatique et la protection de l'environnement, avec des exemples camerounais (lac Tchad, forêt du bassin du Congo). Elle permet d'acquérir le lexique du climat et de maîtriser le futur simple et la phrase conditionnelle de conséquence (si + présent, futur) pour proposer des gestes écocitoyens.
\vspace{4pt}
\begin{multicols}{2}\small
\begin{itemize}
\item À la fin de cette leçon, je sais lire et comprendre un texte espagnol sur le changement climatique
\item À la fin de cette leçon, je sais identifier les causes et les conséquences du changement climatique dans un texte
\item À la fin de cette leçon, je sais employer le lexique du climat et de l'environnement (calentamiento global, sequía, inundación, deforestación, energías renovables, reciclar)
\item À la fin de cette leçon, je sais conjuguer et employer le futur simple pour exprimer des prédictions et des conséquences
\item À la fin de cette leçon, je sais construire une phrase de conséquence avec si + présent, futur
\item À la fin de cette leçon, je sais commenter un texte argumentatif en suivant une méthode claire
\end{itemize}\end{multicols}}

\begin{sommaire}
\begin{multicols}{2}
\begin{enumerate}
\item Texte support : El planeta en peligro
\item Compréhension du texte
\item Lexique thématique : el medioambiente
\item Grammaire : le futur simple
\item Grammaire : si + présent, futur (la conséquence)
\item Méthode du commentaire et production guidée
\item Activité d'intégration
\item Méthodes et exercices résolus
\item Exercices
\item Corrigés des exercices
\item Fiche bilan
\end{enumerate}
\end{multicols}
\end{sommaire}

\begin{objectifs}
\begin{itemize}
\item À la fin de cette leçon, je sais lire et comprendre un texte espagnol sur le changement climatique
\item À la fin de cette leçon, je sais identifier les causes et les conséquences du changement climatique dans un texte
\item À la fin de cette leçon, je sais employer le lexique du climat et de l'environnement (calentamiento global, sequía, inundación, deforestación, energías renovables, reciclar)
\item À la fin de cette leçon, je sais conjuguer et employer le futur simple pour exprimer des prédictions et des conséquences
\item À la fin de cette leçon, je sais construire une phrase de conséquence avec si + présent, futur
\item À la fin de cette leçon, je sais commenter un texte argumentatif en suivant une méthode claire
\item À la fin de cette leçon, je sais proposer des gestes concrets pour protéger l'environnement, à l'oral et à l'écrit
\item À la fin de cette leçon, je sais relier les enjeux climatiques mondiaux à des exemples camerounais (lac Tchad, forêt du Congo)
\end{itemize}
\end{objectifs}

\begin{prerequis}
\begin{itemize}
\item Le présent de l'indicatif des verbes réguliers et irréguliers courants (Seconde)
\item Le lexique de la nature et des saisons (Seconde)
\item Les structures d'opinion : pienso que, me parece que, es importante (Seconde)
\item La phrase hypothétique simple vue en début d'année (si + présent, présent)
\item Les connecteurs logiques de base : porque, por eso, además, sin embargo
\end{itemize}
\end{prerequis}

\newpage

\coursec{Texte support : El planeta en peligro}

\courssub{Situation d'accroche et problématique}

En classe, le professeur projette deux photos côte à côte sur l'écran : le lac Tchad en 1970, immense étendue bleue au milieu des terres, puis le même lac aujourd'hui, réduit à une petite mare entourée de sable. Sa surface a diminué de plus de 90 \%. Un élève lève la main et s'exclame : \esp{¿Qué ha pasado?} « Qu'est-ce qui s'est passé ? ».

Le même jour, de Yaoundé à Madrid en passant par Mexico, des milliers de jeunes manifestent pour le climat avec des pancartes portant le slogan \esp{No hay planeta B} « Il n'y a pas de planète B ».

Comment parler en espagnol du réchauffement climatique, de ses causes et de ses conséquences ? Et surtout, comment proposer des solutions pour protéger notre environnement ? C'est toute la problématique de cette leçon : \textbf{comprendre un texte sur l'environnement et le changement climatique}, afin de pouvoir ensuite lire, commenter et débattre sur ce sujet majeur de notre siècle.

\begin{savaistu}[Le deuxième poumon de la planète]
Le bassin du Congo, dont une partie se trouve au Cameroun (les régions du Sud et de l'Est), abrite la deuxième forêt tropicale du monde après l'Amazonie. Les médias hispanophones l'appellent \esp{el segundo pulmón del planeta} « le deuxième poumon de la planète », car ses arbres absorbent d'énormes quantités de dioxyde de carbone et produisent de l'oxygène. La protéger, c'est protéger le climat de toute la Terre. La Guinée équatoriale, voisine hispanophone du Cameroun, partage elle aussi cette forêt : le sujet de l'environnement est donc un excellent terrain de conversation avec les hispanophones d'Afrique centrale.
\end{savaistu}

\courssub{Lecture globale silencieuse du texte}

\begin{methode}[Comment lire un texte en espagnol : les 5 étapes]
Avant de lire le texte qui suit, rappelle-toi la bonne démarche de lecture :
\begin{enumerate}
\item \textbf{Lire le titre et anticiper} : le titre \esp{El planeta en peligro} annonce déjà le thème (la planète, un danger). Pose-toi la question : de quel danger s'agit-il ?
\item \textbf{Première lecture globale, sans dictionnaire} : lis le texte en entier, en silence, pour saisir l'idée générale. Ne t'arrête pas sur chaque mot inconnu.
\item \textbf{Repérer les mots transparents} : souligne les mots qui ressemblent au français et que tu comprends sans effort.
\item \textbf{Deuxième lecture avec le glossaire} : relis le texte en t'aidant du vocabulaire donné en marge.
\item \textbf{Reformuler l'idée principale} : en trois phrases, avec tes propres mots, dis ce que l'auteur veut nous faire comprendre.
\end{enumerate}
\end{methode}

Lis maintenant le texte en silence, sans dictionnaire. C'est un article de presse original, rédigé dans le style des journaux hispanophones.

\begin{textebox}[El planeta en peligro — artículo de prensa]
El calentamiento global es uno de los problemas más graves de nuestro siglo. La contaminación de las fábricas y de los coches, y la deforestación de los bosques, aumentan los gases de efecto invernadero. Por eso, el clima cambia: hay más sequías en el norte de Camerún y más inundaciones en las ciudades costeras. El lago Chad, por ejemplo, ha perdido más del noventa por ciento de su superficie. La selva del Congo, el segundo pulmón del planeta, también está en peligro. Sin embargo, todavía podemos actuar: si usamos energías renovables, reciclamos los residuos y plantamos árboles, protegeremos el medioambiente para las generaciones futuras.
\end{textebox}

\textbf{Traduction du texte :} « Le réchauffement climatique est l'un des problèmes les plus graves de notre siècle. La pollution des usines et des voitures, ainsi que la déforestation des forêts, augmentent les gaz à effet de serre. C'est pourquoi le climat change : il y a davantage de sécheresses dans le nord du Cameroun et davantage d'inondations dans les villes côtières. Le lac Tchad, par exemple, a perdu plus de quatre-vingt-dix pour cent de sa superficie. La forêt du Congo, le deuxième poumon de la planète, est également en danger. Cependant, nous pouvons encore agir : si nous utilisons des énergies renouvelables, si nous recyclons les déchets et si nous plantons des arbres, nous protégerons l'environnement pour les générations futures. »

\textbf{Glossaire en marge :}

\begin{center}
\begin{tabularx}{\linewidth}{|X|X|}
\hline
\rowcolor{popblueL} \textbf{Español} & \textbf{Français} \\
\hline
el calentamiento global & le réchauffement climatique \\
\hline
el cambio climático & le changement climatique \\
\hline
la sequía & la sécheresse \\
\hline
la inundación & l'inondation \\
\hline
la deforestación & la déforestation \\
\hline
las energías renovables & les énergies renouvelables \\
\hline
reciclar & recycler \\
\hline
el medioambiente & l'environnement \\
\hline
los residuos & les déchets \\
\hline
la fábrica & l'usine \\
\hline
el bosque / la selva & la forêt / la forêt tropicale \\
\hline
los gases de efecto invernadero & les gaz à effet de serre \\
\hline
la superficie & la superficie, la surface \\
\hline
costero, -a & côtier, côtière \\
\hline
sin embargo & cependant, pourtant \\
\hline
\end{tabularx}
\end{center}

\courssub{Lecture à voix haute et prononciation}

Le professeur lit maintenant le texte à voix haute. Écoute bien, puis répète après lui. Trois mots demandent une attention particulière :

\begin{propriete}[Prononciation : trois mots à travailler]
\begin{itemize}
\item \esp{sequía} se prononce « sé-ki-a » : l'accent écrit sur le \textbf{í} oblige à insister sur cette voyelle et à séparer les deux syllabes \emph{quí-a}. Sans accent, on prononcerait « sé-kia » en une seule syllabe.
\item \esp{energía} se prononce « é-nèr-rhi-a » : le \textbf{g} devant \textbf{e} se prononce comme le \textbf{j} espagnol, c'est-à-dire comme un « r » français très grasseyé, et l'accent tonique tombe sur le \textbf{í}.
\item \esp{protección} se prononce « pro-tèk-sion » : le \textbf{c} devant \textbf{i} se prononce « s » en Amérique latine et « th » anglais dans une grande partie de l'Espagne ; l'accent écrit sur le \textbf{ó} marque la syllabe tonique finale.
\end{itemize}
\end{propriete}

\begin{attention}[Les accents écrits changent le sens]
En espagnol, l'accent écrit n'est pas décoratif : \esp{proteccion} sans accent est une faute, et \esp{energia} sans accent se prononcerait « é-nèr-rhia » avec l'accent tonique au mauvais endroit. Retiens : tous les noms en \esp{-ción} portent un accent sur le \textbf{ó} : \esp{protección, contaminación, deforestación, inundación, solución}. Au pluriel, l'accent écrit disparaît : \esp{las inundaciones, las soluciones}.
\end{attention}

\courssub{Repérage des mots transparents}

Les \cle{mots transparents} sont des mots espagnols qui ressemblent tellement au français qu'on les comprend immédiatement. Ils sont tes meilleurs amis en lecture : ils te permettent de saisir le sens d'un texte sans dictionnaire.

\begin{exemplebox}[Les mots transparents du texte]
\begin{center}
\begin{tabularx}{\linewidth}{|X|X|X|}
\hline
\rowcolor{popblueL} \textbf{Español} & \textbf{Français} & \textbf{Remarque} \\
\hline
la protección & la protection & noms en \esp{-ción} = noms en \emph{-tion} \\
\hline
la contaminación & la pollution & même famille que « contamination » \\
\hline
la deforestación & la déforestation & préfixe \esp{des-/de-} = \emph{dé-} \\
\hline
renovables & renouvelables & adjectif en \esp{-able} = \emph{-able} \\
\hline
el problema & le problème & attention au genre : masculin ! \\
\hline
el planeta & la planète & masculin malgré le \textbf{-a} final \\
\hline
las generaciones futuras & les générations futures & presque identique \\
\hline
\end{tabularx}
\end{center}
\end{exemplebox}

\begin{attention}[Faux amis et pièges de genre]
\begin{itemize}
\item \esp{el problema}, \esp{el planeta}, \esp{el clima} sont \textbf{masculins} bien qu'ils finissent en \textbf{-a} : ce sont des mots d'origine grecque. On dit \esp{el problema es grave}, jamais « la problema ».
\item \esp{la contaminación} signifie « la pollution » au sens écologique ; le français « contamination » a un sens plus médical. Pour « pollution », l'espagnol utilise presque toujours \esp{contaminación} (ou \esp{polución} pour la pollution de l'air en ville).
\item \esp{los residuos} signifie « les déchets » ; ne dis pas « los restos », qui désigne plutôt les restes d'un repas.
\item \esp{actuar} signifie « agir » : \esp{todavía podemos actuar} = « nous pouvons encore agir ». Ce n'est pas « jouer un rôle » dans ce contexte.
\end{itemize}
\end{attention}

\courssub{Le mot-clé : medioambiente}

\begin{definition}[El calentamiento global]
\esp{El calentamiento global} désigne l'augmentation progressive de la température moyenne de la planète, provoquée principalement par l'accumulation des \esp{gases de efecto invernadero} (gaz à effet de serre) rejetés par les usines, les voitures et la déforestation. On parle aussi du \esp{cambio climático} (changement climatique) pour désigner l'ensemble des modifications du climat qui en résultent.
\end{definition}

\begin{attention}[Medioambiente : un seul mot]
En français, on dit « le milieu », « l'environnement ». En espagnol, le mot \esp{medioambiente} s'écrit de préférence \textbf{en un seul mot} : \esp{la protección del medioambiente}. La forme en deux mots \esp{el medio ambiente} est également acceptée par la Real Academia, mais jamais « el ambiente medio ». Attention aussi au genre : \esp{el medioambiente} est masculin. Et ne confonds pas \esp{el medioambiente} (la nature, l'environnement) avec \esp{el ambiente} (l'ambiance, l'atmosphère d'un lieu) : \esp{Hay buen ambiente en la clase} « Il y a une bonne ambiance dans la classe ».
\end{attention}

\courssub{Carte mentale et schéma : la logique du texte}

Le texte suit une logique en trois temps : les \textbf{causes}, les \textbf{conséquences}, les \textbf{solutions}. Cette structure est visible dans la carte mentale suivante.

\begin{popfigure}
\begin{tikzpicture}[mindmap, grow cyclic, every node/.style={concept, align=center, font=\small},
root concept/.append style={concept color=popblue, font=\small\bfseries},
level 1/.append style={level distance=3.2cm, sibling angle=120},
level 2/.append style={level distance=2.4cm, sibling angle=45, font=\scriptsize}]
\node[root concept]{Cambio climático}
  child[concept color=poporange]{ node {Causas}
      child { node {contaminación} }
      child { node {deforestación} }
      child { node {gases de efecto invernadero} }
  }
  child[concept color=poppink]{ node {Consecuencias}
      child { node {sequías} }
      child { node {inundaciones} }
      child { node {subida del mar} }
  }
  child[concept color=popgreen]{ node {Soluciones}
      child { node {energías renovables} }
      child { node {reciclar} }
      child { node {plantar árboles} }
  };
\end{tikzpicture}
\legende{Carte mentale du texte : causes, conséquences et solutions du changement climatique.}
\end{popfigure}

Le schéma suivant montre la chaîne de cause à effet : le dioxyde de carbone (CO\textsubscript{2}) provoque le réchauffement, qui provoque à son tour les catastrophes climatiques.

\begin{popfigure}
\begin{tikzpicture}[
nodo/.style={draw=popblue, fill=popblueL, rounded corners, align=center, font=\small, minimum height=0.9cm, inner sep=4pt},
flecha/.style={-{Stealth[length=3mm]}, thick, popdark}]
\node[nodo, fill=poporangeL, draw=poporange, align=center] (co2){CO\textsubscript{2}\\ gases};
\node[nodo, right=1.6cm of co2, align=center] (cal){calentamiento\\ global};
\node[nodo, right=1.6cm of cal, yshift=1.3cm] (seq) {sequía};
\node[nodo, right=1.6cm of cal] (inu) {inundación};
\node[nodo, right=1.6cm of cal, yshift=-1.3cm, align=center] (mar){subida\\ del mar};
\draw[flecha] (co2) -- (cal);
\draw[flecha] (cal) -- (seq);
\draw[flecha] (cal) -- (inu);
\draw[flecha] (cal) -- (mar);
\end{tikzpicture}
\legende{La chaîne de cause à effet : des gaz à effet de serre aux catastrophes climatiques.}
\end{popfigure}

\courssub{Genre, thème et thèse du texte}

Après la deuxième lecture, identifions ensemble les caractéristiques du texte.

\begin{propriete}[Identifier un article de presse]
Un \cle{article de presse} (\esp{artículo de prensa}) se reconnaît à plusieurs indices :
\begin{itemize}
\item un \textbf{titre} court et frappant : \esp{El planeta en peligro} ;
\item un ton \textbf{informatif} : l'auteur donne des faits (le lac Tchad a perdu plus de 90 \% de sa superficie) ;
\item un ton \textbf{argumentatif} : l'auteur veut convaincre (il faut agir, il propose des solutions) ;
\item un vocabulaire du \textbf{domaine public} : environnement, climat, générations futures.
\end{itemize}
\end{propriete}

Appliquons cela à notre texte :

\begin{itemize}
\item \textbf{Le genre} : c'est un article de presse, à la fois informatif et argumentatif.
\item \textbf{Le thème} : le changement climatique et la protection de l'environnement, dans le monde et au Cameroun.
\item \textbf{La thèse de l'auteur} : le réchauffement climatique est un problème grave dont l'homme est responsable, mais il est encore temps d'agir. La phrase clé est la dernière : \esp{todavía podemos actuar} « nous pouvons encore agir ». Le connecteur \esp{sin embargo} « cependant » marque le passage du constat pessimiste à l'espoir.
\end{itemize}

\begin{propriete}[Les connecteurs logiques du texte]
Pour suivre l'argumentation d'un texte, repère ses \cle{connecteurs logiques} :
\begin{itemize}
\item \textbf{la cause} : \esp{por eso} « c'est pourquoi » (introduit ici la conséquence des causes), \esp{porque} « parce que » ;
\item \textbf{l'opposition} : \esp{sin embargo} « cependant », \esp{pero} « mais » ;
\item \textbf{l'exemple} : \esp{por ejemplo} « par exemple » ;
\item \textbf{l'addition} : \esp{también} « aussi, également », \esp{y} « et ».
\end{itemize}
Ces mots-outils te seront indispensables pour tes propres rédactions à l'examen.
\end{propriete}

\courssub{Reformulation orale du texte}

Dernière étape de la lecture : reformuler le texte en trois phrases, avec tes propres mots, sans regarder le texte. C'est un excellent entraînement à l'expression orale et au résumé.

\begin{exoresolu}[Reformuler l'idée principale en trois phrases]
Énoncé : reformule le texte \esp{El planeta en peligro} en trois phrases simples, avec tes propres mots (une phrase pour les causes, une pour les conséquences, une pour les solutions).
\tcblower
Solution détaillée :
\begin{itemize}
\item \textbf{Causes} : \esp{El hombre contamina el aire y corta los bosques, y esto calienta el planeta.} « L'homme pollue l'air et coupe les forêts, et cela réchauffe la planète. »
\item \textbf{Conséquences} : \esp{Por eso hay sequías, inundaciones, y el lago Chad casi desaparece.} « C'est pourquoi il y a des sécheresses, des inondations, et le lac Tchad disparaît presque. »
\item \textbf{Solutions} : \esp{Si reciclamos y plantamos árboles, salvaremos el medioambiente.} « Si nous recyclons et plantons des arbres, nous sauverons l'environnement. »
\end{itemize}
Remarque la structure de la troisième phrase : \esp{si} + présent, puis futur. C'est la structure de la conséquence que nous étudierons en grammaire dans cette leçon.
\end{exoresolu}

\begin{exercice}[À toi de jouer]
\begin{enumerate}
\item Retrouve dans le texte trois mots transparents qui ne figurent pas dans le tableau de la leçon.
\item Vrai ou faux ? Justifie par une phrase du texte : \esp{El lago Chad ha perdido la mitad de su superficie.} « Le lac Tchad a perdu la moitié de sa superficie. »
\item Traduis en français : \esp{Si usamos energías renovables, protegeremos el medioambiente.}
\item Reformule la thèse de l'auteur en une seule phrase commençant par \esp{El autor piensa que...}
\end{enumerate}
\end{exercice}

\begin{corrige}[Exercice : À toi de jouer]
\begin{enumerate}
\item Par exemple : \esp{grave} (grave), \esp{superficie} (superficie), \esp{generaciones} (générations), \esp{futuras} (futures).
\item Faux. Le texte dit : \esp{ha perdido más del noventa por ciento de su superficie} « il a perdu plus de quatre-vingt-dix pour cent de sa superficie », et non la moitié (\esp{la mitad}).
\item « Si nous utilisons des énergies renouvelables, nous protégerons l'environnement. »
\item \esp{El autor piensa que el cambio climático es muy grave, pero que todavía podemos actuar.} « L'auteur pense que le changement climatique est très grave, mais que nous pouvons encore agir. »
\end{enumerate}
\end{corrige}

\begin{aretenir}[L'essentiel de cette section]
\begin{itemize}
\item Le texte \esp{El planeta en peligro} est un article de presse informatif et argumentatif sur le changement climatique.
\item Sa structure : \textbf{causes} (contaminación, deforestación, gases) $\rightarrow$ \textbf{conséquences} (sequías, inundaciones, subida del mar, lac Tchad, forêt du Congo) $\rightarrow$ \textbf{solutions} (energías renovables, reciclar, plantar árboles).
\item La bonne méthode de lecture : titre et anticipation, lecture globale, mots transparents, lecture avec glossaire, reformulation.
\item Vocabulaire-clé : \esp{calentamiento global, cambio climático, sequía, inundación, medioambiente, residuos, reciclar}.
\item Connecteurs à réutiliser : \esp{por eso, sin embargo, por ejemplo, también}.
\item Pièges : \esp{medioambiente} en un seul mot, \esp{el problema / el planeta} masculins, accents obligatoires sur les noms en \esp{-ción} au singulier.
\end{itemize}
\end{aretenir}

\coursec{Compréhension du texte}

Dans la section précédente, nous avons lu ensemble le texte \emph{El planeta en peligro}, un article original sur le changement climatique et la protection de l'environnement, avec des exemples pris au Cameroun et dans le bassin du Congo. Nous avons dégagé le thème général et repéré les mots nouveaux. Il est temps maintenant de vérifier que nous avons \textbf{vraiment compris} le texte : c'est l'étape de la \cle{comprensión del texto}, une compétence centrale à l'évaluation et au Baccalauréat. Pour travailler efficacement, relisons d'abord le texte support.

\begin{textebox}[El planeta en peligro (texte support)]
Actualmente, el cambio climático es uno de los problemas más graves del planeta. La contaminación del aire, de los ríos y de los mares, la deforestación y el uso excesivo de los coches aumentan los gases de efecto invernadero y provocan el calentamiento global.

Las consecuencias ya son visibles en muchos países. En el norte de Camerún, la sequía es cada vez más larga y el lago Chad ha perdido más del noventa por ciento de su superficie desde los años sesenta. Millones de personas que vivían de la pesca y de la agricultura han tenido que abandonar la región. En el sur, en cambio, las lluvias provocan a veces inundaciones que destruyen casas y cosechas.

La selva del Congo, el segundo bosque tropical más grande del mundo después del Amazonas, es muy importante para el planeta: sus árboles absorben el dióxido de carbono y producen oxígeno. Sin embargo, la deforestación avanza rápidamente por la tala de árboles y la agricultura intensiva.

¿Qué podemos hacer? Los gobiernos deben desarrollar las energías renovables, como la energía solar y la energía hidráulica. Pero los ciudadanos también pueden actuar: reciclar la basura, plantar árboles, ahorrar agua y electricidad, y utilizar el transporte público. Si todos colaboramos, protegeremos el medioambiente para las generaciones futuras. Los jóvenes tienen un papel esencial: ellos son el futuro del planeta.
\end{textebox}

\textbf{Glossaire :} \esp{la sequía} : la sécheresse ; \esp{la pesca} : la pêche ; \esp{la cosecha} : la récolte ; \esp{la selva} : la forêt tropicale ; \esp{la tala} : l'abattage (des arbres) ; \esp{ahorrar} : économiser ; \esp{la basura} : les ordures ; \esp{colaborar} : collaborer, participer ; \esp{el efecto invernadero} : l'effet de serre ; \esp{sin embargo} : cependant.

\begin{methode}[Répondre à une question de compréhension]
Au Bac comme en classe, une bonne réponse de compréhension respecte quatre étapes :
\begin{enumerate}
\item \textbf{Lire la question} et souligner les mots-clés (personne, lieu, phénomène demandé).
\item \textbf{Repérer dans le texte} la phrase ou le passage qui contient la réponse.
\item \textbf{Rédiger une phrase complète en espagnol}, en reprenant les mots de la question. On ne répond jamais par un seul mot.
\item \textbf{Justifier par une citation} du texte entre guillemets, introduite par \esp{según el texto} ou \esp{como dice el texto}.
\end{enumerate}
\end{methode}

\begin{propriete}[Les formules de justification]
Pour justifier une réponse, l'espagnol utilise des verbes de parole, comme le français (« selon le texte », « l'auteur affirme que ») :
\begin{itemize}
\item \esp{Según el texto, ...} « Selon le texte, ... »
\item \esp{El autor dice que ...} / \esp{El autor afirma que ...} « L'auteur dit / affirme que ... »
\item \esp{Como dice el texto: ``...''} « Comme le dit le texte : “...” »
\item \esp{El texto menciona que ...} « Le texte mentionne que ... »
\end{itemize}
Après \esp{que}, on utilise l'indicatif, car on rapporte un fait présenté comme certain.
\end{propriete}

\begin{propriete}[Les mots interrogatifs : bien comprendre la question]
Chaque mot interrogatif espagnol porte un \textbf{accent écrit} et appelle un type de réponse précis :
\begin{center}
\begin{tabularx}{\linewidth}{|X|X|X|}
\hline
\rowcolor{popblueL} \textbf{Mot interrogatif} & \textbf{Français} & \textbf{Type de réponse} \\
\hline
\esp{¿De qué...?} & De quoi... ? & le thème \\
\hline
\esp{¿Qué...?} & Que / Quel... ? & une chose, un phénomène \\
\hline
\esp{¿Quién...?} & Qui... ? & une personne \\
\hline
\esp{¿Dónde...?} & Où... ? & un lieu \\
\hline
\esp{¿Cuándo...?} & Quand... ? & un moment \\
\hline
\esp{¿Por qué...?} & Pourquoi... ? & une cause (\esp{porque...}) \\
\hline
\esp{¿Cuál / Cuáles...?} & Quel(le)(s)... ? & un choix parmi plusieurs \\
\hline
\esp{¿Cuánto/a/os/as...?} & Combien de... ? & une quantité \\
\hline
\end{tabularx}
\end{center}
Attention à ne pas confondre \esp{¿por qué?} (« pourquoi », en deux mots, avec accent) et \esp{porque} (« parce que », en un seul mot, sans accent).
\end{propriete}

\courssub{Questions de compréhension globale}

Ces questions portent sur le texte dans son ensemble. Ce sont toujours les premières questions posées à l'oral ou à l'écrit.

\begin{exoresolu}[Compréhension globale]
\textbf{1. ¿De qué habla el texto?} (De quoi parle le texte ?)

\textbf{2. ¿Cuál es la tesis del autor?} (Quelle est la thèse de l'auteur ?)
\tcblower
\textbf{1.} \esp{El texto habla del cambio climático y de la protección del medioambiente, con ejemplos de Camerún y de la selva del Congo.} « Le texte parle du changement climatique et de la protection de l'environnement, avec des exemples du Cameroun et de la forêt du Congo. » Remarque : on reprend les mots de la question (\esp{el texto habla de...}) et on répond par une phrase complète.

\textbf{2.} \esp{La tesis del autor es que el cambio climático es un problema muy grave, pero que todos podemos actuar para proteger el planeta. Según el texto: ``Si todos colaboramos, protegeremos el medioambiente''.} « La thèse de l'auteur est que le changement climatique est un problème très grave, mais que nous pouvons tous agir pour protéger la planète. »
\end{exoresolu}

\courssub{Questions sur les causes et les conséquences}

\textbf{Questions sur les causes}

\esp{¿Qué fenómenos aumentan los gases de efecto invernadero?} « Quels phénomènes augmentent les gaz à effet de serre ? »

Réponse attendue : \esp{Según el texto, la contaminación del aire, de los ríos y de los mares, la deforestación y el uso excesivo de los coches aumentan los gases de efecto invernadero.} « Selon le texte, la pollution de l'air, des rivières et des mers, la déforestation et l'usage excessif des voitures augmentent les gaz à effet de serre. »

\textbf{Questions sur les conséquences}

\esp{¿Qué consecuencias del cambio climático se mencionan en Camerún?} « Quelles conséquences du changement climatique sont mentionnées au Cameroun ? »

Réponse attendue : \esp{En el norte, la sequía es cada vez más larga y el lago Chad ha perdido más del noventa por ciento de su superficie; en el sur, las lluvias provocan inundaciones que destruyen casas y cosechas.} « Au nord, la sécheresse est de plus en plus longue et le lac Tchad a perdu plus de quatre-vingt-dix pour cent de sa superficie ; au sud, les pluies provoquent des inondations qui détruisent des maisons et des récoltes. »

\begin{exemplebox}[Exemples traduits]
\esp{El lago Chad ha perdido más del noventa por ciento de su superficie.} « Le lac Tchad a perdu plus de quatre-vingt-dix pour cent de sa superficie. »

\esp{Las inundaciones destruyen casas y cosechas.} « Les inondations détruisent des maisons et des récoltes. »

\esp{La sequía es cada vez más larga.} « La sécheresse est de plus en plus longue. »
\end{exemplebox}

\courssub{Questions sur les exemples et sur les solutions}

\textbf{Sur les exemples}

\esp{¿Qué ha pasado con el lago Chad?} « Que s'est-il passé avec le lac Tchad ? » — \esp{El lago Chad ha perdido más del noventa por ciento de su superficie desde los años sesenta, y millones de personas han tenido que abandonar la región.} « Le lac Tchad a perdu plus de 90 \% de sa superficie depuis les années soixante, et des millions de personnes ont dû quitter la région. »

\esp{¿Por qué es importante la selva del Congo?} « Pourquoi la forêt du Congo est-elle importante ? » — \esp{Es importante porque sus árboles absorben el dióxido de carbono y producen oxígeno; es el segundo bosque tropical más grande del mundo.} « Elle est importante parce que ses arbres absorbent le dioxyde de carbone et produisent de l'oxygène ; c'est la deuxième forêt tropicale du monde. »

\textbf{Sur les solutions}

\esp{¿Qué podemos hacer para proteger el medioambiente?} « Que pouvons-nous faire pour protéger l'environnement ? » — \esp{Podemos reciclar la basura, plantar árboles, ahorrar agua y electricidad y utilizar el transporte público. Los gobiernos deben desarrollar las energías renovables.} « Nous pouvons recycler les ordures, planter des arbres, économiser l'eau et l'électricité et utiliser les transports en commun. Les gouvernements doivent développer les énergies renouvelables. »

\begin{attention}[Bien employer \esp{para} et \esp{porque} dans les réponses]
\begin{itemize}
\item Pour exprimer le \textbf{but} (« pour protéger »), on utilise \esp{para + infinitif} : \esp{Actuamos para proteger el planeta.} « Nous agissons pour protéger la planète. » Ne dis jamais \esp{por proteger} dans ce sens.
\item Pour exprimer la \textbf{cause} (« parce que »), on utilise \esp{porque + verbe conjugué} : \esp{La selva es importante porque absorbe dióxido de carbono.}
\end{itemize}
\end{attention}

\courssub{Questions de réaction personnelle}

Ces questions n'ont pas de réponse unique : elles évaluent ta capacité à \textbf{t'exprimer} en espagnol. Utilise la première personne et des connecteurs simples : \esp{en mi opinión} « à mon avis », \esp{creo que} « je crois que », \esp{por eso} « c'est pourquoi », \esp{además} « de plus ».

\begin{exemplebox}[Modèles de réponses personnelles]
\esp{¿Y tú, qué haces para proteger el planeta?} — \esp{Yo reciclo las botellas de plástico, apago la luz cuando salgo de mi cuarto y planto árboles con mi familia en el pueblo.} « Moi, je recycle les bouteilles en plastique, j'éteins la lumière quand je sors de ma chambre et je plante des arbres avec ma famille au village. »

\esp{¿Crees que los jóvenes pueden cambiar las cosas?} — \esp{Sí, creo que los jóvenes pueden cambiar las cosas porque son el futuro del planeta y pueden sensibilizar a su familia y a sus amigos.} « Oui, je crois que les jeunes peuvent changer les choses parce qu'ils sont l'avenir de la planète et qu'ils peuvent sensibiliser leur famille et leurs amis. »
\end{exemplebox}

\courssub{Vrai ou faux : l'entraînement type Bac}

L'exercice \esp{verdadero o falso} est classique à l'évaluation : il faut répondre \esp{verdadero} ou \esp{falso} \textbf{et justifier par une citation exacte du texte}. Sans justification, la réponse ne vaut rien.

\begin{center}
\begin{tabularx}{\linewidth}{|X|X|}
\hline
\rowcolor{popblueL} \textbf{Afirmación} & \textbf{Justificación esperada} \\
\hline
1. El lago Chad ha perdido la mitad de su superficie. & Falso : « ha perdido más del noventa por ciento de su superficie » \\
\hline
2. La selva del Congo absorbe dióxido de carbono. & Verdadero : « sus árboles absorben el dióxido de carbono » \\
\hline
3. Solo los gobiernos pueden proteger el medioambiente. & Falso : « los ciudadanos también pueden actuar » \\
\hline
4. En el sur de Camerún hay inundaciones. & Verdadero : « en el sur, las lluvias provocan a veces inundaciones » \\
\hline
\end{tabularx}
\end{center}

\begin{attention}[Pièges de traduction et de compréhension]
\begin{itemize}
\item \textbf{Ne traduis pas mot à mot.} Comprends d'abord le sens global de la phrase, puis reformule en français naturel.
\item \esp{Actualmente} signifie « actuellement, à l'heure actuelle », et non « en fait ». De même, \esp{un problema actual} = « un problème actuel », pas « un problème réel ».
\item \esp{Los años sesenta} = « les années soixante » (et non « les ans soixante »).
\item \esp{El noventa por ciento} = « quatre-vingt-dix pour cent » : attention au calcul mental, \esp{noventa} = 90, pas 19 (\esp{diecinueve}).
\item \esp{Millones de personas} = « des millions de personnes » : \esp{millones de} se construit toujours avec la préposition \esp{de}.
\item \esp{En cambio} = « en revanche, par contre » ; ce n'est pas le nom \esp{un cambio} (« un changement »).
\end{itemize}
\end{attention}

\begin{savaistu}[Le lac Tchad, un drame écologique à nos portes]
Le lac Tchad, situé aux frontières du Tchad, du Cameroun, du Nigeria et du Niger, était autrefois l'un des plus grands lacs d'Afrique. En quelques décennies, il a perdu l'essentiel de sa superficie à cause de la sécheresse et des prélèvements d'eau pour l'agriculture. C'est l'un des exemples les plus cités en Afrique centrale pour parler du \esp{cambio climático}. En espagnol, on écrit \esp{el lago Chad} : le nom du pays, \esp{Chad}, ne prend pas d'accent et se prononce « tchad ».
\end{savaistu}

\courssub{Schéma récapitulatif de la compréhension}

Voici l'organigramme de lecture qui résume la démarche : du titre vers l'idée principale, puis les causes, les conséquences, les solutions, et enfin ton opinion personnelle.

\begin{popfigure}
\begin{center}
\begin{tikzpicture}[
  node distance=0.55cm,
  every node/.style={draw=popblue, thick, rounded corners=3pt, fill=popblueL, align=center, font=\small, inner sep=4pt},
  fleche/.style={-{Stealth[length=2.5mm]}, thick, popdark}
]
\node[fill=popgoldL, draw=popgold] (titulo) {Título : \emph{El planeta en peligro}};
\node[below=of titulo, align=center] (idea){Idea principal : el cambio climático\\es grave, pero podemos actuar};
\node[below left=0.55cm and -0.8cm of idea, align=center] (causas){Causas : contaminación,\\deforestación, coches};
\node[below right=0.55cm and -0.8cm of idea, align=center] (cons){Consecuencias : sequía,\\inundaciones, lago Chad};
\node[below=1.9cm of idea, align=center] (sol){Soluciones : energías renovables,\\reciclar, plantar árboles};
\node[below=of sol, fill=popgreenL, draw=popgreen, align=center] (op){Opinión personal :\\¿Y tú, qué haces?};
\draw[fleche] (titulo) -- (idea);
\draw[fleche] (idea) -- (causas);
\draw[fleche] (idea) -- (cons);
\draw[fleche] (causas) -- (sol);
\draw[fleche] (cons) -- (sol);
\draw[fleche] (sol) -- (op);
\end{tikzpicture}
\end{center}
\legende{Organigramme de compréhension du texte \emph{El planeta en peligro}}
\end{popfigure}

\begin{experience}[Entretien écologique en binômes]
Par groupes de deux, jouez la scène suivante : un journaliste de la radio de Yaoundé interviewe un élève sur l'environnement. Le journaliste pose au moins quatre questions avec les mots interrogatifs étudiés (\esp{¿qué?}, \esp{¿por qué?}, \esp{¿dónde?}, \esp{¿cuándo?}) ; l'élève répond par des phrases complètes en s'inspirant du texte. Exemple de départ : \esp{— ¿Qué haces para proteger el medioambiente? — Yo reciclo la basura y ahorro agua en casa.} Puis échangez les rôles.
\end{experience}

\courssub{Entraînement}

\begin{exercice}[Compréhension du texte]
Réponds par des phrases complètes en espagnol, avec une justification quand elle est demandée.

1. ¿De qué habla el texto?

2. ¿Qué fenómenos aumentan los gases de efecto invernadero?

3. ¿Qué consecuencias del cambio climático se mencionan en el norte de Camerún?

4. ¿Por qué es importante la selva del Congo?

5. ¿Qué pueden hacer los ciudadanos para proteger el medioambiente?

6. Verdadero o falso (justifica con una cita) : \esp{Los jóvenes no pueden hacer nada por el planeta.}

7. ¿Y tú, qué haces para proteger el planeta?
\end{exercice}

\begin{corrige}[Exercice : Compréhension du texte]
1. \esp{El texto habla del cambio climático, de sus causas y consecuencias, y de las soluciones para proteger el medioambiente.}

2. \esp{La contaminación, la deforestación y el uso excesivo de los coches aumentan los gases de efecto invernadero.}

3. \esp{En el norte, la sequía es cada vez más larga y el lago Chad ha perdido más del noventa por ciento de su superficie.}

4. \esp{Es importante porque sus árboles absorben el dióxido de carbono y producen oxígeno.}

5. \esp{Los ciudadanos pueden reciclar la basura, plantar árboles, ahorrar agua y electricidad y utilizar el transporte público.}

6. Falso : \esp{como dice el texto, ``los jóvenes tienen un papel esencial: ellos son el futuro del planeta''.}

7. Réponse libre, par exemple : \esp{Yo apago la luz, reciclo el plástico y planto árboles en mi barrio.}
\end{corrige}

\begin{aretenir}[À retenir]
Pour réussir la compréhension d'un texte : lis deux fois, souligne les mots-clés, réponds toujours par une \textbf{phrase complète en espagnol}, reprends les mots de la question et justifie avec \esp{según el texto} + une citation entre guillemets. Méfie-toi des faux amis comme \esp{actualmente} (« actuellement ») et ne traduis jamais mot à mot. Retiens aussi la différence entre \esp{¿por qué?} (« pourquoi ») et \esp{porque} (« parce que »).
\end{aretenir}

\coursec{Lexique thématique : el medioambiente}

Dans la section précédente, nous avons lu et compris le texte \esp{El planeta en peligro}. Nous y avons rencontré des mots comme \esp{calentamiento global}, \esp{sequía} ou \esp{deforestación}. Pour bien parler de l'environnement, il ne suffit pas de comprendre ces mots : il faut savoir les employer. Cette section organise donc tout le vocabulaire du texte en quatre grandes familles, faciles à mémoriser et à réutiliser à l'oral comme à l'écrit.

\courssub{Observer d'abord : un mini-texte de réinvestissement}

Avant de classer le vocabulaire, lisons ce court article de presse original. Souligne mentalement tous les mots qui se rapportent à l'environnement.

\begin{textebox}[El clima cambia en el norte de Camerún]
En el norte de Camerún, cerca del lago Chad, el clima cambia rápidamente. La sequía es cada vez más larga y el agua del lago disminuye año tras año. Los pescadores y los agricultores sufren las consecuencias del calentamiento global: los peces desaparecen y la tierra está seca.

En el sur del país, el problema es diferente: la deforestación amenaza el bosque del Congo, uno de los pulmones verdes del planeta. Cada año, miles de árboles desaparecen por culpa de la explotación de la madera.

Sin embargo, hay esperanza. Los gobiernos, las ONG y los ciudadanos toman medidas para proteger el medioambiente. En las escuelas de Yaoundé, los alumnos aprenden a reciclar los residuos y a plantar árboles. En algunas aldeas, los paneles solares producen energía limpia gracias al sol.

Los ecologistas repiten un mensaje claro: hay que reducir la contaminación, ahorrar agua y energía, y utilizar las energías renovables. Si todos luchamos contra el cambio climático, salvaremos el planeta para las generaciones futuras. Como dice el famoso lema: ¡No hay planeta B!
\end{textebox}

\textbf{Glossaire :} \esp{la sequía} = la sécheresse ; \esp{disminuye} = diminue ; \esp{los pescadores} = les pêcheurs ; \esp{los pulmones verdes} = les poumons verts ; \esp{la madera} = le bois ; \esp{las aldeas} = les villages ; \esp{el lema} = le slogan.

\textbf{Questions de compréhension :}
\begin{enumerate}
\item ¿Qué problemas climáticos sufren el norte y el sur de Camerún?
\item ¿Qué hacen los alumnos de Yaoundé para proteger el medioambiente?
\item ¿Qué soluciones proponen los ecologistas?
\end{enumerate}

\courssub{Champ lexical 1 : los problemas}

Voici les mots qui désignent les dangers qui menacent la planète. Apprends-les avec leur article : en espagnol comme en français, le genre est essentiel.

\begin{center}
\begin{tabularx}{\linewidth}{|l|X|X|}\hline
\rowcolor{popblueL} \textbf{Español} & \textbf{Français} & \textbf{Ejemplo} \\ \hline
el calentamiento global & le réchauffement climatique & \esp{El calentamiento global aumenta.} \\ \hline
la contaminación & la pollution & \esp{La contaminación del aire es grave.} \\ \hline
la sequía & la sécheresse & \esp{La sequía destruye las cosechas.} \\ \hline
la inundación & l'inondation & \esp{Las inundaciones destruyen casas.} \\ \hline
la deforestación & la déforestation & \esp{La deforestación amenaza la selva.} \\ \hline
los residuos / las basuras & les déchets / les ordures & \esp{Hay muchos residuos en la ciudad.} \\ \hline
el efecto invernadero & l'effet de serre & \esp{El efecto invernadero calienta la Tierra.} \\ \hline
el agujero de la capa de ozono & le trou de la couche d'ozone & \esp{El agujero de la capa de ozono es peligroso.} \\ \hline
\end{tabularx}
\end{center}

\begin{attention}[Les accents qui sauvent des points]
Le mot \esp{la sequía} (la sécheresse) prend un accent sur le \esp{í} : on prononce « sé-ki-a » en insistant sur le \esp{i}. Écrire \esp{sequia} sans accent est une faute classique des francophones. Même vigilance pour \esp{energía}, \esp{río}, \esp{protección}, \esp{contaminación} : tous les mots en \esp{-ción} portent un accent sur le \esp{ó}.
\end{attention}

\courssub{Champ lexical 2 : la naturaleza}

\begin{center}
\begin{tabularx}{\linewidth}{|l|X|}\hline
\rowcolor{popgreenL} \textbf{Español} & \textbf{Français} \\ \hline
el medioambiente & l'environnement \\ \hline
el planeta / la Tierra & la planète / la Terre \\ \hline
el bosque / la selva & la forêt / la forêt tropicale \\ \hline
el lago / el río / el mar & le lac / la rivière, le fleuve / la mer \\ \hline
el clima & le climat \\ \hline
la tierra & la terre (le sol) \\ \hline
\end{tabularx}
\end{center}

\begin{attention}[Faux amis et pièges de genre]
\begin{itemize}
\item \esp{el medioambiente} s'écrit en un seul mot et signifie « l'environnement » (au sens écologique). Ne dis pas \esp{el ambiente} pour parler de la planète : \esp{el ambiente} désigne plutôt « l'ambiance, l'atmosphère » d'un lieu.
\item \esp{La Tierra} avec majuscule = la planète Terre ; \esp{la tierra} avec minuscule = le sol, la terre que l'on cultive.
\item Attention au genre : \esp{el clima}, \esp{el planeta}, \esp{el problema}, \esp{el sistema}, \esp{el tema}, \esp{el programa}, \esp{el mapa}, \esp{el día} sont masculins malgré leur terminaison en \esp{-a}.
\end{itemize}
\end{attention}

\courssub{Champ lexical 3 : las soluciones}

\begin{center}
\begin{tabularx}{\linewidth}{|l|X|X|}\hline
\rowcolor{popgreenL} \textbf{Español} & \textbf{Français} & \textbf{Ejemplo} \\ \hline
las energías renovables & les énergies renouvelables & \esp{Las energías renovables no contaminan.} \\ \hline
la energía solar / eólica & l'énergie solaire / éolienne & \esp{La energía solar es limpia.} \\ \hline
reciclar & recycler & \esp{Reciclamos el plástico y el papel.} \\ \hline
reutilizar & réutiliser & \esp{Reutilizamos las bolsas.} \\ \hline
reducir & réduire & \esp{Hay que reducir los residuos.} \\ \hline
plantar árboles & planter des arbres & \esp{Los alumnos plantan árboles.} \\ \hline
proteger / cuidar & protéger / prendre soin de & \esp{Debemos proteger la selva.} \\ \hline
ahorrar agua y energía & économiser l'eau et l'énergie & \esp{Debemos ahorrar agua.} \\ \hline
\end{tabularx}
\end{center}

\begin{definition}[Las energías renovables]
Les \esp{energías renovables} sont des énergies produites à partir de sources naturelles inépuisables : le soleil (\esp{la energía solar}), le vent (\esp{la energía eólica}) et l'eau (\esp{la energía hidráulica}). Contrairement au pétrole (\esp{el petróleo}) et au charbon (\esp{el carbón}), elles ne s'épuisent pas et polluent peu.
\end{definition}

\begin{exemplebox}[Des phrases à réutiliser]
\begin{itemize}
\item \esp{Hay que reciclar los residuos.} = « Il faut recycler les déchets. »
\item \esp{El bosque está en peligro.} = « La forêt est en danger. »
\item \esp{Debemos ahorrar agua.} = « Nous devons économiser l'eau. »
\item \esp{Los paneles solares producen energía limpia.} = « Les panneaux solaires produisent une énergie propre. »
\end{itemize}
\end{exemplebox}

\courssub{Champ lexical 4 : los actores}

Qui agit pour la planète ? Quatre groupes d'acteurs :

\begin{itemize}
\item \esp{los ecologistas} : les écologistes, ceux qui défendent la nature ;
\item \esp{las ONG} (organizaciones no gubernamentales) : les ONG ;
\item \esp{los gobiernos} : les gouvernements, qui prennent des lois ;
\item \esp{los ciudadanos} : les citoyens, c'est-à-dire nous tous ;
\item \esp{las generaciones futuras} : les générations futures, pour qui nous protégeons la planète.
\end{itemize}

\begin{exemplebox}
\esp{Los gobiernos toman medidas contra la contaminación.} = « Les gouvernements prennent des mesures contre la pollution. »\\
\esp{Los ciudadanos deben luchar contra el cambio climático.} = « Les citoyens doivent lutter contre le changement climatique. »
\end{exemplebox}

\courssub{Verbes-clés et expressions utiles}

\begin{center}
\begin{tabularx}{\linewidth}{|l|X|l|}\hline
\rowcolor{poporangeL} \textbf{Verbo} & \textbf{Français} & \textbf{Famille de mots} \\ \hline
contaminar & polluer & la contaminación, contaminado \\ \hline
proteger & protéger & la protección, protegido \\ \hline
reciclar & recycler & el reciclaje, reciclable \\ \hline
reducir & réduire & la reducción \\ \hline
aumentar & augmenter & el aumento \\ \hline
desaparecer & disparaître & la desaparición \\ \hline
amenazar & menacer & la amenaza \\ \hline
salvar & sauver & la salvación \\ \hline
cuidar & prendre soin de & el cuidado \\ \hline
\end{tabularx}
\end{center}

\begin{propriete}[Quatre expressions à connaître par cœur]
\begin{itemize}
\item \esp{estar en peligro (de extinción)} = être en danger (d'extinction) : \esp{Los elefantes están en peligro de extinción.} « Les éléphants sont en danger d'extinction. »
\item \esp{luchar contra} + nom = lutter contre : \esp{Luchamos contra la deforestación.}
\item \esp{tomar medidas} = prendre des mesures : \esp{El gobierno toma medidas.}
\item \esp{hay que} + infinitif = il faut + infinitif (tournure impersonnelle) : \esp{Hay que cuidar el planeta.} « Il faut prendre soin de la planète. »
\end{itemize}
\end{propriete}

\begin{attention}[Attention aux verbes \esp{proteger} et \esp{reducir}]
\esp{Proteger} est irrégulier à la 1\iere personne du présent : \esp{yo protejo} (avec un \esp{j}), mais \esp{tú proteges}, \esp{él protege}. De même, \esp{reducir} donne \esp{yo reduzco} (cédille). Pour « sauvegarder un fichier », on utilise \esp{guardar} ; \esp{salvar} s'emploie pour « sauver une vie, la planète » : \esp{salvar el planeta} = « sauver la planète ».
\end{attention}

\courssub{Carte mentale : tout le lexique en un coup d'œil}

\begin{popfigure}
\begin{center}
\begin{tikzpicture}[mindmap, grow cyclic,
  every node/.style={concept, align=center, font=\small},
  root concept/.append style={fill=popgreen, font=\bfseries},
  level 1/.append style={level distance=3.4cm, sibling angle=90},
  level 2/.append style={level distance=2.2cm, sibling angle=45, font=\scriptsize}]
\node[root concept]{El medioambiente}
  child[concept color=poporange!70]{ node {Problemas}
    child { node {contaminación} }
    child { node {sequía} }
    child { node {deforestación} }
  }
  child[concept color=popblue!60]{ node {Naturaleza}
    child { node {el bosque} }
    child { node {el río} }
    child { node {el clima} }
  }
  child[concept color=popgreen!60]{ node {Soluciones}
    child { node {reciclar} }
    child { node {plantar árboles} }
    child { node {energías renovables} }
  }
  child[concept color=poppurple!60]{ node {Actores}
    child { node {las ONG} }
    child { node {los gobiernos} }
    child { node {los ciudadanos} }
  };
\end{tikzpicture}
\end{center}
\legende{Carte mentale du lexique de l'environnement : quatre familles de mots.}
\end{popfigure}

\begin{popfigure}
\begin{center}
\begin{tikzpicture}[
  prob/.style={draw=poporange, fill=poporangeL, rounded corners, align=center, font=\small, minimum width=3.8cm, minimum height=0.85cm},
  sol/.style={draw=popgreen, fill=popgreenL, rounded corners, align=center, font=\small, minimum width=3.8cm, minimum height=0.85cm},
  tit/.style={font=\bfseries},
  flecha/.style={-{Stealth[length=3mm]}, thick, popmuted}]
\node[tit, poporange] at (-2.8,0.9) {Problemas};
\node[tit, popgreen] at (2.8,0.9) {Soluciones};
\node[prob] (p1) at (-2.8,0) {la contaminación};
\node[sol] (s1) at (2.8,0) {reciclar los residuos};
\node[prob] (p2) at (-2.8,-1.3) {la deforestación};
\node[sol] (s2) at (2.8,-1.3) {plantar árboles};
\node[prob] (p3) at (-2.8,-2.6) {el calentamiento global};
\node[sol] (s3) at (2.8,-2.6) {energías renovables};
\node[prob] (p4) at (-2.8,-3.9) {la sequía};
\node[sol] (s4) at (2.8,-3.9) {ahorrar agua};
\draw[flecha] (p1) -- (s1);
\draw[flecha] (p2) -- (s2);
\draw[flecha] (p3) -- (s3);
\draw[flecha] (p4) -- (s4);
\end{tikzpicture}
\end{center}
\legende{À chaque problème, sa solution : des associations à mémoriser.}
\end{popfigure}

\courssub{Méthode et entraînement}

\begin{methode}[Comment mémoriser durablement ce lexique]
\begin{enumerate}
\item Regroupe les mots par \cle{familles} : \esp{contaminar} (verbe) → \esp{la contaminación} (nom) → \esp{contaminado} (adjectif). Un seul effort, trois mots appris.
\item Classe chaque mot nouveau dans une colonne : \esp{Problemas}, \esp{Naturaleza}, \esp{Soluciones}, \esp{Actores}.
\item Fabrique des \cle{phrases personnelles} sur ton quartier ou ton village : \esp{En mi barrio hay mucha basura; hay que reciclar.} Le cerveau retient ce qui te concerne.
\item Réutilise les expressions toutes faites (\esp{hay que}, \esp{estar en peligro}, \esp{luchar contra}) : elles donnent des points à l'oral comme à l'écrit.
\end{enumerate}
\end{methode}

\begin{exoresolu}[Classer le vocabulaire]
Énoncé : classe ces douze mots dans les deux colonnes \esp{Problemas} et \esp{Soluciones} : \esp{la sequía} — \esp{reciclar} — \esp{la deforestación} — \esp{plantar árboles} — \esp{la contaminación} — \esp{ahorrar agua} — \esp{la inundación} — \esp{las energías renovables} — \esp{los residuos} — \esp{proteger el bosque} — \esp{el efecto invernadero} — \esp{reutilizar}.
\tcblower
Solution détaillée :
\begin{center}
\begin{tabularx}{\linewidth}{|X|X|}\hline
\rowcolor{poporangeL} \textbf{Problemas} & \textbf{Soluciones} \\ \hline
la sequía & reciclar \\ \hline
la deforestación & plantar árboles \\ \hline
la contaminación & ahorrar agua \\ \hline
la inundación & las energías renovables \\ \hline
los residuos & proteger el bosque \\ \hline
el efecto invernadero & reutilizar \\ \hline
\end{tabularx}
\end{center}
Méthode : demande-toi si le mot désigne un \cle{danger} pour la nature (problème) ou une \cle{action positive} (solution). \esp{Los residuos} est un problème ; \esp{reciclar} est la solution correspondante.
\end{exoresolu}

\begin{exercice}[À toi de jouer]
\begin{enumerate}
\item Complète avec le mot qui convient : \esp{sequía}, \esp{contaminación}, \esp{reciclar}, \esp{renovable}, \esp{peligro}.
\begin{enumerate}
\item En el norte de Camerún, la \ldots\ destruye las cosechas.
\item Hay que \ldots\ el papel y el plástico.
\item La energía solar es una energía \ldots
\item Los coches producen mucha \ldots\ del aire.
\item El lago Chad está en \ldots.
\end{enumerate}
\item Traduis en espagnol : a) « Il faut économiser l'eau. » b) « Les gouvernements doivent lutter contre la déforestation. » c) « Les éléphants sont en danger d'extinction. »
\end{enumerate}
\end{exercice}

\begin{corrige}[Exercice : À toi de jouer]
1. a) \esp{sequía} ; b) \esp{reciclar} ; c) \esp{renovable} (accord au singulier : \esp{una energía renovable}) ; d) \esp{contaminación} ; e) \esp{peligro}.\\
2. a) \esp{Hay que ahorrar agua.} b) \esp{Los gobiernos deben luchar contra la deforestación.} c) \esp{Los elefantes están en peligro de extinción.}
\end{corrige}

\begin{savaistu}[« No hay planeta B »]
\esp{¡No hay planeta B!} (« Il n'y a pas de planète B ! ») est le slogan des marches pour le climat dans tout le monde hispanophone, de Madrid à Mexico. Les jeunes l'écrivent sur leurs pancartes lors des grandes manifestations du vendredi. Retiens aussi le jeu de mots espagnol : \esp{planeta} rime avec \esp{meta} (« objectif ») — \esp{Salvar el planeta es nuestra meta} : « Sauver la planète est notre objectif ».
\end{savaistu}

\begin{aretenir}[L'essentiel du lexique]
Quatre familles : les \cle{problèmes} (\esp{contaminación, sequía, deforestación}), la \cle{nature} (\esp{medioambiente, bosque, río, clima}), les \cle{solutions} (\esp{reciclar, plantar árboles, energías renovables, ahorrar}) et les \cle{acteurs} (\esp{ecologistas, ONG, gobiernos, ciudadanos}). Quatre expressions clés : \esp{estar en peligro}, \esp{luchar contra}, \esp{tomar medidas}, \esp{hay que} + infinitif. Et n'oublie jamais les accents : \esp{sequía, energía, río, protección}.
\end{aretenir}

\coursec{Grammaire : le futur simple}

Dans la section précédente, nous avons enrichi notre vocabulaire sur le thème de \esp{el medioambiente} : \esp{calentamiento global}, \esp{contaminación}, \esp{sequía}, \esp{deforestación}, \esp{energías renovables}, \esp{reciclar}. Mais un vocabulaire ne suffit pas : pour parler de l'avenir de la planète, pour annoncer ce qui \emph{se passera} si nous ne changeons rien, il faut un temps verbal précis. Ce temps, c'est le \cle{futur simple} (\esp{el futuro simple}). C'est lui que nous allons observer, comprendre et maîtriser maintenant.

\courssub{Observation : le futur dans le texte}

Reprenons trois phrases tirées de notre texte support \emph{El planeta en peligro} et observons les verbes en gras.

\begin{exemplebox}[Phrases du texte à observer]
\begin{itemize}
\item \esp{Si actuamos ahora, \textbf{protegeremos} el medioambiente.} « Si nous agissons maintenant, nous \textbf{protégerons} l'environnement. »
\item \esp{El clima \textbf{cambiará} si no reducimos la contaminación.} « Le climat \textbf{changera} si nous ne réduisons pas la pollution. »
\item \esp{\textbf{Habrá} más sequías en el norte de Camerún.} « Il y \textbf{aura} plus de sécheresses dans le nord du Cameroun. »
\end{itemize}
\end{exemplebox}

\textbf{Questions d'observation.}
\begin{enumerate}
\item De quel infinitif viennent \esp{protegeremos}, \esp{cambiará} et \esp{habrá} ? Réponse : \esp{proteger}, \esp{cambiar}, \esp{haber}.
\item Que remarque-t-on ? Les terminaisons \esp{-emos} et \esp{-á} se sont ajoutées \emph{à l'infinitif entier}, et non à un radical coupé comme en français.
\item Pourquoi \esp{habrá} ne ressemble-t-il pas à \esp{haber} ? Parce que \esp{haber} a un radical irrégulier : \esp{habr-}.
\end{enumerate}

Voici un court texte original qui utilise massivement le futur simple. Lisez-le, soulignez mentalement tous les futurs, puis vérifiez avec le glossaire.

\begin{textebox}[El futuro de nuestro planeta]
Los científicos dicen que el clima cambiará mucho en los próximos años. Si no protegemos los bosques, la deforestación aumentará y muchos animales perderán su casa. En el norte de Camerún, habrá más sequías y el lago Chad será todavía más pequeño. En otras regiones, las lluvias serán más fuertes y habrá inundaciones.

Pero todavía podemos actuar. Si reciclamos más, contaminaremos menos. Si plantamos árboles, el aire será más limpio. Los gobiernos invertirán en energías renovables como el sol y el viento. Los jóvenes también participarán : en mi escuela de Yaundé, nosotros reciclaremos el plástico y plantaremos flores. Mis amigos dicen : « Yo no tiraré basura al suelo y usaré menos bolsas de plástico ».

El futuro del planeta dependerá de nosotros. Si todos hacemos un pequeño esfuerzo, la Tierra será un lugar mejor para nuestros hijos. ¿Qué harás tú mañana para proteger el medioambiente ?
\end{textebox}

\textbf{Glossaire.} \esp{los próximos años} : les prochaines années ; \esp{aumentará} : augmentera ; \esp{perderán} : perdront ; \esp{las lluvias} : les pluies ; \esp{invertirán} : investiront ; \esp{tirar basura} : jeter des ordures ; \esp{el esfuerzo} : l'effort ; \esp{dependerá de} : dépendra de.

\textbf{Questions.} 1) Relevez cinq verbes au futur dans le texte et donnez leur infinitif. 2) Quels problèmes climatiques sont annoncés pour le Cameroun ? 3) Quelles actions les jeunes feront-ils ?

\courssub{Formation du futur simple : une règle unique}

\begin{propriete}[Formation : infinitif + terminaisons]
Le futur simple espagnol se forme en ajoutant les terminaisons suivantes \textbf{directement à l'infinitif entier}, pour \textbf{tous} les verbes, qu'ils se terminent en \esp{-ar}, \esp{-er} ou \esp{-ir} :

\begin{center}
\begin{tabular}{|l|l|l|}
\hline
\rowcolor{popblueL} Personne & Terminaison & Exemple avec \esp{hablar} \\
\hline
\esp{yo} & \esp{-é} & \esp{hablaré} \\
\hline
\esp{tú} & \esp{-ás} & \esp{hablarás} \\
\hline
\esp{él / ella / usted} & \esp{-á} & \esp{hablará} \\
\hline
\esp{nosotros/-as} & \esp{-emos} & \esp{hablaremos} \\
\hline
\esp{vosotros/-as} & \esp{-éis} & \esp{hablaréis} \\
\hline
\esp{ellos / ellas / ustedes} & \esp{-án} & \esp{hablarán} \\
\hline
\end{tabular}
\end{center}

Chez les verbes irréguliers, \textbf{seul le radical change} ; les terminaisons restent exactement les mêmes. Il n'y a donc qu'une seule conjugaison à apprendre.
\end{propriete}

\textbf{Tableau complet des trois groupes.}

\begin{center}
\begin{tabular}{|l|l|l|l|}
\hline
\rowcolor{popblueL} Personne & \esp{hablar} (parler) & \esp{comer} (manger) & \esp{vivir} (vivre) \\
\hline
\esp{yo} & \esp{hablaré} & \esp{comeré} & \esp{viviré} \\
\hline
\esp{tú} & \esp{hablarás} & \esp{comerás} & \esp{vivirás} \\
\hline
\esp{él/ella/ud.} & \esp{hablará} & \esp{comerá} & \esp{vivirá} \\
\hline
\esp{nosotros} & \esp{hablaremos} & \esp{comeremos} & \esp{viviremos} \\
\hline
\esp{vosotros} & \esp{hablaréis} & \esp{comeréis} & \esp{viviréis} \\
\hline
\esp{ellos/uds.} & \esp{hablarán} & \esp{comerán} & \esp{vivirán} \\
\hline
\end{tabular}
\end{center}

\textbf{Conjugaison de deux verbes de notre leçon.}

\begin{center}
\begin{tabular}{|l|l|l|}
\hline
\rowcolor{popgreenL} Personne & \esp{proteger} (protéger) & \esp{reciclar} (recycler) \\
\hline
\esp{yo} & \esp{protegeré} & \esp{reciclaré} \\
\hline
\esp{tú} & \esp{protegerás} & \esp{reciclarás} \\
\hline
\esp{él/ella/ud.} & \esp{protegerá} & \esp{reciclará} \\
\hline
\esp{nosotros} & \esp{protegeremos} & \esp{reciclaremos} \\
\hline
\esp{vosotros} & \esp{protegeréis} & \esp{reciclaréis} \\
\hline
\esp{ellos/uds.} & \esp{protegerán} & \esp{reciclarán} \\
\hline
\end{tabular}
\end{center}

\begin{aretenir}[L'accent, la clé du futur]
Toutes les personnes du futur portent un \textbf{accent aigu obligatoire} sur la terminaison, \textbf{sauf} la première personne du pluriel : \esp{hablaré}, \esp{hablarás}, \esp{hablará}, \esp{hablaremos} (sans accent), \esp{hablaréis}, \esp{hablarán}. Sans accent, le sens change : \esp{hablara} est un imparfait du subjonctif, pas un futur !
\end{aretenir}

\courssub{Comparaison avec le français}

\begin{center}
\begin{tabularx}{\linewidth}{|X|X|X|}
\hline
\rowcolor{popblueL} Français & Español & Remarque \\
\hline
je parler\textbf{ai} & \esp{hablar\textbf{é}} & En français, on coupe l'infinitif (\emph{parler-}) ; en espagnol, l'infinitif reste \textbf{entier} : \esp{hablar} + \esp{é}. \\
\hline
nous finir\textbf{ons} & \esp{vivir\textbf{emos}} & Même logique : \esp{vivir} + \esp{emos}, sans couper le \esp{-ir}. \\
\hline
il y aura & \esp{habrá} & Radical irrégulier \esp{habr-} : à mémoriser. \\
\hline
tu feras & \esp{harás} & Radical irrégulier \esp{har-} : à mémoriser. \\
\hline
\end{tabularx}
\end{center}

L'espagnol est donc plus \emph{régulier} que le français dans sa formation : une fois les dix radicaux irréguliers connus, tout le reste est mécanique.

\begin{popfigure}
\begin{tikzpicture}[font=\small]
\draw[-{Stealth}, very thick, popblue] (0,0) -- (10.5,0);
\fill[popgreen] (1.2,0) circle (0.12);
\node[above, align=center, popgreen] at (1.2,0.15) {\textbf{ahora}\\ \esp{reciclo}\\ « je recycle »};
\fill[poporange] (5.2,0) circle (0.12);
\node[above, align=center, poporange] at (5.2,0.15) {\textbf{mañana}\\ \esp{reciclaré}\\ « je recyclerai »};
\fill[poppurple] (9.2,0) circle (0.12);
\node[above, align=center, poppurple] at (9.2,0.15) {\textbf{en el futuro}\\ \esp{reciclaremos}\\ « nous recyclerons »};
\node[below, popmuted] at (1.2,-0.15) {présent};
\node[below, popmuted] at (5.2,-0.15) {futur simple};
\node[below, popmuted] at (9.2,-0.15) {futur simple};
\end{tikzpicture}
\legende{Frise chronologique : du présent (\esp{reciclo}) au futur simple (\esp{reciclaré}, \esp{reciclaremos}).}
\end{popfigure}

\courssub{Les dix radicaux irréguliers}

Une dizaine de verbes très fréquents modifient leur radical au futur. Il faut les apprendre par cœur : ils reviennent dans presque tous les textes.

\begin{popfigure}
\begin{tikzpicture}[font=\small]
\node[draw=popblue, fill=popblueL, rounded corners, minimum width=3.2cm, minimum height=0.8cm] (t) at (0,0) {\textbf{Radicaux irréguliers}};
\node[draw=popgreen, fill=popgreenL, rounded corners, align=center, minimum width=4.6cm] (a) at (-4.6,-2.2) {\textbf{Perte de voyelle}\\ \esp{haber} $\to$ \esp{habr-}\\ \esp{poder} $\to$ \esp{podr-}\\ \esp{querer} $\to$ \esp{querr-}\\ \esp{saber} $\to$ \esp{sabr-}};
\node[draw=poporange, fill=poporangeL, rounded corners, align=center, minimum width=4.6cm] (b) at (0,-2.2) {\textbf{Voyelle $\to$ d}\\ \esp{poner} $\to$ \esp{pondr-}\\ \esp{salir} $\to$ \esp{saldr-}\\ \esp{tener} $\to$ \esp{tendr-}\\ \esp{venir} $\to$ \esp{vendr-}};
\node[draw=poppurple, fill=poppinkL, rounded corners, align=center, minimum width=4.6cm] (c) at (4.6,-2.2) {\textbf{Changement total}\\ \esp{decir} $\to$ \esp{dir-}\\ \esp{hacer} $\to$ \esp{har-}};
\draw[-{Stealth}, thick, popblue] (t) -- (a);
\draw[-{Stealth}, thick, popblue] (t) -- (b);
\draw[-{Stealth}, thick, popblue] (t) -- (c);
\end{tikzpicture}
\legende{Carte mentale : les dix radicaux irréguliers du futur, classés en trois familles.}
\end{popfigure}

\begin{center}
\begin{tabularx}{\linewidth}{|l|l|X|}
\hline
\rowcolor{popblueL} Verbe & Radical & Exemple traduit \\
\hline
\esp{decir} (dire) & \esp{dir-} & \esp{El profesor dirá la verdad.} « Le professeur dira la vérité. » \\
\hline
\esp{hacer} (faire) & \esp{har-} & \esp{¿Qué harás mañana ?} « Que feras-tu demain ? » \\
\hline
\esp{haber} (avoir, il y a) & \esp{habr-} & \esp{Habrá más inundaciones.} « Il y aura plus d'inondations. » \\
\hline
\esp{poder} (pouvoir) & \esp{podr-} & \esp{Podremos viajar.} « Nous pourrons voyager. » \\
\hline
\esp{querer} (vouloir) & \esp{querr-} & \esp{Ella querrá ayudar.} « Elle voudra aider. » \\
\hline
\esp{saber} (savoir) & \esp{sabr-} & \esp{Sabré la respuesta.} « Je saurai la réponse. » \\
\hline
\esp{poner} (mettre) & \esp{pondr-} & \esp{Pondré la mesa.} « Je mettrai la table. » \\
\hline
\esp{salir} (sortir) & \esp{saldr-} & \esp{Saldrán temprano.} « Ils sortiront tôt. » \\
\hline
\esp{tener} (avoir) & \esp{tendr-} & \esp{Tendré veinte años.} « J'aurai vingt ans. » \\
\hline
\esp{venir} (venir) & \esp{vendr-} & \esp{Vendrán a Douala.} « Ils viendront à Douala. » \\
\hline
\end{tabularx}
\end{center}

\begin{attention}[Pièges mortels : \esp{habrá} et \esp{haré}]
\begin{itemize}
\item \esp{Habrá} (« il y aura ») vient de \esp{haber} $\to$ radical \esp{habr-}. N'écrivez \textbf{jamais} \esp{haberá} : ce mot n'existe pas.
\item \esp{Haré} (« je ferai ») vient de \esp{hacer} $\to$ radical \esp{har-}. N'écrivez \textbf{jamais} \esp{haceré}.
\item Ne confondez pas le \textbf{futur} \esp{hablaré} (« je parlerai ») et le \textbf{conditionnel} \esp{hablaría} (« je parlerais »). Le futur a un accent aigu et pas de \esp{-í-} ; le conditionnel a un \esp{-í-} avec accent. \esp{Cantaré} = « je chanterai » ; \esp{cantaría} = « je chanterais ».
\end{itemize}
\end{attention}

\courssub{Emplois du futur simple}

\begin{definition}[Les trois emplois essentiels]
\begin{enumerate}
\item \textbf{La prédiction} : annoncer ce qui arrivera. \esp{Habrá más inundaciones en el futuro.} « Il y aura plus d'inondations à l'avenir. »
\item \textbf{La conséquence logique} : après \esp{si} + présent, le futur exprime la conséquence. \esp{Si contaminamos, el planeta sufrirá.} « Si nous polluons, la planète souffrira. »
\item \textbf{La promesse ou l'engagement} : s'engager à agir. \esp{Reciclaremos más y plantaremos árboles.} « Nous recyclerons davantage et nous planterons des arbres. »
\end{enumerate}
\end{definition}

\begin{exemplebox}[Exemples traduits]
\begin{itemize}
\item \esp{Si contaminamos, el planeta sufrirá.} « Si nous polluons, la planète souffrira. »
\item \esp{Mañana plantaré un árbol.} « Demain je planterai un arbre. »
\item \esp{Habrá más sequías en el norte.} « Il y aura plus de sécheresses dans le nord. »
\item \esp{Los gobiernos invertirán en energías renovables.} « Les gouvernements investiront dans les énergies renouvelables. »
\item \esp{Si cuidamos el bosque, los animales vivirán en paz.} « Si nous protégeons la forêt, les animaux vivront en paix. »
\item \esp{El lago Chad será más pequeño si no llueve.} « Le lac Tchad sera plus petit s'il ne pleut pas. »
\end{itemize}
\end{exemplebox}

\begin{methode}[Vérifier un futur en trois étapes]
\begin{enumerate}
\item \textbf{Retrouver l'infinitif} du verbe à conjuguer (ex. « je ferai » $\to$ \esp{hacer}).
\item \textbf{Vérifier si le radical est irrégulier} : le verbe fait-il partie de la liste des dix (\esp{decir, hacer, haber, poder, querer, saber, poner, salir, tener, venir}) ? Si oui, prendre le radical irrégulier (\esp{hacer} $\to$ \esp{har-}) ; sinon, garder l'infinitif entier.
\item \textbf{Ajouter la terminaison} avec l'accent au bon endroit : \esp{-é, -ás, -á, -emos, -éis, -án}. Résultat : \esp{haré}. Vérification finale : \esp{hablaremos} est la seule forme sans accent.
\end{enumerate}
\end{methode}

\begin{exoresolu}[Conjuguez au futur simple]
Énoncé. Mettez les verbes entre parenthèses au futur simple : 1) \esp{Nosotros (proteger) el medioambiente.} 2) \esp{Si no llueve, (haber) una gran sequía.} 3) \esp{¿Qué (hacer) tú para proteger el planeta ?} 4) \esp{Mis amigos (reciclar) el plástico.} 5) \esp{Yo (venir) a la reunión mañana.}
\tcblower
Solution détaillée. 1) \esp{protegeremos} : verbe régulier, infinitif entier + \esp{-emos}, pas d'accent à \esp{nosotros}. 2) \esp{habrá} : \esp{haber} est irrégulier, radical \esp{habr-} + \esp{-á} ; attention, jamais \esp{haberá}. 3) \esp{harás} : \esp{hacer} est irrégulier, radical \esp{har-} + \esp{-ás}. 4) \esp{reciclarán} : verbe régulier, infinitif + \esp{-án}. 5) \esp{vendré} : \esp{venir} est irrégulier, radical \esp{vendr-} + \esp{-é}. Traductions : 1) « Nous protégerons l'environnement. » 2) « S'il ne pleut pas, il y aura une grande sécheresse. » 3) « Que feras-tu pour protéger la planète ? » 4) « Mes amis recycleront le plastique. » 5) « Je viendrai à la réunion demain. »
\end{exoresolu}

\begin{exercice}[À vous de jouer]
Traduisez en espagnol : 1) Demain, nous planterons des arbres à Yaoundé. 2) S'ils polluent l'eau, les poissons mourront (\esp{morir}). 3) Il y aura moins de déchets si nous recyclons. 4) Je pourrai expliquer la leçon à mes camarades. 5) Que direz-vous au ministre de l'Environnement ?
\end{exercice}

\begin{corrige}[Exercice « À vous de jouer »]
1) \esp{Mañana plantaremos árboles en Yaundé.} 2) \esp{Si contaminan el agua, los peces morirán.} 3) \esp{Habrá menos basura si reciclamos.} (radical \esp{habr-} !) 4) \esp{Podré explicar la lección a mis compañeros.} (radical \esp{podr-}) 5) \esp{¿Qué dirán ustedes al ministro del Medioambiente ?} ou, à la deuxième personne : \esp{¿Qué diréis al ministro... ?} (radical \esp{dir-}).
\end{corrige}

\begin{savaistu}[Le futur simple au Bac]
À l'épreuve d'espagnol du Baccalauréat A, le futur simple est très souvent exigé dans les sujets de production écrite sur l'environnement. La question classique est : \esp{¿Qué harás para proteger el planeta ?} « Que feras-tu pour protéger la planète ? » Pour obtenir une excellente note, répondez avec au moins cinq futurs variés, dont deux irréguliers : \esp{Reciclaré, plantaré árboles, haré un esfuerzo, habrá menos contaminación gracias a nosotros, y vendré a ayudar cada semana.} Un futur correct et accentué impressionne toujours le correcteur !
\end{savaistu}

\coursec{Grammaire : \texorpdfstring{\esp{si} + présent $\rightarrow$ futur}{si + présent → futur} (la conséquence)}

Après avoir maîtrisé la formation et les emplois du \cle{futur simple} (section précédente), nous abordons maintenant l'une des structures les plus utiles pour exprimer une \cle{condition réalisable} et sa \cle{conséquence probable} : la construction \textbf{\esp{si} + présent de l'indicatif $\rightarrow$ futur simple}. C'est l'outil grammatical indispensable pour parler d'écologie, de prévisions climatiques ou d'engagements citoyens : \esp{« Si actuamos ahora, salvaremos el planeta. »}

\courssub{Observation à partir du texte support}

Reprenons la phrase extraite de notre texte \emph{El planeta en peligro} :
\begin{exemplebox}[Exemple du texte]
\esp{Si usamos energías renovables, protegeremos el medioambiente.}
« Si nous utilisons des énergies renouvelables, nous protégerons l'environnement. »
\end{exemplebox}

\noindent Analysons la structure :
\begin{itemize}
    \item La subordonnée conditionnelle (introduite par \esp{si}) est au \textbf{présent de l'indicatif} : \esp{usamos} (verbe \emph{usar}, 1\textsuperscript{re} personne du pluriel).
    \item La principale exprime la conséquence au \textbf{futur simple} : \esp{protegeremos} (verbe \emph{proteger}, 1\textsuperscript{re} personne du pluriel, radical régulier + terminaison \emph{-emos}).
\end{itemize}
Cette phrase illustre parfaitement une \textbf{condition réelle} (l'utilisation d'énergies renouvelables est possible) et sa \textbf{conséquence probable/quasi certaine} (la protection de l'environnement en découle logiquement).

\courssub{Règle fondamentale : \texorpdfstring{\esp{Si} + présent $\rightarrow$ futur}{Si + présent → futur}}

\begin{propriete}[Structure de la condition réelle et de sa conséquence]
Pour exprimer une \textbf{hypothèse réalisable dans le présent ou le futur immédiat} et sa \textbf{conséquence future certaine ou très probable}, l'espagnol utilise exactement la même structure que le français :
\medskip

\begin{center}
\begin{tabularx}{\linewidth}{|X|X|X|}
\hline
\rowcolor{popblueL}
\textbf{Français} & \textbf{Espagnol} & \textbf{Niveau de langue} \\
\hline
Si + \textbf{présent de l'indicatif} & \esp{Si} + \textbf{presente de indicativo} & Condition réalisable \\
\hline
$\rightarrow$ \textbf{futur simple} & $\rightarrow$ \textbf{futuro simple} & Conséquence probable \\
\hline
\end{tabularx}
\end{center}

\medskip
\noindent \textbf{Formule mnémotechnique :} \cle{\esp{Si} + presente $\rightarrow$ futuro}.
\end{propriete}

\begin{aretenir}[Lexique clé pour s'exprimer sur le climat]
\begin{center}
\begin{tabularx}{\linewidth}{|l|X|l|}
\hline
\rowcolor{popgreenL}
\textbf{Espagnol} & \textbf{Français} & \textbf{Exemple} \\
\hline
\esp{el calentamiento global} & le réchauffement climatique & \esp{El calentamiento global es una amenaza.} \\
\esp{la contaminación} & la pollution & \esp{La contaminación del aire aumenta.} \\
\esp{la sequía} & la sécheresse & \esp{La sequía destruye las cosechas.} \\
\esp{la inundación} & l'inondation & \esp{Las inundaciones son frecuentes en Douala.} \\
\esp{la deforestación} & la déforestation & \esp{La deforestación avanza en el Congo.} \\
\esp{las energías renovables} & les énergies renouvelables & \esp{Invertimos en energías renovables.} \\
\esp{reciclar} & recycler & \esp{Debemos reciclar el plástico.} \\
\esp{proteger / salvar} & protéger / sauver & \esp{Protegeremos la selva.} \\
\hline
\end{tabularx}
\end{center}
\end{aretenir}

\begin{exemplebox}[Exemples variés (contexte camerounais et hispanophone)]
\esp{Si contaminamos el río Wouri, los peces morirán.} \\
« Si nous polluons le fleuve Wouri, les poissons mourront. » \\[4pt]
\esp{Si no actuamos ahora, el lago Chad desaparecerá.} \\
« Si nous n'agissons pas maintenant, le lac Tchad disparaîtra. » \\[4pt]
\esp{Si plantamos árboles en la selva del Congo, frenaremos la deforestación.} \\
« Si nous plantons des arbres dans la forêt du Congo, nous freinerons la déforestation. » \\[4pt]
\esp{Si ahorras agua en casa, ayudarás al planeta.} \\
« Si tu économises l'eau à la maison, tu aideras la planète. » \\[4pt]
\esp{Si el gobierno invierte en energías solares, reduciremos la contaminación.} \\
« Si le gouvernement investit dans le solaire, nous réduirons la pollution. »
\end{exemplebox}

\courssub{Ordre des propositions : la virgule}

En espagnol comme en français, l'ordre des deux propositions peut être inversé sans changer le sens. \textbf{La ponctuation change en revanche} :

\begin{propriete}[Position de \texorpdfstring{\esp{si}}{si} et virgule]
\begin{itemize}
    \item \textbf{\esp{Si} en tête} : on sépare les deux propositions par une \textbf{virgule}.
    \item \textbf{\esp{Si} en seconde position} : \textbf{pas de virgule} avant \esp{si}.
\end{itemize}
\end{propriete}

\begin{exemplebox}[Inversion sans changement de sens]
\esp{Protegeremos el planeta si reciclamos.} \\
« Nous protégerons la planète si nous recyclons. » (pas de virgule) \\[4pt]
\esp{Si reciclamos, protegeremos el planeta.} \\
« Si nous recyclons, nous protégerons la planète. » (virgule obligatoire)
\end{exemplebox}

\courssub{Illustration : Schéma de la structure conditionnelle}

\begin{popfigure}
\begin{tikzpicture}[
    node distance=1.5cm and 2cm,
    box/.style={draw=popblue, thick, rounded corners, fill=popblueL, minimum width=3.8cm, minimum height=1.8cm, align=center, font=\sffamily},
    arrow/.style={-Stealth, thick, poporange, shorten >=2pt, shorten <=2pt}
]
\node[box, align=center] (si){\textbf{Si + Presente de Indicativo}\\\small{(Condition réalisable)}\\\esp{Si reciclamos}};
\node[box, right=of si, align=center] (futuro){\textbf{Futuro Simple}\\\small{(Conséquence probable)}\\\esp{protegeremos}};
\draw[arrow] (si.east) -- (futuro.west) node[midway, above, font=\small\sffamily\bfseries, popdark] {$\rightarrow$ entraîne};
\end{tikzpicture}
\legende{Schéma de la relation condition $\rightarrow$ conséquence : \esp{Si reciclamos $\rightarrow$ protegeremos}.}
\end{popfigure}

\courssub{Comparaison systématique Français / Espagnol}

C'est un point de convergence majeur qui doit vous rassurer : \textbf{la structure est rigoureusement identique dans les deux langues}. Le piège ne vient pas de la différence, mais de l'oubli de la règle commune.

\begin{center}
\begin{tabularx}{\linewidth}{|X|X|X|}
\hline
\rowcolor{popblueL}
\textbf{Français} & \textbf{Espagnol} & \textbf{Observation} \\
\hline
Si nous recyclons, nous protégerons. & \esp{Si reciclamos, protegeremos.} & Structure identique : Si + Présent $\rightarrow$ Futur. \\
\hline
Nous protégerons la planète si nous recyclons. & \esp{Protegeremos el planeta si reciclamos.} & Inversion identique : pas de virgule avant \emph{si} / \esp{si}. \\
\hline
Si je vais à Yaoundé, je verrai le match. & \esp{Si voy a Yaundé, veré el partido.} & Verbes irréguliers au futur (\esp{voy, veré}) mais même logique. \\
\hline
S'il pleut, il y aura des inondations. & \esp{Si llueve, habrá inundaciones.} & \esp{Hay} $\rightarrow$ \esp{habrá} (verbe \emph{haver} impersonnel). \\
\hline
\end{tabularx}
\end{center}

\courssub{Variante utile pour le Bac : \texorpdfstring{\esp{Si} + présent $\rightarrow$ impératif}{Si + présent → impératif}}

Souvent, la conséquence d'une condition n'est pas une simple prévision, mais un \textbf{conseil}, un \textbf{ordre} ou une \textbf{exhortation}. On utilise alors l'\textbf{impératif} dans la principale. Cette structure est très fréquente dans les campagnes écologiques et peut tomber à l'examen (compréhension ou expression).

\begin{propriete}[\texorpdfstring{\esp{Si} + présent $\rightarrow$ impératif}{Si + présent → impératif} (Conseil / Ordre conditionnel)]
\esp{Si + presente de indicativo} $\rightarrow$ \textbf{Imperativo afirmativo}.
\end{propriete}

\begin{exemplebox}[Exemples style campagne / Bac]
\esp{Si quieres salvar el planeta, recicla.} \\
« Si tu veux sauver la planète, recycle. » (Impératif \emph{tú}) \\[4pt]
\esp{Si quieren beber agua limpia, no contaminen los ríos.} \\
« Si vous voulez boire de l'eau propre, ne polluez pas les rivières. » (Impératif \emph{ustedes} / \emph{vous}) \\[4pt]
\esp{Si vamos a la playa, recojamos la basura.} \\
« Si nous allons à la plage, ramassons les déchets. » (Impératif \emph{nosotros} / \emph{nous})
\end{exemplebox}

\courssub{Les deux pièges mortels (Erreurs fréquentes des francophones)}

\begin{attention}[Piège n°1 : JAMAIS de futur après \texorpdfstring{\esp{si}}{si}]
C'est l'erreur \textbf{n°1} qui fait perdre des points bêtement. En espagnol \textbf{comme en français}, le futur est \textbf{interdit} dans la proposition introduite par \esp{si} (condition).
\begin{itemize}
    \item \textcolor{popgreen}{\textbf{Correct :}} \esp{Si reciclamos, protegeremos.} (Présent après \esp{si})
    \item \textcolor{popred}{\textbf{Incorrect :}} \esp{Si reciclaremos, protegeremos.} (\textsc{FAUX} — Futur après \esp{si})
\end{itemize}
\noindent \textbf{Rappel français :} On ne dit pas « Si nous recyclerons... » mais « Si nous recyclons... ». La règle est \textbf{strictement la même}. Ne traduisez pas mot à mot le futur français « nous protégerons » en le mettant après \esp{si} !
\end{attention}

\begin{attention}[Piège n°2 : \texorpdfstring{\esp{si}}{si} (sans accent) $\neq$ \texorpdfstring{\esp{sí}}{sí} (avec accent)]
Ne confondez jamais la conjonction de condition et l'adverbe d'affirmation.
\begin{itemize}
    \item \esp{si} (sans accent) = « si » (conjonction). \esp{Si llueve, no voy.}
    \item \esp{sí} (avec accent) = « oui » (réponse) ou « soi » (pronom réfléchi tonique). \esp{¡Sí, quiero!} / \esp{Él lo hizo por sí mismo.}
\end{itemize}
\noindent \textbf{Exemple combiné :} \esp{— ¿Reciclas? — Sí, si reciclas, ganarás.} (« — Tu recycles ? — Oui, si tu recycles, tu gagneras. »)
\end{attention}

\courssub{Méthode : Construire une phrase de conséquence réelle}

\begin{methode}[Écrire une phrase \texorpdfstring{\esp{Si} + presente $\rightarrow$ futuro}{Si + présent → futur} sans erreur]
\begin{enumerate}
    \item \textbf{Identifiez la condition} (ce qui déclenche l'action) : elle suit \esp{si}. $\rightarrow$ Mettez le verbe au \textbf{présent de l'indicatif}.
    \item \textbf{Identifiez la conséquence} (ce qui arrivera si la condition est remplie) : c'est la principale. $\rightarrow$ Mettez le verbe au \textbf{futur simple}.
    \item \textbf{Vérifiez l'ordre} : Si \esp{si} est au début $\rightarrow$ \textbf{virgule}. Si \esp{si} est au milieu $\rightarrow$ \textbf{pas de virgule}.
    \item \textbf{Contrôlez les irrégularités du futur} : \esp{saldré, pondré, tendré, vendré, haré, sabré, querré, podré, cabré, habrá, diré, estaré}. (Voir section précédente).
    \item \textbf{Relisez} : Y a-t-il un futur après \esp{si} ? \textbf{Non.} Y a-t-il un accent sur \esp{sí} sans raison ? \textbf{Non.}
\end{enumerate}
\end{methode}

\courssub{Méthode : Commenter un texte argumentatif sur l'environnement}

\begin{methode}[Étapes pour un commentaire de texte réussi (Bac / Évaluation)]
\begin{enumerate}
    \item \textbf{Contexte} : Identifiez la source (article, discours, campagne), le thème (climat, pollution, biodiversité) et la thèse défendue.
    \item \textbf{Structure} : Repérez l'introduction (constat), le développement (arguments \texttt{si + présent $\rightarrow$ futur}, données, exemples locaux), la conclusion (appel à l'action, souvent \texttt{si + présent $\rightarrow$ impératif}).
    \item \textbf{Lexique} : Surlignez le vocabulaire thématique (voir boîte \emph{Lexique clé}) et les connecteurs logiques (\esp{sin embargo, por tanto, además, en conclusión}).
    \item \textbf{Grammaire en action} : Relevez les structures \texttt{si + présent $\rightarrow$ futur/impératif} ; analysez pourquoi l'auteur les choisit (probabilité, exhortation).
    \item \textbf{Expression personnelle} : Rédigez un paragraphe \emph{« Y yo, ¿qué hago? »} en réemployant 3 structures étudiées.
\end{enumerate}
\end{methode}

\courssub{Exercice résolu : Thème (Français $\rightarrow$ Espagnol)}

\begin{exoresolu}[Thème : 5 phrases sur l'environnement]
\textbf{Énoncé :} Traduisez en espagnol en appliquant la règle \esp{Si} + présent $\rightarrow$ futur.
\begin{enumerate}
    \item Si nous plantons des arbres, nous sauverons la forêt.
    \item S'il pleut moins, il y aura des sécheresses.
    \item Si vous (vous) recyclez le plastique, vous réduirez les déchets.
    \item Nous protégerons le lac Tchad si nous agissons maintenant.
    \item Si le gouvernement interdit le sac plastique, la ville sera plus propre.
\end{enumerate}

\tcblower
\textbf{Solution détaillée :}

\begin{enumerate}
    \item \esp{Si plantamos árboles, salvaremos el bosque.}
    \begin{itemize}
        \item \esp{plantamos} (présent \emph{plantar}, 1\textsuperscript{re} pl.) — \esp{salvaremos} (futur \emph{salvar}, radical régulier).
        \item Vocabulaire : \esp{el bosque} (la forêt, terme général) ou \esp{la selva} (forêt tropicale).
    \end{itemize}
    \item \esp{Si llueve menos, habrá sequías.}
    \begin{itemize}
        \item \esp{llueve} (présent \emph{llover}, 3\textsuperscript{e} sg., diphtongue \emph{ue}) — \esp{habrá} (futur \emph{haver} impersonnel, \textbf{irrégulier}).
        \item \esp{sequías} (sécheresses, pluriel).
    \end{itemize}
    \item \esp{Si reciclan el plástico, reducirán los residuos.}
    \begin{itemize}
        \item Contexte « vous » (formel \emph{ustedes}, standard au Cameroun / Amérique latine / Guinée équatoriale) : \esp{reciclan} (présent), \esp{reducirán} (futur \emph{reducir}).
        \item \emph{Variante péninsulaire (vosotros) :} \esp{Si recicláis el plástico, reduciréis los residuos.}
        \item \esp{los residuos} = les déchets / les ordures.
    \end{itemize}
    \item \esp{Protegaremos el lago Chad si actuamos ahora.}
    \begin{itemize}
        \item Inversion : \esp{si} en seconde position $\rightarrow$ \textbf{pas de virgule}.
        \item \esp{Protegaremos} (futur \emph{proteger}) — \esp{actuamos} (présent \emph{actuar}).
    \end{itemize}
    \item \esp{Si el gobierno prohíbe las bolsas plásticas, la ciudad estará más limpia.}
    \begin{itemize}
        \item \esp{prohíbe} (présent \emph{prohibir}, 3\textsuperscript{e} sg., accent sur \emph{í} pour casser la diphtongue) — \esp{estará} (futur \emph{estar}, \textbf{irrégulier} : radical \emph{estarr-} $\rightarrow$ \esp{estará}).
        \item \esp{las bolsas plásticas} = les sacs plastiques. \esp{más limpia} = plus propre (accord avec \esp{ciudad} fém. sg.).
        \item \textbf{Choix \emph{estar} :} état resultant / aspect physique après l'action, non caractéristique essentielle (\emph{ser}).
    \end{itemize}
\end{enumerate}
\end{exoresolu}

\courssub{Texte d'entraînement : \emph{El futuro está en nuestras manos}}

\begin{textebox}[Texte original — Niveau A2+/B1]
\esp{
El cambio climático ya no es una amenaza lejana; es una realidad en Camerún y en todo el mundo. En el norte, la sequía avanza y el lago Chad se reduce cada año. En el sur, la deforestación de la selva del Congo amenaza la biodiversidad y el clima global. Las inundaciones en Douala son más frecuentes y destructivas.

Sin embargo, no todo está perdido. Si cambiamos nuestros hábitos, frenaremos el calentamiento global. Si usamos el transporte público o la bicicleta, emitiremos menos CO2. Si reciclamos el plástico y el papel, no contaminaremos los ríos. Si plantamos un árbol cada año, el aire será más puro.

Los jóvenes cameruneses tienen la palabra. Si exigen energías renovables, el gobierno invertirá en paneles solares. Si educan a sus hermanos pequeños, la próxima generación respetará la naturaleza. Si quieres salvar el planeta, actúa ahora. El futuro no se escribe solo; se construye. ¿Y tú, qué harás mañana por el planeta?
}

\medskip
\noindent \textbf{Glossaire des mots difficiles :}
\begin{itemize}
    \item \esp{amenaza} = menace ; \esp{lejana} = lointaine ; \esp{se reduce} = se réduit / diminue (verbe pronominal \emph{reducirse}).
    \item \esp{amenaza la biodiversidad} = menace la biodiversité (suj. \emph{la deforestación}, verbe sg.).
    \item \esp{hábito} = habitude ; \esp{frenaremos} = nous freinerons (futur \emph{frenar}).
    \item \esp{emitiremos} = nous émettons (futur \emph{emitir}) ; \esp{CO2} = CO2 (dioxyde de carbone).
    \item \esp{exigen} = ils/elles exigent (présent \emph{exigir}) ; \esp{invertirá} = il/elle investira (futur \emph{invertir}).
    \item \esp{paneles solares} = panneaux solaires ; \esp{se construye} = se construit / se bâtit (passif réfléchi).
    \item \esp{actúa} = agis (impératif \emph{tú} de \emph{actuar}).
\end{itemize}

\medskip
\noindent \textbf{Questions de compréhension / analyse grammaticale :}
\begin{enumerate}
    \item Relevez \textbf{les 7 phrases} construites sur le modèle \esp{Si} + présent $\rightarrow$ futur (6) ou impératif (1) dans le texte.
    \item Pour la phrase \esp{« Si plantamos un árbol cada año, el aire será más puro »}, identifiez le verbe au présent et le verbe au futur. Pourquoi \esp{será} est-il irrégulier ?
    \item Transformez la phrase \esp{« Si exigen energías renovables, el gobierno invertirá en paneles solares »} en mettant \esp{si} en seconde position (attention à la virgule !).
    \item Trouvez dans le texte l'unique exemple de \esp{Si} + présent $\rightarrow$ \textbf{impératif} et précisez la personne concernée.
\end{enumerate}
\end{textebox}

\courssub{Exercices d'application}

\begin{exercice}[Mise en pratique : Complétez avec le bon temps]
Conjuguez le verbe entre parenthèses au \textbf{présent} (après \esp{si}) ou au \textbf{futur simple} (conséquence).
\begin{enumerate}
    \item Si nosotros \_\_\_\_\_\_\_\_\_\_ (contaminar) el aire, nosotros \_\_\_\_\_\_\_\_\_\_ (respirar) aire sucio.
    \item Si tú \_\_\_\_\_\_\_\_\_\_ (querer) ayudar, tú \_\_\_\_\_\_\_\_\_\_ (poder) empezar hoy.
    \item \_\_\_\_\_\_\_\_\_\_ (haber) menos basura si la gente \_\_\_\_\_\_\_\_\_\_ (reciclar) más.
    \item Yo \_\_\_\_\_\_\_\_\_\_ (comprar) una bici si yo \_\_\_\_\_\_\_\_\_\_ (tener) dinero.
    \item Si llover \_\_\_\_\_\_\_\_\_\_ (llover) mucho, el río \_\_\_\_\_\_\_\_\_\_ (desbordarse).
\end{enumerate}
\end{exercice}

\begin{corrige}[Exercice : Mise en pratique]
\begin{enumerate}
    \item \esp{contaminamos} (présent) — \esp{respiraremos} (futur \emph{respirar}).
    \item \esp{quieres} (présent \emph{querer}, diphtongue) — \esp{podrás} (futur \emph{poder}, \textbf{irrégulier} : \emph{podr-}).
    \item \esp{Habrá} (futur \emph{haver} impersonnel, \textbf{irrégulier}) — \esp{recicla} (présent, \emph{la gente} = sg.). \textbf{Note :} \esp{la gente} est singulier en espagnol $\rightarrow$ \esp{recicla}.
    \item \esp{compraré} (futur \emph{comprar}) — \esp{tengo} (présent \emph{tener}, \textbf{irrégulier} : \emph{tengo}).
    \item \esp{llueve} (présent \emph{llover}, diphtongue) — \esp{se desbordará} (futur \emph{desbordarse}, pronominal).
\end{enumerate}
\end{corrige}

\begin{exercice}[Thème guidé : Engagements pour le Bac]
Traduisez les phrases suivantes en espagnol. Utilisez la structure \esp{Si} + présent $\rightarrow$ futur (ou impératif si c'est un conseil).
\begin{enumerate}
    \item Si nous économisons l'électricité, nous paierons moins cher. (\esp{la electricidad}, \esp{pagar menos})
    \item Si tu veux un avenir meilleur, plante un arbre. (Impératif \emph{tú})
    \item Le Cameroun sera plus vert si nous protégeons nos forêts. (Inversion, pas de virgule)
    \item S'ils (ils) coupent les arbres, la terre sera sèche. (\esp{cortar}, \esp{la tierra}, \esp{seco})
    \item Si nous n'utilisons pas la voiture, nous marcherons plus. (\esp{usar el coche}, \esp{caminar más})
\end{enumerate}
\end{exercice}

\begin{corrige}[Exercice : Thème guidé]
\begin{enumerate}
    \item \esp{Si ahorramos electricidad, pagaremos menos.}
    \item \esp{Si quieres un futuro mejor, planta un árbol.} (Impératif \emph{planta}).
    \item \esp{Camerún será más verde si protegemos nuestros bosques.} (Pas de virgule avant \esp{si}).
    \item \esp{Si cortan los árboles, la tierra estará seca.} (\esp{estará} car \esp{la tierra} sg ; \esp{seca} accord fém. sg). \textbf{Attention :} \esp{estar} au futur (\textbf{irrégulier} \esp{estará}) car état resultant, pas \esp{ser}.
    \item \esp{Si no usamos el coche, caminaremos más.}
\end{enumerate}
\end{corrige}

\courssub{Expression orale guidée : Débat citoyen}

\begin{experience}[Jeu de rôle : « La voz de la juventud »]
\textbf{Consigne :} En binômes. L'un est un journaliste de \emph{CRTV} / \emph{Radio Nacional de Guinea Ecuatorial}, l'autre un jeune militant écolo.
\begin{itemize}
    \item \textbf{Journaliste :} Posez 3 questions sur l'avenir du climat au Cameroun (utilisez \esp{¿Qué pasará si...?}, \esp{¿Crees que...?}).
    \item \textbf{Militant :} Répondez en utilisant \textbf{au moins 3 phrases} \esp{Si + présent $\rightarrow$ futur} et \textbf{1 phrase} \esp{Si + présent $\rightarrow$ impératif} pour lancer un appel à l'action.
    \item \textbf{Exemple :} \esp{— ¿Qué pasará si no protegemos el bosque? — Si no lo protegemos, perderemos biodiversidad. ¡Si quieres ayudar, planta un árbol hoy!}
\end{itemize}
\textbf{Durée :} 3 min par binôme. Les meilleurs passages sont présentés à la classe.
\end{experience}

\courssub{Synthèse visuelle : Tableau comparatif Français / Espagnol (TikZ)}

\begin{popfigure}
\begin{tikzpicture}
\begin{scope}[every node/.style={font=\small\sffamily}]
% Header
\node[draw=popblue, fill=popblueL, minimum width=5.5cm, minimum height=0.9cm, align=center] (fr) at (0,0) {\textbf{Français}};
\node[draw=popblue, fill=popblueL, minimum width=5.5cm, minimum height=0.9cm, align=center] (es) at (6,0) {\textbf{Español}};
% Structure
\node[draw=popgreen, fill=popgreenL, minimum width=5.5cm, minimum height=1.2cm, align=center, anchor=north] (fr1) at (0,-0.45) {\textbf{Si} + \textcolor{popdark}{Présent indicatif}\\(\textit{Condition réalisable})\\(Si nous recyclons)};
\node[draw=popgreen, fill=popgreenL, minimum width=5.5cm, minimum height=1.2cm, align=center, anchor=north] (es1) at (6,-0.45) {\textbf{Si} + \textcolor{popdark}{Presente de indicativo}\\(\textit{Condición realizable})\\(\esp{Si reciclamos})};
% Arrow
\draw[-Stealth, very thick, poporange] (fr1.south) -- ++(0,-0.3) -| (es1.south) node[pos=0.25, left, font=\tiny] {Même logique};
% Consequence
\node[draw=poppurple, fill=poppurpleL, minimum width=5.5cm, minimum height=1.2cm, align=center, anchor=north] (fr2) at (0,-2.1) {\textbf{Futur simple}\\(\textit{Conséquence probable})\\(nous protégerons)};
\node[draw=poppurple, fill=poppurpleL, minimum width=5.5cm, minimum height=1.2cm, align=center, anchor=north] (es2) at (6,-2.1) {\textbf{Futuro simple}\\(\textit{Consecuencia probable})\\(\esp{protegeremos})};
% Inversion
\node[draw=poppink, fill=poppinkL, minimum width=11.5cm, minimum height=1cm, align=center, anchor=north] (inv) at (3,-3.7) {\textbf{Inversion :} Pas de virgule avant « si » / \esp{si} \\(\esp{Protegeremos el planeta si reciclamos} = Nous protégerons la planète si nous recyclons)};
% Trap
\node[draw=popred, fill=popred!15, minimum width=11.5cm, minimum height=0.9cm, align=center, anchor=north] (trap) at (3,-5.0) {\textbf{\textcolor{popred}{PIÈGE :}} Jamais de futur après \esp{si} / \emph{si} \textcolor{popred}{(Si reciclaremos $\times$)} \quad | \quad \esp{si} (condition) $\neq$ \esp{sí} (oui)};
\end{scope}
\end{tikzpicture}
\legende{Tableau récapitulatif : convergence totale FR/ES sur la structure \emph{Si + Présent $\rightarrow$ Futur}.}
\end{popfigure}

\courssub{Pour aller plus loin : Contexte culturel — La Guinée équatoriale}

\begin{savaistu}[L'espagnol vert de Guinée équatoriale]
La Guinée équatoriale, seul pays africain hispanophone, est en première ligne face au changement climatique. Sa forêt tropicale (partie du bassin du Congo) est un \emph{sumidero de carbono} (puits de carbone) vital. Le pays mise sur le \emph{desarrollo sostenible} (développement durable) et les \emph{energías renovables} pour diversifier son économie pétrolière. À Malabo et Bata, des campagnes utilisent exactement la structure étudiée : \esp{« Si cuidas el bosque, cuidas tu futuro »} (Si tu protèges la forêt, tu protèges ton avenir). Maîtriser \esp{si} + présent $\rightarrow$ futur, c'est aussi pouvoir lire et comprendre les engagements écologiques de notre voisin équatorial.
\end{savaistu}

\coursec{Méthode du commentaire et production guidée}

Après avoir maîtrisé le futur simple et la structure \esp{si + présent + futur} pour exprimer la conséquence, nous allons maintenant mettre ces outils au service de deux compétences majeures du Bac : le \textbf{commentaire de texte} et la \textbf{production écrite guidée}. Le texte support \emph{El planeta en peligro} (étudié en sections 1 et 2) servira de base à notre commentaire modèle.

\courssub{La méthode du commentaire de texte en 4 étapes}

Au Bac d'espagnol (LV2), l'épreuve de commentaire suit une structure rigoureuse. Voici la démarche pas à pas, illustrée sur notre texte.

\begin{methode}[Commenter un texte de presse sur l'environnement]
\textbf{Étape 1 : Présenter le texte} (Introducción)
\begin{enumerate}
    \item \textbf{Nature :} Article de presse / Reportage / Éditorial / Extrait de roman ?
    \item \textbf{Thème :} Le changement climatique, la protection de l'environnement.
    \item \textbf{Source / Auteur / Date :} Si connus (sinon : \esp{« un artículo de prensa reciente »}).
    \item \textbf{Annonce du plan :} \esp{« En primer lugar analizaremos las causas... luego las consecuencias... y finalmente las soluciones propuestas. »}
\end{enumerate}

\textbf{Étape 2 : Dégager l'idée principale et la thèse} (Idea principal y tesis)
\begin{enumerate}
    \item Lire le texte une première fois pour la compréhension globale.
    \item Formuler en une phrase : \esp{« El texto denuncia la gravedad del calentamiento global y llama a la acción inmediata. »}
    \item Identifier la \textbf{thèse} de l'auteur : est-il alarmiste ? optimiste ? critique envers les gouvernements ?
\end{enumerate}

\textbf{Étape 3 : Analyser les arguments avec citations} (Análisis de los argumentos)
\begin{enumerate}
    \item Repérer les \textbf{arguments} (souvent 3 : Causes, Conséquences, Solutions).
    \item Pour chaque argument : \textbf{Citer} (entre guillemets \esp{« ... »}) $\rightarrow$ \textbf{Expliquer} avec ses mots $\rightarrow$ \textbf{Relier} à la thèse.
    \item Utiliser les connecteurs logiques (voir tableau ci-dessous).
\end{enumerate}

\textbf{Étape 4 : Donner son avis argumenté} (Opinión personal)
\begin{enumerate}
    \item Prendre position : \esp{Estoy de acuerdo / No estoy de acuerdo / En mi opinión...}
    \item Justifier avec \textbf{un exemple personnel ou local} (Cameroun : lac Tchad, déforestation, inondations à Douala).
    \item Employer \textbf{le futur} et \textbf{si + présent + futur} pour montrer la maîtrise grammaticale.
\end{enumerate}
\end{methode}

\begin{attention}[Pièges à éviter au commentaire]
\begin{itemize}
    \item \textbf{Le résumé :} Ne pas raconter le texte. Il faut \emph{analyser} (comment l'auteur construit son argumentation).
    \item \textbf{L'absence de citations :} Un argument sans citation du texte vaut 0 point. Citez \emph{toujours} : \esp{« la deforestación », « el nivel del mar subirá »}.
    \item \textbf{Le calque « Je pense que oui » :} $\rightarrow$ \esp{Pienso que sí} (pas \esp{pienso que oui}).
    \item \textbf{Les accents oubliés :} \esp{opinión, conclusión, protegeré, actuaré, habrá, será}. Un futur sans accent change le sens ou est fautif.
    \item \textbf{L'oubli du « si » + indicatif :} \esp{Si no actuamos, el planeta sufrirá} (pas de subjonctif après \esp{si} pour une condition réelle).
\end{itemize}
\end{attention}

\courssub{Plan type et connecteurs logiques}

Pour structurer le commentaire, on suit ce plan canonique. Le tableau ci-dessous récapitule les connecteurs indispensables pour articuler les parties et les arguments.

\begin{center}
\begin{tabularx}{\linewidth}{|X|X|X|}
\hline
\rowcolor{popblueL}
\textbf{Partie du commentaire} & \textbf{Connecteurs espagnols} & \textbf{Équivalent français} \\
\hline
\textbf{Introducción} & \esp{El texto que vamos a comentar es... ; Trata de... ; El autor denuncia...} & Le texte que nous allons commenter est... ; Il traite de... ; L'auteur dénonce... \\
\hline
\textbf{Idea principal} & \esp{La idea principal es... ; El autor sostiene que... ; La tesis central es...} & L'idée principale est... ; L'auteur soutient que... ; La thèse centrale est... \\
\hline
\textbf{Análisis (Arguments)} & \esp{En primer lugar / Primero / En cuanto a las causas...} & En premier lieu / Premièrement / Quant aux causes... \\
& \esp{Además / También / Por otra parte...} & De plus / Aussi / D'autre part... \\
& \esp{Por ejemplo / Como se ve en el texto: «...»} & Par exemple / Comme on le voit dans le texte : «...» \\
& \esp{Sin embargo / No obstante / En cambio...} & Cependant / Néanmoins / En revanche... \\
& \esp{Por eso / Por lo tanto / Como consecuencia...} & C'est pourquoi / Par conséquent / Comme conséquence... \\
\hline
\textbf{Opinión personal} & \esp{En mi opinión / A mi modo de ver / Pienso que...} & À mon avis / À mon sens / Je pense que... \\
& \esp{Estoy de acuerdo / No estoy de acuerdo porque...} & Je suis d'accord / Je ne suis pas d'accord car... \\
& \esp{En mi país, Camerún, por ejemplo...} & Dans mon pays, le Cameroun, par exemple... \\
\hline
\textbf{Conclusión} & \esp{En conclusión / Para terminar / En definitiva...} & En conclusion / Pour terminer / En définitive... \\
& \esp{El texto nos invita a... ; Es fundamental que...} & Le texte nous invite à... ; Il est fondamental que... \\
\hline
\end{tabularx}
\end{center}

\begin{savaistu}[L'astuce du Bac : la question « Donnez votre avis »]
Au Bac, la dernière question de compréhension ou l'instruction de production est souvent : \esp{« ¿Cuál es tu opinión sobre el tema ? »} ou \esp{« ¿Qué propones tú ? »}.
C'est une \textbf{aubaine} : elle rapporte des points faciles si vous respectez la « règle des 3 » :
\begin{enumerate}
    \item Ouvrir avec \esp{En mi opinión,} (ou \esp{A mi juicio,}).
    \item Utiliser \textbf{au moins un futur} (\esp{protegeré, reciclaré, cambiaré}).
    \item Inclure \textbf{une phrase avec \esp{si + présent + futur}} (\esp{Si todos reciclamos, el mundo será mejor}).
\end{enumerate}
En 2-3 phrases bien construites, vous validez la compétence « Expression écrite » et la grammaire du module.
\end{savaistu}

\courssub{Commentaire modèle : « El planeta en peligro »}

Voici un commentaire modèle (niveau B1, environ 11 lignes), rédigé comme attendu en classe ou au Bac. Il sert de référence pour l'analyse en classe.

\begin{textebox}[Commentaire modèle — El planeta en peligro]
El texto « El planeta en peligro » es un artículo de prensa que trata del cambio climático. El autor explica las causas del calentamiento global, como la contaminación industrial y la deforestación masiva, y sus consecuencias dramáticas: las sequías prolongadas, las inundaciones y la subida del nivel del mar.

En primer lugar, el texto destaca que la actividad humana es la responsable principal. Cita el ejemplo del lago Chad, que ha perdido el 90\% de su superficie, un caso muy cercano a Camerún. Además, menciona la selva del Congo como pulmón verde amenazado. Por eso, el autor propone soluciones urgentes: usar energías renovables, reciclar y proteger los bosques.

En mi opinión, el texto tiene toda la razón. Es alarmante ver cómo el lago Chad desaparece. En Camerún, si no protegemos nuestros bosques, sufriremos más sequías. Por eso, yo actuaré: reciclaré los plásticos, plantaré árboles y usaré transporte público. En conclusión, el texto nos invita a actuar ahora porque si no actuamos hoy, el planeta sufrirá mañana.
\end{textebox}

\begin{exemplebox}[Analyse du commentaire modèle]
\begin{itemize}
    \item \textbf{Introducción :} Nature (artículo de prensa), thème (cambio climático), thèse implicite (urgence d'agir).
    \item \textbf{Idea principal :} \esp{« El autor explica las causas... y sus consecuencias... »}
    \item \textbf{Análisis :} 3 arguments (Causes, Conséquences locales, Solutions) avec citations implicites (\esp{lago Chad, selva del Congo, energías renovables}) et connecteurs (\esp{En primer lugar, Además, Por eso}).
    \item \textbf{Opinión personal :} \esp{En mi opinión}, exemple local (Camerún), \textbf{4 futurs} (\esp{actuaré, reciclaré, plantaré, usaré}), \textbf{2 structures \esp{si}} (\esp{si no protegemos..., si no actuamos...}).
    \item \textbf{Conclusión :} \esp{En conclusión}, reprise de la thèse + ouverture.
\end{itemize}
\end{exemplebox}

\courssub{Production guidée : « ¿Qué harás tú para proteger el medioambiente ? »}

Consigne : Rédigez un paragraphe cohérent (80-100 mots) en répondant à cette question. Vous devez \textbf{impérativement} inclure :
\begin{enumerate}
    \item Au moins \textbf{3 verbes au futur simple} (différents si possible, 1\up{re} personne du singulier).
    \item Au moins \textbf{2 phrases avec \esp{si + présent + futur}}.
    \item Au moins \textbf{3 mots du lexique du climat} (voir section 3 : \esp{calentamiento global, contaminación, sequía, inundación, deforestación, energías renovables, reciclar}).
    \item Au moins \textbf{2 connecteurs} (ex. \esp{por eso, además, en conclusión}).
\end{enumerate}

\begin{methode}[Escalier de la production réussie]
Suivez ces 4 marches pour construire votre paragraphe sans erreur :
\begin{enumerate}
    \item \textbf{Mots-clés (Lexique) :} Choisissez 3 mots : \esp{reciclar, sequía, energías renovables, contaminación, deforestación, inundación}.
    \item \textbf{Verbes d'action (Futur) :} Conjuguez 3 verbes à la 1\up{re} personne du singulier : \esp{protegeré, reduciré, plantaré, cambiaré, usaré, separaré}.
    \item \textbf{Conditions (Si + Présent + Futur) :} Créez 2 liens cause-effet : \esp{Si reciclo, reduciré la basura. / Si planto árboles, habrá más oxígeno.}
    \item \textbf{Assemblage (Connecteurs) :} Reliez le tout : \esp{En mi opinión, protegeré el medioambiente. Primero, reciclaré el plástico. Además, si separo la basura, habrá menos contaminación. También plantaré árboles porque si reforestamos, evitaremos la sequía. En conclusión, pequeños gestos cambian el mundo.}
\end{enumerate}
\end{methode}

\begin{popfigure}
\begin{tikzpicture}[
    node distance=1.2cm and 1.5cm,
    box/.style={draw=popblue, thick, rounded corners, fill=popblueL, text width=3.2cm, align=center, minimum height=1.8cm, font=\small},
    arrow/.style={-Stealth, thick, popblue}
]
\node[box, align=center] (lex){\textbf{1. LEXIQUE}\\(3 mots min.)\\\esp{reciclar, sequía, energías renovables}};
\node[box, right=of lex, align=center] (fut){\textbf{2. FUTUR}\\(3 verbes min.)\\\esp{protegeré, reciclaré, plantaré}};
\node[box, right=of fut, align=center] (si){\textbf{3. SI + PRES/FUT}\\(2 phrases min.)\\\esp{Si reciclo, reduciré...}};
\node[box, right=of si, align=center] (par){\textbf{4. PARAGRAPHE}\\(Connecteurs + Cohérence)\\\esp{En mi opinión... Por eso...}};

\draw[arrow] (lex) -- node[above, font=\tiny, popdark] {choisir} (fut);
\draw[arrow] (fut) -- node[above, font=\tiny, popdark] {conjuguer} (si);
\draw[arrow] (si) -- node[above, font=\tiny, popdark] {construire} (par);

\node[draw=popgreen, thick, rounded corners, fill=popgreenL, fit=(lex)(fut)(si)(par), inner sep=5pt] (frame) {};
\node[above=2pt of frame.north, anchor=south, font=\bfseries\small, popgreen] {ESCALIER DE LA PRODUCTION RÉUSSIE};
\end{tikzpicture}
\legende{Schéma de construction du paragraphe : du lexique à la production finale.}
\end{popfigure}

\begin{exoresolu}[Production guidée — Écriture d'un paragraphe]
\textbf{Énoncé :} Rédigez le paragraphe demandé en respectant les 4 contraintes (3 futurs, 2 \esp{si}, 3 mots lexique, 2 connecteurs). Utilisez l'escalier ci-dessus.

\tcblower

\textbf{Proposition de rédaction (élève type) :}
\esp{En mi opinión, protegeré el medioambiente con acciones concretas. Primero, reciclaré todo el plástico y el papel en mi casa. Además, si separo los residuos, reduciré la contaminación del suelo. También plantaré árboles en mi barrio porque si reforestamos la ciudad, evitaremos la sequía y el calor extremo. En conclusión, cada gesto cuenta para salvar el planeta.}

\textbf{Vérification des contraintes :}
\begin{itemize}
    \item \textbf{3 Futurs (1\up{re} pers. sg.) :} \esp{protegeré, reciclaré, plantaré, reduciré} (\checkmark)
    \item \textbf{2 \esp{si} + prés/fut :} \esp{si separo..., reduciré} / \esp{si reforestamos..., evitaremos} (\checkmark)
    \item \textbf{3 Lexique :} \esp{medioambiente, contaminación, sequía} (\checkmark)
    \item \textbf{Connecteurs :} \esp{En mi opinión, Primero, Además, También, En conclusión} (\checkmark)
    \item \textbf{Accents :} Tous présents (\esp{protegeré, reciclaré, plantaré, reduciré, evitaré}) (\checkmark)
\end{itemize}
\end{exoresolu}

\courssub{Grille d'auto-évaluation avant de rendre la copie}

Avant de considérer votre production comme finie, cochez chaque case. Si une case n'est pas cochée, relisez la règle correspondante (sections 3, 4, 5) et corrigez.

\begin{center}
\begin{tabularx}{\linewidth}{|>{\hsize=0.8\hsize}X|>{\hsize=1.2\hsize}X|>{\hsize=1.2\hsize}X|>{\hsize=0.8\hsize}X|}
\hline
\rowcolor{popblueL}
\textbf{Critère} & \textbf{Indicateur de réussite} & \textbf{Où vérifier ?} & \textbf{OK ?} \\
\hline
\textbf{Lexique climat} & 3 mots minimum utilisés correctement (\esp{calentamiento global, sequía, reciclar...}) & Section 3 (Liste + Exercices) & $\square$ \\
\hline
\textbf{Futur simple} & 3 verbes à la 1\up{re} pers. sg. avec accents (\esp{-é, -ás, -á, -emos, -éis, -án}) & Section 4 (Tableau + Irréguliers) & $\square$ \\
\hline
\textbf{Structure \esp{si}} & 2 phrases \esp{si + indicatif présent + futur} (pas de subjonctif !) & Section 5 (Règle + Exemples) & $\square$ \\
\hline
\textbf{Connecteurs} & Au moins 2 connecteurs logiques (\esp{por eso, además, sin embargo, en conclusión}) & Tableau ci-dessus (Plan type) & $\square$ \\
\hline
\textbf{Cohérence} & Le paragraphe répond à la question, suit un fil logique (Idée $\rightarrow$ Action $\rightarrow$ Conséquence) & Méthode « Escalier » & $\square$ \\
\hline
\textbf{Orthographe} & Accents vérifiés (\esp{opinión, conclusión, protegeré, acción}), \esp{ñ} respecté (\esp{mañana, España}) &  box (Pièges) & $\square$ \\
\hline
\end{tabularx}
\end{center}

\begin{exemplebox}[Exemples traduits pour l'oral / rédaction]
\begin{itemize}
    \item \esp{En mi opinión, todos debemos actuar ya.} $\rightarrow$ « À mon avis, nous devons tous agir maintenant. »
    \item \esp{En conclusión, el texto nos invita a proteger el planeta.} $\rightarrow$ « En conclusion, le texte nous invite à protéger la planète. »
    \item \esp{A mi modo de ver, la solución está en las energías renovables.} $\rightarrow$ « À mon sens, la solution réside dans les énergies renouvelables. »
    \item \esp{Por eso, yo cambiaré mis hábitos.} $\rightarrow$ « C'est pourquoi, je changerai mes habitudes. »
\end{itemize}
\end{exemplebox}

\begin{exercice}[Entraînement au commentaire — Expression écrite]
Sur la base du texte \emph{El planeta en peligro} et de la méthode en 4 étapes, rédigez un commentaire complet (environ 150-180 mots) en respectant le plan : \textbf{Introducción / Idea principal / Análisis (3 arguments) / Opinión personal / Conclusión}.
\textbf{Contraintes obligatoires :}
\begin{enumerate}
    \item Au moins 3 citations du texte entre guillemets.
    \item Au moins 4 connecteurs différents.
    \item Une partie « Opinión personal » avec un exemple camerounais, 2 futurs et 1 phrase \esp{si + présent + futur}.
\end{enumerate}
\end{exercice}

\begin{corrige}[Exercice — Commentaire complet]
\textbf{Proposition de corrigé type (niveau B1/Bac) :}

\esp{Introducción:} El texto « El planeta en peligro » es un artículo de prensa reciente que trata del cambio climático. Su autor denuncia la situación crítica del medioambiente y propone soluciones urgentes.

\esp{Idea principal:} La tesis central es que la actividad humana causa el calentamiento global y que debemos actuar ya para evitar catástrofes.

\esp{Análisis:} En primer lugar, el autor explica las causas: « la contaminación industrial y la deforestación ». Además, detalla las consecuencias: « sequías prolongadas » e « inundaciones devastadoras ». Por otra parte, cita el caso del lago Chad: « ha perdido el 90\% de su superficie ». Por eso, propone soluciones: « energías renovables » y « reciclar ».

\esp{Opinión personal:} En mi opinión, el texto tiene razón. En Camerún, vemos las consecuencias: si no protegemos la selva del Congo, tendremos más sequías. Yo actuaré: reciclaré y plantaré árboles.

\esp{Conclusión:} En conclusión, el texto es una llamada a la acción. Si todos actuamos ahora, el planeta tendrá futuro.
\end{corrige}

\newpage

\coursec{Activité d'intégration}

\coursec{Lección EA05 : La protección del medioambiente y el cambio climático}

\courssub{Activité d'intégration : Concours d'affiches « Semana Verde »}

\courssub{Objectifs de l'activité}

À la fin de cette activité d'intégration, tu seras capable de :
\begin{itemize}
\item réutiliser le \cle{lexique} du changement climatique (\esp{calentamiento global, contaminación, sequía, inundación, deforestación, energías renovables, reciclar}) dans une production authentique ;
\item employer correctement le \cle{futur simple} pour annoncer des engagements et des conséquences ;
\item construire des phrases de conséquence avec la structure \esp{si + presente, futuro} ;
\item créer, en groupe, une \cle{affiche} de sensibilisation et la présenter oralement en une minute ;
\item voter et justifier ton choix avec des arguments simples en espagnol.
\end{itemize}

\courssub{Situation : l'annonce du concours}

Ton lycée de Yaoundé participe à la \esp{Semana Verde}, une semaine consacrée à la protection de l'environnement. Le club environnement du lycée, \esp{Jóvenes por el clima}, organise un concours d'affiches. La proviseure a annoncé la nouvelle au tableau d'affichage :

\begin{textebox}[Anuncio del club « Jóvenes por el clima »]
¡ATENCIÓN, ALUMNOS DE PRIMERA!

El club de medioambiente « Jóvenes por el clima » organiza un concurso de carteles para la Semana Verde del liceo.

El tema de este año : « Si cuidamos el planeta hoy, viviremos mejor mañana ».

Cada grupo de cuatro alumnos debe crear un cartel con :
— un eslogan original ;
— tres problemas del medioambiente (la deforestación, la contaminación, la sequía, las inundaciones...) ;
— tres compromisos concretos en futuro : « Reciclaremos..., Plantaremos..., Ahorraremos... ».

Cada grupo presentará su cartel a la clase en un minuto. Después, la clase votará el mejor cartel. Los tres mejores carteles serán expuestos en el pasillo del liceo durante toda la Semana Verde.

¡Participad! El planeta os necesita.
\end{textebox}

\emph{Traduction :} « Attention, élèves de Première ! Le club environnement "Jeunes pour le climat" organise un concours d'affiches pour la Semaine verte du lycée. Le thème de cette année : "Si nous prenons soin de la planète aujourd'hui, nous vivrons mieux demain". Chaque groupe de quatre élèves doit créer une affiche avec : un slogan original ; trois problèmes de l'environnement ; trois engagements concrets au futur. Chaque groupe présentera son affiche à la classe en une minute. Ensuite, la classe votera la meilleure affiche. Les trois meilleures affiches seront exposées dans le couloir du lycée pendant toute la Semaine verte. Participez ! La planète a besoin de vous. »

\courssub{Consignes de réalisation}

Travaillez par groupes de quatre. Suivez les étapes dans l'ordre.

\begin{enumerate}
\item \textbf{Compréhension du document.} Lisez l'annonce du club et répondez par écrit, en espagnol, à trois questions : \esp{¿Quién organiza el concurso? ¿Cuál es el tema de este año? ¿Qué tres elementos debe tener el cartel?}
\item \textbf{Remue-méninges lexical.} Dans votre groupe, listez au moins huit mots du lexique étudié dans la leçon (\esp{calentamiento global, contaminación, sequía, inundación, deforestación, energías renovables, reciclar, basura, árboles, río...}) et classez-les en deux colonnes : \esp{problemas} et \esp{soluciones}.
\item \textbf{Choix des trois problèmes.} Choisissez trois problèmes qui existent réellement au Cameroun ou dans le monde (par exemple la déforestation du bassin du Congo, la sécheresse autour du lac Tchad, la pollution plastique à Douala). Rédigez une phrase au présent pour chaque problème : \esp{La deforestación destruye los bosques.}
\item \textbf{Les trois engagements au futur.} Rédigez trois engagements de votre groupe au futur simple, à la première personne du pluriel : \esp{Reciclaremos el plástico. Plantaremos árboles en el patio del liceo.}
\item \textbf{Le slogan.} Créez un slogan court et frappant avec la structure \esp{si + presente, futuro} : \esp{Si reciclamos hoy, viviremos mejor mañana.}
\item \textbf{Fabrication de l'affiche.} Disposez sur la feuille : le slogan en grand en haut, les trois problèmes avec un petit dessin ou symbole chacun, les trois engagements en bas. Écrivez lisiblement et vérifiez les accents.
\item \textbf{Présentation orale.} Préparez un discours d'une minute : chaque membre du groupe dit au moins deux phrases (salutation, présentation du slogan, un problème, un engagement, phrase de conclusion).
\item \textbf{Le vote.} Écoutez les autres groupes, puis votez pour la meilleure affiche (pas la vôtre !) et justifiez votre vote avec une phrase : \esp{Voto por el grupo tres porque su eslogan es muy original.}
\end{enumerate}

\begin{methode}[Réussir sa présentation orale d'une minute]
\begin{enumerate}
\item \textbf{Saluer et identifier le groupe :} \esp{Buenos días a todos. Somos el grupo [Nombre].}
\item \textbf{Annoncer le slogan :} \esp{Nuestro eslogan es : « [Slogan] ».}
\item \textbf{Présenter les 3 problèmes (tour de table) :} \esp{El primer problema es... ; El segundo problema es... ; El tercer problema es...}
\item \textbf{Annoncer les 3 engagements (tour de table) :} \esp{Nos comprometemos a... ; Reciclaremos... ; Plantaremos...}
\item \textbf{Conclure par une phrase de conséquence motivante :} \esp{Si todos participamos, el planeta será más verde.}
\item \textbf{Appel au vote et remerciements :} \esp{¡Voten por nosotros! ¡Gracias por su atención!}
\end{enumerate}
\end{methode}

\begin{popfigure}
\begin{tikzpicture}[node distance=1.8cm, >=Stealth, font=\small]
\node[draw, rounded corners, fill=popblueL, align=center, text width=4.5cm] (si) {\textbf{Si} + \textcolor{popdark}{Présent de l'indicatif}\\(condition / cause réelle)};
\node[draw, rounded corners, fill=poporangeL, align=center, text width=4.5cm, right=of si] (futur) {\textcolor{popdark}{Futur simple}\\(conséquence future)};
\draw[->, thick, popblue] (si) -- node[above, font=\small\bfseries] {entraîne} (futur);
\node[below=1.2cm of si, align=center, text width=4.5cm] (ex1) {\esp{Si contaminamos los ríos,}};
\node[below=1.2cm of futur, align=center, text width=4.5cm] (ex2) {\esp{no tendremos agua limpia.}};
\end{tikzpicture}
\legende{Structure de la conséquence : \esp{Si + présent indicatif, futur simple}. Jamais de futur après \esp{si}.}
\end{popfigure}

\begin{attention}[Pièges à éviter dans votre affiche]
\begin{itemize}
\item Après \esp{si} de condition, \textbf{jamais de futur} : on écrit \esp{Si reciclamos, viviremos mejor} et non \esp{si reciclaremos}.
\item \textbf{Futur simple} : une seule forme, terminaisons sur l'infinitif entier : \esp{reciclaré, plantaremos, protegerán}. Attention aux irréguliers fréquents : \esp{hacer → haremos}, \esp{poder → podremos}, \esp{haber → habrá}, \esp{saber → sabremos}, \esp{poner → pondremos}, \esp{salir → saldremos}, \esp{tener → tendremos}, \esp{venir → vendremos}.
\item \esp{El medioambiente} s'écrit en un seul mot (on accepte aussi \esp{medio ambiente} en deux mots).
\item \textbf{Accents obligatoires} : \esp{contaminación, sequía, energía, árboles, inundación, población, región}.
\item \esp{Basura} (ordures) pas d'accent. \esp{Río} (accent sur le \emph{í} pour casser le diphtongue).
\end{itemize}
\end{attention}

\courssub{Ressources à mobiliser}

\begin{aretenir}[Boîte à outils de l'écocitoyen]
\textbf{Lexique essentiel :}
\begin{center}
\begin{tabularx}{\linewidth}{|X|X|X|}
\hline
\rowcolor{popblueL}
\textbf{Problème (Problema)} & \textbf{Solution / Action (Solución)} & \textbf{Autres termes} \\
\hline
\esp{el calentamiento global} & \esp{las energías renovables} & \esp{el planeta} \\
\esp{la contaminación} & \esp{reciclar} & \esp{los residuos / la basura} \\
\esp{la sequía} & \esp{proteger} & \esp{el agua} \\
\esp{la inundación} & \esp{plantar (árboles)} & \esp{el bosque / la selva} \\
\esp{la deforestación} & \esp{ahorrar (agua, luz)} & \esp{el lago / el río} \\
\hline
\end{tabularx}
\end{center}

\textbf{Futur simple (réguliers) :} infinitif + \esp{-é, -ás, -á, -emos, -éis, -án}.
\begin{center}
\begin{tabular}{|l|l|l|l|}
\hline
\rowcolor{popblueL}
\textbf{Personne} & \textbf{\esp{reciclar}} & \textbf{\esp{plantar}} & \textbf{\esp{proteger}} \\
\hline
\esp{yo} & \esp{reciclaré} & \esp{plantaré} & \esp{protegeré} \\
\esp{tú} & \esp{reciclarás} & \esp{plantarás} & \esp{protegerás} \\
\esp{él/ella/usted} & \esp{reciclará} & \esp{plantará} & \esp{protegerá} \\
\esp{nosotros/as} & \esp{reciclaremos} & \esp{plantaremos} & \esp{protegeremos} \\
\esp{vosotros/as} & \esp{reciclaréis} & \esp{plantaréis} & \esp{protegeréis} \\
\esp{ellos/ellas/ustedes} & \esp{reciclarán} & \esp{plantarán} & \esp{protegerán} \\
\hline
\end{tabular}
\end{center}

\textbf{Emplois du futur :}
\begin{itemize}
\item \textbf{Prévision / Certitude :} \esp{El lago Chad se secará si no llueve.} (Le lac Tchad séchera s'il ne pleut pas.)
\item \textbf{Promesse / Engagement :} \esp{Plantaremos diez árboles en el patio.} (Nous planterons dix arbres dans la cour.)
\end{itemize}

\textbf{Conséquence (si + présent, futur) :} \esp{Si contaminamos los ríos, no tendremos agua limpia.} « Si nous polluons les rivières, nous n'aurons pas d'eau propre. »

\textbf{Expressions pour présenter :} \esp{Buenos días, somos el grupo... ; Nuestro eslogan es... ; El primer problema es... ; Nos comprometemos a... ; ¡Gracias por su atención!}
\end{aretenir}

\begin{savaistu}[Le Cameroun face au changement climatique]
Le Cameroun est très vulnérable : le \textbf{lac Tchad} a perdu 90\% de sa surface en 50 ans (\esp{sequía}), la \textbf{forêt du bassin du Congo} (2e poumon vert mondial) subit la \esp{deforestación} pour le bois et l'agriculture, et \textbf{Douala} fait face à des \esp{inundaciones} récurrentes et à la \esp{contaminación} plastique. La \esp{Gran Muralla Verde} panafricaine traverse le nord du pays pour lutter contre la désertification. La Guinée équatoriale, pays hispanophone voisin, partage ces défis forestiers et côtiers.
\end{savaistu}

\courssub{Proposition de production et corrigé commenté}

\begin{exoresolu}[Modèle d'affiche et de discours du groupe « Los Verdes »]
Énoncé : rédige l'affiche complète d'un groupe (slogan, trois problèmes, trois engagements) et le discours de présentation d'une minute.
\tcblower
\textbf{Contenu de l'affiche :}

\esp{SI CUIDAMOS EL PLANETA HOY, VIVIREMOS MEJOR MAÑANA}

\esp{Los problemas :}

\esp{1. La deforestación destruye los bosques del Congo.}

\esp{2. La contaminación del plástico ensucia las calles de Douala.}

\esp{3. La sequía seca el lago Chad.}

\esp{Nuestros compromisos :}

\esp{1. Reciclaremos el plástico y el papel en el liceo.}

\esp{2. Plantaremos veinte árboles en el patio.}

\esp{3. Usaremos menos agua y menos electricidad.}

\textbf{Discours de présentation :}

\esp{Buenos días a todos. Somos el grupo « Los Verdes ». Nuestro eslogan es : « Si cuidamos el planeta hoy, viviremos mejor mañana ». En nuestro cartel hay tres problemas : la deforestación de los bosques del Congo, la contaminación del plástico en Douala y la sequía del lago Chad. Pero hay soluciones : nosotros reciclaremos el plástico y el papel, plantaremos veinte árboles en el patio del liceo y usaremos menos agua. Si todos participamos, el planeta será más verde. ¡Voten por nosotros! ¡Gracias por su atención!}

\emph{Traduction du discours :} « Bonjour à tous. Nous sommes le groupe "Les Verts". Notre slogan est : "Si nous prenons soin de la planète aujourd'hui, nous vivrons mieux demain". Sur notre affiche, il y a trois problèmes : la déforestation des forêts du Congo, la pollution du plastique à Douala et la sécheresse du lac Tchad. Mais il y a des solutions : nous recyclerons le plastique et le papier, nous planterons vingt arbres dans la cour du lycée et nous utiliserons moins d'eau. Si nous participons tous, la planète sera plus verte. Votez pour nous ! Merci de votre attention ! »

\textbf{Commentaire des choix de langue :}
\begin{itemize}
\item Le slogan applique la structure de la conséquence : \esp{si + presente} (\esp{cuidamos}) \esp{+ futuro} (\esp{viviremos}). C'est la structure grammaticale centrale de la leçon.
\item Les problèmes sont au présent de vérité générale : \esp{destruye, ensucia, seca}. Le présent convient pour décrire des faits actuels et permanents.
\item Les engagements sont tous au futur simple à la première personne du pluriel (\esp{reciclaremos, plantaremos, usaremos}) : c'est l'emploi du futur de promesse, parfait pour un manifeste écocitoyen.
\item Le discours se termine par une phrase de conséquence motivante (\esp{Si todos participamos, el planeta será más verde}) et deux formules de communication orale (\esp{¡Voten por nosotros! ¡Gracias por su atención!}).
\item Le lexique de la leçon est entièrement mobilisé : \esp{deforestación, contaminación, sequía, reciclar, planeta}.
\end{itemize}
\end{exoresolu}

\courssub{Bilan de l'activité}

\begin{aretenir}[Ce que j'ai appris à faire]
Dans cette activité, j'ai :
\begin{itemize}
\item lu et compris une annonce authentique en espagnol ;
\item réutilisé le lexique du climat dans un contexte réel et camerounais (lac Tchad, forêt du Congo, Douala) ;
\item conjugué le futur simple pour exprimer des engagements : \esp{reciclaremos, plantaremos, protegeremos} ;
\item construit des slogans avec \esp{si + presente, futuro} sans mettre de futur après \esp{si} ;
\item pris la parole en public pendant une minute et justifié un vote.
\end{itemize}
Auto-évaluation : je sais citer cinq mots du lexique du climat, conjuguer trois verbes au futur simple (dont un irrégulier) et écrire un slogan correct avec \esp{si}. Si un point reste fragile, je revois la partie grammaire de la leçon avant l'évaluation.
\end{aretenir}

\newpage

\coursec{Méthodes et exercices résolus}

\coursec{La protección del medioambiente y el cambio climático}

\courssub{Texte support : El planeta en peligro}
\begin{textebox}[Texte original — Niveau A2+/B1 — \textasciitilde 170 mots]
Actualmente, el calentamiento global es el gran reto de la humanidad. La deforestación de la selva del Congo, la contaminación industrial y el uso de energías no renovables aumentan el efecto invernadero. Como consecuencia, las temperaturas suben y el clima cambia.

En Camerún, los efectos son visibles. El lago Chad, compartido con Chad, Níger y Nigeria, se reduce año tras año; ha perdido más del 90\% de su agua. En el Norte, la sequía prolongada destruye las cosechas de mijo y maíz. En Douala, las inundaciones destruyen casas cerca del río Wouri.

Si no actuamos ahora, la situación empeorará. Perderemos biodiversidad y habrá más conflictos por el agua. Pero hay soluciones: debemos proteger los bosques, reciclar los residuos y usar energías renovables como la solar. Si cambiamos nuestros hábitos, el futuro será mejor. ¡Actuemos ya!
\end{textebox}

\begin{exemplebox}[Traduction et glossaire des mots clés]
\textbf{Traduction :} Actuellement, le réchauffement climatique est le grand défi de l'humanité. La déforestation de la forêt du Congo, la pollution industrielle et l'utilisation d'énergies non renouvelables augmentent l'effet de serre. Conséquence : les températures montent et le climat change.

Au Cameroun, les effets sont visibles. Le lac Tchad, partagé avec le Tchad, le Niger et le Nigeria, se réduit année après année ; il a perdu plus de 90\% de son eau. Au Nord, la sécheresse prolongée détruit les récoltes de mil et de maïs. À Douala, les inondations détruisent les maisons près du fleuve Wouri.

Si nous n'agissons pas maintenant, la situation empire. Nous perdrons la biodiversité et il y aura plus de conflits pour l'eau. Mais il y a des solutions : nous devons protéger les forêts, recycler les déchets et utiliser des énergies renouvelables comme le solaire. Si nous changeons nos habitudes, l'avenir sera meilleur. Agissons dès maintenant !

\textbf{Glossaire :} \esp{el calentamiento global} (le réchauffement climatique) ; \esp{la deforestación} (la déforestation) ; \esp{la selva} (la jungle, forêt tropicale) ; \esp{la contaminación} (la pollution) ; \esp{el efecto invernadero} (l'effet de serre) ; \esp{la sequía} (la sécheresse) ; \esp{la inundación} (l'inondation) ; \esp{la biodiversidad} (la biodiversité) ; \esp{reciclar} (recycler) ; \esp{las energías renovables} (les énergies renouvelables) ; \esp{el residuo} (le déchet) ; \esp{el hábito} (l'habitude).
\end{exemplebox}

\courssub{Compréhension du texte}
\begin{exercice}[Questions de compréhension — \textit{El planeta en peligro}]
\textbf{Consigne :} Répondez en espagnol, phrases complètes.
\begin{enumerate}
\item \textbf{Lecture globale :} Quel est le thème principal ? Quel est le ton du texte (alarmiste, informatif, mobilisateur) ? Citez un connecteur logique et un marqueur temporel.
\item \textbf{Compréhension explicite :} Citez les \textbf{trois causes principales} du réchauffement climatique mentionnées dans le premier paragraphe.
\item \textbf{Compréhension implicite + Contexte :} Le texte affirme : \esp{« El lago Chad se reduce año tras año »}. Expliquez cette phrase en la reliant au contexte camerounais (2-3 lignes).
\item \textbf{Lexique :} Trouvez dans le texte les synonymes ou définitions pour : \textit{pollution} ; \textit{forêt tropicale} ; \textit{énergies propres} ; \textit{perdre de l'eau/volume}.
\item \textbf{Grammaire en contexte :} Relevez \textbf{trois verbes au futur simple} dans le texte et donnez leur infinitif.
\item \textbf{Expression écrite (60-80 mots) :} \esp{¿Qué gestos cotidianos puedes hacer tú para proteger el medioambiente en tu ciudad (Yaoundé, Douala, etc.) ? Utiliza el futuro simple y conectores (Primero, Además, Si...).}
\end{enumerate}
\end{exercice}

\begin{corrige}[Exercice Compréhension]
\begin{enumerate}
\item \textbf{Global :} Thème : changement climatique et solutions. Ton : informatif et mobilisateur (appel à l'action final). Connecteur : \esp{Como consecuencia, Pero, Si}. Marqueur temporel : \esp{Actualmente, ahora, año tras año}.
\item \textbf{Causes (explicites) :} 1. \esp{La deforestación de la selva del Congo.} 2. \esp{La contaminación industrial.} 3. \esp{El uso de energías no renovables.}
\item \textbf{Implicite + Contexte :} \esp{El lago Chad, compartido por Camerún, Chad, Níger y Nigeria, ha perdido más del 90\% de su superficie en 50 años. El calentamiento global provoca sequías prolongadas y la evaporación del agua; además, la sobreexplotación para la agricultura y la ganadería agrava el problema. Esto causa conflictos entre pueblos y migración forzada.}
\item \textbf{Lexique :} \esp{contaminación} (pollution) ; \esp{selva} (forêt tropicale) ; \esp{energías renovables} (énergies propres) ; \esp{se reduce} (perdre du volume/rétrécir).
\item \textbf{Futur simple :} \esp{empeorará} (inf. \esp{empeorar}) ; \esp{perderemos} (inf. \esp{perder}) ; \esp{habrá} (inf. \esp{haber} — impersonnel, forme irrégulière). \textit{Aussi :} \esp{cambiará} (inf. \esp{cambiar}) ; \esp{será} (inf. \esp{ser}).
\item \textbf{Production (modèle) :} \esp{En mi ciudad, Douala, yo reciclaré las botellas de plástico y no tiraré basura en los arroyos. Además, usaré el transporte compartido (\"bendskin\" o taxi) en vez de moto individual para reducir la contaminación. También plantaré un árbol en el patio de mi escuela. Si todos hacemos esto, el aire será más limpio.}
\end{enumerate}
\end{corrige}

\courssub{Lexique thématique : El medioambiente}
\begin{definition}[Vocabulaire de base — Noms, verbes, adjectifs]
\begin{center}
\begin{tabularx}{\linewidth}{|X|X|X|}\hline
\rowcolor{popblueL} \textbf{Problèmes (Sustantivos)} & \textbf{Causes / Actions négatives (Verbos)} & \textbf{Solutions / Actions positives (Verbos)} \\ \hline
\esp{el calentamiento global} & \esp{contaminar} (polluer) & \esp{reciclar} (recycler) \\ \hline
\esp{la contaminación} (del aire, agua, suelo) & \esp{talar} (abattre des arbres) & \esp{reutilizar} (réutiliser) \\ \hline
\esp{la sequía} & \esp{quemar} (brûler \textit{combustibles fósiles}) & \esp{reducir} (réduire \textit{residuos}) \\ \hline
\esp{la inundación} & \esp{desperdiciar} (gaspiller \textit{agua, energía}) & \esp{plantar} (planter \textit{árboles}) \\ \hline
\esp{la deforestación} & \esp{emitir} (émettre \textit{gases}) & \esp{ahorrar} (économiser \textit{energía, agua}) \\ \hline
\esp{la desertificación} &  & \esp{usar energías renovables} (solar, eólica) \\ \hline
\esp{el efecto invernadero} &  & \esp{proteger} (protéger) \\ \hline
\end{tabularx}
\end{center}
\vspace{0.3cm}
\textbf{Adjectifs utiles :} \esp{renovable} (invariable au sg, \esp{-s} au pl), \esp{sostenible}, \esp{ecológico}, \esp{contaminado}, \esp{sequioso} (aride), \esp{inundado}.
\end{definition}

\begin{aretenir}[Pièges orthographiques et faux-amis — \textbf{À apprendre par cœur}]
\begin{itemize}
\item \textbf{\esp{Contaminación}} = pollution (environnement). \textit{Pas} \esp{*polución} (anglicisme rare). \esp{Infección} = infection (médical).
\item \textbf{\esp{Seqüía}} : accent sur le \cle{í} (hiatus \esp{u-í} \rightarrow \esp{ú-i} tonique). \esp{Inundación} : pas d'accent sur le \esp{u}.
\item \textbf{\esp{Bosque}} = forêt (tempérée) ; \textbf{\esp{selva}} = jungle / forêt tropicale humide (Congo, Amazonie).
\item \textbf{\esp{Medioambiente}} (un mot, masc.) ou \esp{medio ambiente} (deux mots, admis par la RAE). Préférer un mot.
\item \textbf{Accords :} \esp{energía renovable} (sg) $\rightarrow$ \esp{energ\textbf{í}as renov\textbf{ables}} (pl). L'adjectif en \cle{-e} prend \cle{-s} au pluriel, pas \cle{-es}.
\item \textbf{\esp{Haber} impersonnel :} \esp{Hay} (présent) $\rightarrow$ \esp{Habrá} (futur). Invariable. \textit{Jamais} \esp{*Habrán} ou \esp{*Habremos} pour « il y aura / il y avait ».
\end{itemize}
\end{aretenir}

\courssub{Grammaire : Le futur simple (Futuro Imperfecto)}
\begin{propriete}[Formation, irréguliers et emplois — Niveau A2+/B1]
\textbf{1. Formation (réguliers) :} \cle{Infinitif} + \cle{terminaisons uniformes} (3 groupes). \textbf{Accent tonique sur la dernière syllabe} (sauf \esp{nosotros} : \esp{-emos}).
\begin{center}
\begin{tabular}{|l|c|c|c|}\hline
\rowcolor{popblueL} \textbf{Personne} & \textbf{-AR (hablar)} & \textbf{-ER (comer)} & \textbf{-IR (vivir)} \\ \hline
Yo & hablar\textbf{é} & comer\textbf{é} & vivir\textbf{é} \\ \hline
Tú & hablar\textbf{ás} & comer\textbf{ás} & vivir\textbf{ás} \\ \hline
Él/Ella/Ud. & hablar\textbf{á} & comer\textbf{á} & vivir\textbf{á} \\ \hline
Nosotros & hablar\textbf{emos} & comer\textbf{emos} & vivir\textbf{emos} \\ \hline
Vosotros & hablar\textbf{éis} & comer\textbf{éis} & vivir\textbf{éis} \\ \hline
Ellos/Uds. & hablar\textbf{án} & comer\textbf{án} & vivir\textbf{án} \\ \hline
\end{tabular}
\end{center}

\textbf{2. Irréguliers fréquents (radical modifié + terminaisons régulières) :} À mémoriser par familles.
\begin{center}
\begin{tabularx}{\linewidth}{|X|X|X|}\hline
\rowcolor{popblueL} \textbf{Infinitif} & \textbf{Radical futur} & \textbf{Exemple 1ère sg / 3ème pl} \\ \hline
\esp{Caber} & \esp{cabr-} & \esp{cabré / cabrán} \\ \hline
\esp{Poner} & \esp{pondr-} & \esp{pondré / pondrán} \\ \hline
\esp{Saber} & \esp{sabr-} & \esp{sabré / sabrán} \\ \hline
\esp{Salir} & \esp{saldr-} & \esp{saldré / saldrán} \\ \hline
\esp{Tener} & \esp{tendr-} & \esp{tendré / tendrán} \\ \hline
\esp{Valer} & \esp{valdr-} & \esp{valdré / valdrán} \\ \hline
\esp{Venir} & \esp{vendr-} & \esp{vendré / vendrán} \\ \hline
\esp{Haber} & \esp{habr-} & \esp{habré / habrán} (\textit{impers.} \rightarrow \esp{habrá}) \\ \hline
\esp{Hacer} & \esp{har-} & \esp{haré / harán} \\ \hline
\esp{Querer} & \esp{querr-} & \esp{querré / querrán} \\ \hline
\esp{Decir} & \esp{dir-} & \esp{diré / dirán} \\ \hline
\esp{Poder} & \esp{podr-} & \esp{podré / podrán} \\ \hline
\end{tabularx}
\end{center}

\textbf{3. Emplois (niveau A2+/B1) :}
\begin{itemize}
\item Action future certaine / programmée : \esp{Mañana plantaremos árboles en el colegio.}
\item Prédiction / Hypothèse probable : \esp{El clima cambiará mucho en 20 años.}
\item Promesse / Engagement : \esp{Te prometo que reciclaré todo el plástico.}
\item \textbf{Attention :} \textbf{Pas pour intention immédiate} (proche futur) $\rightarrow$ utiliser \esp{ir a + infinitif} : \esp{Voy a reciclar ahora.}
\end{itemize}
\end{propriete}

\begin{exoresolu}[Futur simple : Conjugaison, Transformation, Thème]
\textbf{Énoncé :}
\begin{enumerate}
\item \textbf{Conjuguez au futur simple} (1ère pers. sg. \textit{yo} et 3ème pers. pl. \textit{ellos/ustedes}) : \esp{contaminar, poder, haber, reciclar, hacer, secar}.
\item \textbf{Transformez le présent en futur simple} :
\begin{enumerate}
\item \esp{La fábrica contamina el río.}
\item \esp{Podemos salvar el bosque.}
\item \esp{Hay mucha basura en la calle.}
\end{enumerate}
\item \textbf{Thème (Français $\rightarrow$ Espagnol) :} Traduisez en utilisant le futur simple.
\begin{enumerate}
\item Nous protégerons la forêt tropicale.
\item Il y aura plus de sécheresses si nous ne changeons pas.
\item Pourras-tu (tú) venir planter des arbres samedi ?
\end{enumerate}
\end{enumerate}
\tcblower
\textbf{Solution détaillée :}
\begin{enumerate}
\item \textbf{Conjugaison (vérification radicaux) :}
\begin{center}
\begin{tabular}{|l|c|c|}\hline
\rowcolor{popblueL} \textbf{Verbe} & \textbf{Yo} & \textbf{Ellos/Uds.} \\ \hline
\esp{contaminar} (reg.) & \esp{contaminar\textbf{é}} & \esp{contaminar\textbf{án}} \\ \hline
\esp{poder} (irr. \cle{podr-}) & \esp{podr\textbf{é}} & \esp{podr\textbf{án}} \\ \hline
\esp{haber} (irr. \cle{habr-}) & \esp{habr\textbf{é}} & \esp{habr\textbf{án}} \\ \hline
\esp{reciclar} (reg.) & \esp{reciclar\textbf{é}} & \esp{reciclar\textbf{án}} \\ \hline
\esp{hacer} (irr. \cle{har-}) & \esp{har\textbf{é}} & \esp{har\textbf{án}} \\ \hline
\esp{secar} (reg.) & \esp{secar\textbf{é}} & \esp{secar\textbf{án}} \\ \hline
\end{tabular}
\end{center}
\item \textbf{Transformation (présent $\rightarrow$ futur) :}
\begin{enumerate}
\item \esp{La fábrica contaminar\textbf{á} el río.} (\esp{contamina} 3sg $\rightarrow$ \esp{contaminará}).
\item \esp{Podr\textbf{emos} salvar el bosque.} (\esp{Podemos} 1pl irr. $\rightarrow$ \esp{Podremos}).
\item \esp{Habrá mucha basura en la calle.} (\esp{Hay} = \esp{Haber} 3sg présent $\rightarrow$ \esp{Habrá} 3sg futur). \textbf{Rappel :} \esp{Hay} n'existe qu'au présent ; futur = \esp{Habrá} (invariable).
\end{enumerate}
\item \textbf{Thème :}
\begin{enumerate}
\item \esp{Proteg\textbf{eremos} la selva tropical.} (\esp{Proteger} -er, 1pl $\rightarrow$ \esp{-emos}).
\item \esp{Habr\textbf{á} más sequías si no cambiamos.} (\esp{Habrá} 3sg ; \esp{sequías} pl. fém. ; \esp{si no cambiamos} = \cle{si + présent indicatif}).
\item \esp{\textbf{Podrás} venir a plantar árboles el sábado ?} (\esp{Poder} irr. \cle{podr-}, 2sg \esp{tú} $\rightarrow$ \esp{podrás} ; \esp{venir} + \esp{plantar} infinitifs).
\end{enumerate}
\end{enumerate}
\textbf{Point de vigilance :} Ne pas confondre \esp{Habrá} (futur) et \esp{Habría} (conditionnel). L'accent est sur le \cle{á} final.
\end{exoresolu}

\courssub{Grammaire : Structure de conséquence réelle « Si + Présent, Futur »}
\begin{propriete}[Règle, ordre des propositions et distinction avec l'irréel]
\textbf{Structure figée (Condition réelle future possible) :}
\cle{Si} + \textbf{Indicatif Présent} (la condition) $\rightarrow$ \textbf{Futur Simple} (la conséquence).

\textbf{Ordre des propositions :}
\begin{itemize}
\item \esp{Si + Presente, Futuro.} (Virgule si \esp{Si} en tête).
\item \esp{Futuro si + Presente.} (Pas de virgule).
\end{itemize}

\textbf{Distinction cruciale (Niveau Bac) :}
\begin{center}
\begin{tabularx}{\linewidth}{|X|X|X|}\hline
\rowcolor{popblueL} \textbf{Type de condition} & \textbf{Structure} & \textbf{Exemple} \\ \hline
\textbf{Réel futur (Cette leçon)} & \cle{Si + Présent Indicatif, Futur} & \esp{Si reciclas, ayudarás al planeta.} \\ \hline
Irréel présent (Niveau sup.) & \cle{Si + Imparfait Subjonctif, Conditionnel} & \esp{Si reciclaras, ayudarías al planeta.} \\ \hline
Irréel passé (Niveau sup.) & \cle{Si + Plus-que-parfait Subj., Cond. Passé} & \esp{Si hubieras reciclado, habrías ayudado.} \\ \hline
\end{tabularx}
\end{center}

\textbf{Interdictions formelles :}
\begin{enumerate}
\item \textbf{Jamais} de futur après \esp{Si} : \esp{*Si yo reciclaré...} $\times$
\item \textbf{Jamais} de conditionnel après \esp{Si} dans cette structure : \esp{*Si yo reciclaría...} $\times$
\end{enumerate}

\textbf{Verbes fréquents au présent dans la \esp{Si}-proposition (thème environnement) :} \esp{contaminar, talar, quemar, reciclar, ahorrar, usar, proteger, cambiar, actuar, hacer, separar (la basura), plantar.}
\end{propriete}

\begin{exoresolu}[Structure « Si + Présent, Futur » : Transformation et Production]
\textbf{Énoncé :}
\begin{enumerate}
\item \textbf{Associez} et conjuguez le verbe entre parenthèses au temps qui convient (Présent ou Futur).
\begin{enumerate}
\item Si nosotros \_\_\_\_\_\_ (proteger) los bosques, \_\_\_\_\_\_ (haber) más oxígeno.
\item Si la temperatura \_\_\_\_\_\_ (subir), los glaciares \_\_\_\_\_\_ (derretirse).
\item Si tú \_\_\_\_\_\_ (no / reciclar), el vertedero \_\_\_\_\_\_ (crecer).
\end{enumerate}
\item \textbf{Complétez} la conséquence logique (Futur) pour ces conditions réelles au Cameroun :
\begin{enumerate}
\item Si llueve mucho en Douala en agosto, \_\_\_\_\_\_\_\_\_\_\_\_.
\item Si talamos los árboles del Monte Camerún, \_\_\_\_\_\_\_\_\_\_\_\_.
\item Si usamos energías renovables en los pueblos del Norte, \_\_\_\_\_\_\_\_\_\_\_\_.
\end{enumerate}
\item \textbf{Inversion :} Réécrivez la phrase 1a en commençant par la conséquence (sans \esp{si} au début).
\end{enumerate}
\tcblower
\textbf{Solution détaillée :}
\begin{enumerate}
\item \textbf{Association + Conjugaison (Règle : Si + Présent, Futur) :}
\begin{enumerate}
\item Si nosotros \textbf{protegemos} (présent ind. 1pl), \textbf{habrá} (futur 3sg \esp{hay}) más oxígeno.
\item Si la temperatura \textbf{sube} (présent ind. 3sg), los glaciares \textbf{se derretirán} (futur 3pl pronominal).
\item Si tú \textbf{no reciclas} (présent ind. 2sg), el vertedero \textbf{crecerá} (futur 3sg).
\end{enumerate}
\item \textbf{Production guidée (Exemples corrects attendus) :}
\begin{enumerate}
\item \esp{... habrá inundaciones en los barrios bajos.} (Vocabulaire : \esp{inundaciones}, \esp{barrios bajos}).
\item \esp{... habrá erosión del suelo y deslizamientos de tierra.} (\esp{erosión}, \esp{deslizamientos}).
\item \esp{... tendrán electricidad limpia y barata.} (\esp{tendrán} futur \esp{tener} irr. \cle{tendr-}, \esp{eléctricidad limpia}).
\end{enumerate}
\item \textbf{Inversion (ordre inverse, pas de virgule) :}
\esp{Habrá más oxígeno si nosotros protegemos los bosques.}
\end{enumerate}
\textbf{Justification :} La proposition introduite par \esp{Si} exprime une condition \textbf{réelle et future possible} $\rightarrow$ \textbf{Indicatif Présent}. La principale exprime la conséquence certaine $\rightarrow$ \textbf{Futur Simple}. L'inversion ne change pas les temps.
\end{exoresolu}

\courssub{Méthodologie : Commenter un texte et produire un texte argumentatif}
\begin{methode}[Commenter un texte de presse sur l'environnement]
\textbf{Objectif :} Répondre à des questions de compréhension (explicite, implicite, vocabulaire) et produire un commentaire personnel structuré en espagnol.
\begin{enumerate}
\item \textbf{Lecture globale :} Identifier la source (presse, ONG, blog), le thème (changement climatique, pollution, solutions), le ton (alarmiste, informatif, mobilisateur) et la structure (introduction, causes, conséquences, solutions, conclusion).
\item \textbf{Lecture détaillée :} Surligner les \cle{mots-clés} du thème (vocabulaire de la leçon), les \cle{connecteurs logiques} (\esp{por tanto, además, sin embargo, por ejemplo}) et les \cle{marqueurs temporels} (\esp{actualmente, en el futuro, dentro de 10 años}).
\item \textbf{Questions de compréhension :}
\begin{itemize}
\item \textit{Explicites :} La réponse est dans le texte (recopier ou reformuler).
\item \textit{Implicites :} Déduire le sens (pourquoi ? quelle intention ?).
\item \textit{Lexicales :} Expliquer un mot du texte par un synonyme, un antonyme ou une définition en espagnol.
\end{itemize}
\item \textbf{Commentaire personnel (Expresión escrita) :} Structurer en 3 parties :
\begin{enumerate}
\item \textit{Introduction :} Présenter le document + annoncer le point de vue (\esp{En mi opinión, este texto refleja la realidad de...}).
\item \textit{Développement (2 arguments) :} Argumenter avec \cle{connecteurs} (\esp{Primero, además, por un lado... por otro lado}) + exemples locaux (Cameroun, Guinée équatoriale) ou personnels.
\item \textit{Conclusion :} Ouvrir sur une action future (\esp{Por eso, debemos actuar ahora; si no cambiamos, el planeta sufrirá más}).
\end{enumerate}
\end{enumerate}
\end{methode}

\begin{methode}[Rédiger un paragraphe d'argumentation : « Protéger l'environnement »]
\textbf{Objectif :} Produire un texte cohérent (100-120 mots) proposant des solutions concrètes, en mobilisant lexique, futur simple et structure \esp{Si + Présent, Futur}.
\begin{enumerate}
\item \textbf{Analyser le sujet :} Type (article blog, discours, lettre au maire), public (jeunes, autorités), ton (engagé, persuasif). Mots-clés imposés : \esp{calentamiento global, reciclar, energías renovables, sequía, futuro}.
\item \textbf{Plan « 3 paragraphes / 3 arguments » :}
\begin{enumerate}
\item \textit{Diagnostic (Présent) :} \esp{Actualmente, nuestro planeta sufre... La contaminación y la deforestación causan...}
\item \textit{Solutions individuelles (Futur + Si+Pres/Fut) :} \esp{Primero, yo reciclaré... Si ahorramos agua, habrá más para todos.}
\item \textit{Solutions collectives / Appel à l'action (Futur + Devoir) :} \esp{Los gobiernos deben invertir en energías renovables. Si actuamos ahora, el futuro será mejor.}
\end{enumerate}
\item \textbf{Connecteurs logiques obligatoires :} \esp{Actualmente, Por eso, Por tanto, Además, Sin embargo, En conclusión, Para terminar.}
\item \textbf{Relecture (Check-list) :}
\begin{itemize}
\item Accents sur futurs (\esp{haré, habrá, podré}) et \esp{sequía, acción, nación}.
\item Accords (adj. fém. pl. : \esp{energías renovables}).
\item Structure \esp{Si + Pres, Fut} correcte (pas de futur après \esp{si}).
\item Pas de français déguisé (\esp{*proteger el medio ambiente} $\checkmark$, pas \esp{*proteger l'environnement}).
\end{itemize}
\end{enumerate}
\end{methode}

\courssub{Expression écrite guidée et exercices résolus}
\begin{exoresolu}[Compréhension et commentaire : « El planeta en peligro »]
\textbf{Énoncé :} À partir du texte support \textit{El planeta en peligro} (lu en classe), répondez en espagnol :
\begin{enumerate}
\item Quels sont les \textbf{trois causes principales} du réchauffement climatique citées dans le texte ?
\item Le texte affirme : \esp{« El lago Chad se reduce año tras año »}. Expliquez cette phrase en la reliant au contexte camerounais (2-3 lignes).
\item Trouvez dans le texte \textbf{trois verbes au futur simple} et donnez leur infinitif.
\item \textbf{Expression écrite (60-80 mots) :} \esp{¿Qué gestos cotidianos puedes hacer tú para proteger el medioambiente en tu ciudad (Yaoundé, Douala, etc.) ? Utiliza el futuro simple y conectores.}
\end{enumerate}
\tcblower
\textbf{Solution détaillée :}
\begin{enumerate}
\item \textbf{Causes (repérage explicite) :}
\begin{itemize}
\item La \cle{deforestación} (coupes abusives, agriculture itinérante) \textit{— texte : \esp{deforestación de la selva del Congo}}.
\item La \cle{contaminación} industrielle (\textit{texte : \esp{contaminación industrial}}).
\item L'utilisation d'\cle{énergies non renouvelables} (\textit{texte : \esp{uso de energías no renovables}}) au détriment des \esp{energías renovables}.
\end{itemize}
\item \textbf{Explication implicite + contexte :}
\esp{El lago Chad, compartido por Camerún, Chad, Níger y Nigeria, ha perdido más del 90\% de su superficie en 50 años. El calentamiento global provoca sequías prolongadas y la evaporación del agua; además, la sobreexplotación para la agricultura y la ganadería agrava el problema. Esto causa conflictos entre pueblos y migración forzada.}
\item \textbf{Verbes au futur (reconnaissance morphologique) :}
\begin{itemize}
\item \esp{empeorará} (inf. \esp{empeorar}) — radical \esp{empeor-} + \esp{-ará}.
\item \esp{perderemos} (inf. \esp{perder}) — radical \esp{perder-} + \esp{-emos}.
\item \esp{habrá} (inf. \esp{haber}) — forme irrégulière (radical \esp{habr-}), impersonnel.
\end{itemize}
\item \textbf{Production écrite (modèle) :}
\esp{En mi ciudad, Douala, yo reciclaré las botellas de plástico y no tiraré basura en los arroyos. Además, usaré el transporte compartido (\"bendskin\" o taxi) en vez de moto individual para reducir la contaminación. También plantaré un árbol en el patio de mi escuela. Si todos hacemos esto, el aire será más limpio.}
\end{enumerate}
\textbf{Justification grammaticale :} \esp{reciclaré, usaré, plantaré} = futur simple 1ère pers. sg. (réguliers). \esp{será} = futur irrégulier de \esp{ser}. \esp{Si todos hacemos..., el aire será...} = structure \cle{si + présent indicatif, futur} (condition réelle future).
\end{exoresolu}

\begin{exoresolu}[Lexique en contexte : Complétion et Thème]
\textbf{Énoncé :}
\begin{enumerate}
\item \textbf{Complétez} avec le mot juste (forme correcte) : \esp{calentamiento, contaminación, sequía, inundación, deforestación, renovables, reciclar}.
\begin{enumerate}
\item La \_\_\_\_\_\_ del aire en las grandes ciudades causa enfermedades respiratorias.
\item En el norte de Camerún, la \_\_\_\_\_\_ prolongada destruye las cosechas de mijo y maíz.
\item La \_\_\_\_\_\_ de la selva del Congo libera mucho CO₂ a la atmósfera.
\item Debemos invertir en energías \_\_\_\_\_\_ como la solar y la eólica.
\item Si separamos la basura, podremos \_\_\_\_\_\_ el papel, el vidrio y el plástico.
\end{enumerate}
\item \textbf{Thème (Français $\rightarrow$ Espagnol) :} Traduisez en utilisant le vocabulaire ci-dessus.
\begin{enumerate}
\item Le réchauffement climatique est un danger pour la planète.
\item Les inondations détruisent les maisons près du fleuve Wouri.
\item Nous devons protéger l'environnement pour les générations futures.
\end{enumerate}
\end{enumerate}
\tcblower
\textbf{Solution détaillée :}
\begin{enumerate}
\item \textbf{Vocabulaire en contexte (accords genre/nombre) :}
\begin{enumerate}
\item \esp{contaminación} (fém. sg., suit \esp{La}).
\item \esp{sequía} (fém. sg., accent sur \cle{í} ! \esp{La sequía}).
\item \esp{deforestación} (fém. sg.).
\item \esp{renovables} (fém. pl., accord avec \esp{energías}).
\item \esp{reciclar} (infinitif après \esp{podremos}).
\end{enumerate}
\item \textbf{Thème (traduction vers l'espagnol) :}
\begin{enumerate}
\item \esp{El calentamiento global es un peligro para el planeta.} (\cle{calentamiento global} = terme figé ; \esp{peligro} + \esp{para}).
\item \esp{Las inundaciones destruyen las casas cerca del río Wouri.} (\esp{inundación} $\rightarrow$ pl. \esp{inundaciones} ; \esp{cerca de} + \esp{el} = \esp{del}).
\item \esp{Debemos proteger el medioambiente para las generaciones futuras.} (\esp{Debemos} = \esp{deber} 1ère pl. présent obligation ; \esp{proteger} infinitif ; \esp{medioambiente} masc. sg.).
\end{enumerate}
\end{enumerate}
\textbf{Pièges évités :} \esp{sequía} (accent), \esp{inundaciones} (pluriel en \cle{-es}), \esp{del} (contraction), \esp{medioambiente} (un mot).
\end{exoresolu}

\begin{exoresolu}[Expression écrite guidée : Article pour le journal du lycée]
\textbf{Énoncé :} Vous écrivez un article pour le journal de votre lycée (\esp{« El Eco Estudiantil »}) intitulé : \esp{« Jóvenes cameruneses por el clima : ¡Actuemos ya ! »} (100-120 mots).
\textbf{Consignes :}
\begin{itemize}
\item Décrire un problème local (Douala/Yaoundé/Nord).
\item Proposer 2 gestes quotidiens (futur simple \esp{yo/nosotros}).
\item Utiliser \textbf{au moins deux phrases} avec \esp{Si + Présent, Futur}.
\item Terminer par un appel à l'action collectif.
\end{itemize}
\tcblower
\textbf{Production type (Modèle corrigé - 115 mots) :}
\begin{textebox}[Modèle de rédaction — \textit{El Eco Estudiantil}]
Actualmente, Camerún sufre los efectos del calentamiento global. En el Norte, la sequía destruye las cosechas; en Douala, la contaminación del río Wouri y las inundaciones amenazan la salud. Por eso, nosotros, los jóvenes, debemos actuar ya.

Primero, yo reciclaré todo el plástico y el papel en mi casa y en el colegio. Además, usaré mi botella reutilizable para no comprar botellas de un solo uso. Si separamos la basura, los vertederos no crecerán tan rápido.

Segundo, plantaremos árboles en los patios de las escuelas. Si protegemos los bosques como el de Korup, frenaremos la deforestación.

En conclusión, pido al gobierno que invierta en energías renovables. ¡Juntos construiremos un futuro sostenible!
\end{textebox}
\textbf{Analyse grammaticale du modèle :}
\begin{itemize}
\item \textbf{Lexique :} \esp{calentamiento global, sequía, contaminación, inundaciones, deforestación, energías renovables, futuro sostenible}.
\item \textbf{Futur simple (1ère pers. sg/pl) :} \esp{reciclaré, usaré, plantaremos, construiremos}.
\item \textbf{Si + Présent, Futur (x2) :} \esp{Si separamos..., no crecerán...} ; \esp{Si protegemos..., frenaremos...}.
\item \textbf{Connecteurs :} \esp{Actualmente, Por eso, Primero, Además, Segundo, En conclusión}.
\item \textbf{Contexte camerounais :} Nord (sécheresse), Douala (Wouri, inondations), Korup (forêt), jeunes/école.
\end{itemize}
\end{exoresolu}

\courssub{Points de vigilance pour l'évaluation}
\begin{attention}[Les 5 erreurs qui coûtent des points au Bac]
\begin{enumerate}
\item \textbf{Accent manquant sur le futur :} \esp{*reciclaré} $\rightarrow$ \esp{reciclar\textbf{é}} ; \esp{*habra} $\rightarrow$ \esp{habr\textbf{á}} ; \esp{*podra} $\rightarrow$ \esp{podr\textbf{á}}. \textit{Règle : Accent tonique sur la dernière voyelle (sauf nosotros -emos).}
\item \textbf{Futur après \esp{Si} (Calque français « Si je ferai... ») :} \esp{*Si yo reciclaré, ayudaré} $\times$ $\rightarrow$ \esp{Si yo \textbf{reciclo}, ayudaré} $\checkmark$. \textit{Jamais de futur dans la proposition conditionnelle introduite par \esp{Si}.}
\item \textbf{Faux-ami \esp{Contaminación} :} Ne pas écrire \esp{*polución} (anglicisme/peu usité) pour « pollution » environnemental ; le mot standard est \esp{contaminación}. Attention : \esp{infección} = infection (médical).
\item \textbf{Accord adjectif \esp{Renovable} :} \esp{*energía renovable} (sg) mais \esp{energ\textbf{í}as renov\textbf{ables}} (pl). \textit{L'adjectif en \cle{-e} prend \cle{-s} au pluriel, pas \cle{-es}.}
\item \textbf{Verbe \esp{Haber} (hay) au futur :} \esp{Habrá} $\checkmark$ (il y aura) mais \esp{Habrá muchos árboles} (invariable). \textit{Piège :} \esp{Hay} (présent) $\rightarrow$ \esp{Habrá} (futur). Ne pas conjuguer \esp{Haber} comme un verbe normal (\esp{*Habremos} $\times$ pour « il y aura »).
\end{enumerate}
\end{attention}

\newpage

\coursec{Exercices}

\courssub{Je vérifie mes connaissances}

\begin{exercice}[QCM : le lexique du climat]
Entoure la bonne réponse.
\begin{enumerate}
\item \esp{El calentamiento global} signifie :
\begin{enumerate}
\item le refroidissement de la planète
\item l'augmentation de la température de la Terre
\item la protection des forêts
\end{enumerate}
\item Une \esp{sequía} est :
\begin{enumerate}
\item une longue période sans pluie
\item une grande quantité d'eau
\item une tempête de sable
\end{enumerate}
\item \esp{Reciclar} veut dire :
\begin{enumerate}
\item jeter les déchets dans la rue
\item brûler les ordures
\item réutiliser les déchets pour fabriquer de nouveaux produits
\end{enumerate}
\item Les \esp{energías renovables} sont par exemple :
\begin{enumerate}
\item le pétrole et le charbon
\item le soleil et le vent
\item le bois et le gaz
\end{enumerate}
\end{enumerate}
\end{exercice}

\begin{exercice}[Vrai ou faux ?]
Dis si les phrases suivantes sont vraies (V) ou fausses (F) et justifie ta réponse par une courte phrase en espagnol.
\begin{enumerate}
\item \esp{La deforestación es buena para los animales.}
\item \esp{Las inundaciones pueden destruir las casas y los cultivos.}
\item \esp{El lago Chad crece cada año gracias a las lluvias.}
\item \esp{Reciclar los residuos ayuda a proteger el medioambiente.}
\end{enumerate}
\end{exercice}

\begin{exercice}[Le bon mot]
Complète chaque phrase avec le mot qui convient : \esp{contaminación}, \esp{deforestación}, \esp{inundación}, \esp{sequía}, \esp{reciclar}.
\begin{enumerate}
\item À Douala, après les grandes pluies, il y a souvent une \ldots dans certains quartiers.
\item La \ldots de la forêt du bassin du Congo menace les gorilles et les éléphants.
\item Dans l'Extrême-Nord du Cameroun, la \ldots rend l'agriculture très difficile.
\item La \ldots de l'air vient des voitures, des usines et des feux de brousse.
\item Pour protéger la planète, nous devons \ldots le plastique et le papier.
\end{enumerate}
\end{exercice}

\begin{exercice}[Conjugaison : le futur simple]
Conjugue les verbes entre parenthèses au futur simple.
\begin{enumerate}
\item Mañana, nosotros \ldots (plantar) diez árboles en el patio del liceo.
\item Si no llueve, los campesinos no \ldots (poder) cultivar el maíz.
\item ¿Tú \ldots (venir) a la manifestación por el clima el sábado ?
\item El próximo año \ldots (haber) menos plástico en nuestra ciudad.
\item Los científicos \ldots (decir) la verdad sobre el clima.
\end{enumerate}
\end{exercice}

\courssub{Je m'entraîne}

\begin{exercice}[Thème : traduis en espagnol]
Traduis les phrases suivantes en espagnol. Attention aux futurs et aux phrases avec \esp{si}.
\begin{enumerate}
\item Si nous ne protégeons pas la forêt, elle disparaîtra.
\item Il y aura plus de sécheresses et les animaux souffriront.
\item Nous planterons des arbres et nous recyclerons les déchets.
\item Si tu éteins la lumière, tu économiseras de l'énergie.
\end{enumerate}
\end{exercice}

\begin{exercice}[Transforme les phrases]
Transforme chaque phrase en une phrase de conséquence avec \esp{si + presente, futuro}, comme dans le modèle.\par
\emph{Modèle :} \esp{Contaminas el río. Los peces mueren.} $\rightarrow$ \esp{Si contaminas el río, los peces morirán.}
\begin{enumerate}
\item Cortamos los árboles. Los animales pierden su casa.
\item Usas el coche todos los días. Contaminas el aire.
\item Reciclamos el plástico. Protegemos el medioambiente.
\item No llueve durante meses. Hay una sequía terrible.
\end{enumerate}
\end{exercice}

\begin{exercice}[Texte à trous]
Complète le texte suivant avec les mots de la liste : \esp{calentamiento global}, \esp{sequía}, \esp{inundación}, \esp{energías renovables}, \esp{reciclar}, \esp{deforestación}.
\begin{textebox}[Un texto para completar]
El \ldots\ es un problema muy grave para nuestro planeta : las temperaturas suben cada año. En el norte de Camerún, la \ldots\ destruye las cosechas porque no llueve durante muchos meses. En cambio, en las ciudades del sur, una \ldots\ puede cubrir las calles de agua después de las lluvias. La \ldots\ de los bosques mata a muchos animales. Para luchar contra estos problemas, debemos usar \ldots\ como el sol y el viento, y también debemos \ldots\ nuestros residuos todos los días.
\end{textebox}
\end{exercice}

\begin{exercice}[Appariements]
Relie chaque mot espagnol à sa traduction française.
\begin{center}
\begin{tabularx}{\linewidth}{|X|X|}
\hline
\rowcolor{popblueL} \textbf{Español} & \textbf{Français} \\
\hline
1. \esp{la sequía} & a. les énergies renouvelables \\
\hline
2. \esp{la inundación} & b. recycler \\
\hline
3. \esp{la deforestación} & c. la sécheresse \\
\hline
4. \esp{las energías renovables} & d. l'inondation \\
\hline
5. \esp{reciclar} & e. la déforestation \\
\hline
\end{tabularx}
\end{center}
\end{exercice}

\courssub{Je me prépare au Bac}

\begin{exercice}[Compréhension écrite : la forêt amazonienne]
Lis le texte suivant, puis réponds aux questions.
\begin{textebox}[La Amazonía respira con dificultad]
La Amazonía es el bosque más grande del mundo. Millones de animales y de plantas viven allí, y muchas familias indígenas dependen de sus árboles y de sus ríos. Pero hoy, el bosque está en peligro. Cada año, los hombres cortan miles de árboles para cultivar la tierra, criar vacas o buscar oro. Esta deforestación tiene consecuencias terribles : los animales pierden su casa, el aire se contamina y el clima del planeta cambia.

Los científicos dicen que la situación es urgente. Si no protegemos la Amazonía, desaparecerán muchas especies de plantas y de animales. Además, habrá más sequías y más incendios en la región. Algunos gobiernos y organizaciones ya trabajan para salvar el bosque : plantan árboles, crean parques naturales y educan a los jóvenes.

Todos podemos ayudar. Si consumimos menos carne y menos papel, reduciremos la presión sobre el bosque. Si reciclamos y usamos energías renovables, protegeremos el planeta entero. La Amazonía es el pulmón de la Tierra : salvarla es salvar nuestra propia vida.
\end{textebox}
\textbf{Glossaire :} \esp{criar vacas} = élever des vaches ; \esp{el oro} = l'or ; \esp{el incendio} = l'incendie ; \esp{el pulmón} = le poumon.\par
\vspace{0.3em}
\textbf{I. Vrai ou faux ?} Justifie chaque réponse par une citation exacte du texte.
\begin{enumerate}
\item La Amazonía es un bosque pequeño del sur de Europa.
\item Los hombres cortan árboles solamente para construir escuelas.
\item Algunas organizaciones ya trabajan para salvar el bosque.
\end{enumerate}
\textbf{II. Questions ouvertes.} Réponds en espagnol, avec des phrases complètes.
\begin{enumerate}
\item ¿Por qué está en peligro la Amazonía ? Da dos razones.
\item Según el texto, ¿qué pasará si no protegemos la Amazonía ?
\item ¿Qué podemos hacer todos para ayudar al bosque ?
\end{enumerate}
\textbf{III. Questions de langue.}
\begin{enumerate}
\item Trouve dans le texte deux verbes conjugués au futur simple et recopie les phrases.
\item Trouve dans le texte une phrase avec \esp{si} et explique sa construction.
\end{enumerate}
\end{exercice}

\begin{exercice}[Thème et version]
\textbf{A. Thème.} Traduis en espagnol le texte suivant (attention aux futurs et à la phrase avec \esp{si}) :\par
\vspace{0.3em}
« Si nous ne protégeons pas la forêt, elle disparaîtra. Il y aura plus de sécheresses et les animaux souffriront. Nous planterons des arbres et nous recyclerons les déchets. »\par
\vspace{0.5em}
\textbf{B. Version.} Traduis en français le texte suivant.
\begin{textebox}[El lago Chad]
El lago Chad, en el norte de Camerún, es cada vez más pequeño. Los científicos dicen que, si no llueve más, el lago desaparecerá casi totalmente. Los pescadores no podrán trabajar y las familias sufrirán. Por eso, los países de la región plantarán árboles y protegerán el agua del lago.
\end{textebox}
\end{exercice}

\begin{exercice}[Production écrite type Bac]
Rédige un texte en espagnol de \textbf{80 à 100 mots} sur le sujet suivant :\par
\vspace{0.3em}
\esp{¿Qué harás tú y tu familia para proteger el medioambiente?}\par
\vspace{0.3em}
\textbf{Consignes obligatoires :}
\begin{itemize}
\item utilise au moins \textbf{trois mots du lexique} de la leçon (\esp{reciclar}, \esp{energías renovables}, \esp{contaminación}, \esp{sequía}, \esp{deforestación}\ldots) ;
\item emploie au moins \textbf{trois verbes au futur simple} ;
\item écris au moins \textbf{deux phrases avec \esp{si}} (si + présent, futur) ;
\item soigne l'orthographe et les accents.
\end{itemize}
\end{exercice}

\newpage

\coursec{Corrigés des exercices}

\begin{corrige}[Exercice 1 — QCM : le lexique du climat]
\begin{enumerate}
\item Réponse \textbf{b} : \esp{El calentamiento global} signifie l'augmentation de la température de la Terre. \esp{Calentar} = chauffer ; \esp{calentamiento} = réchauffement.
\item Réponse \textbf{a} : une \esp{sequía} est une longue période sans pluie. Attention au faux ami : \esp{sequía} ne veut pas dire « séquence ».
\item Réponse \textbf{c} : \esp{reciclar} = réutiliser les déchets pour fabriquer de nouveaux produits.
\item Réponse \textbf{b} : les \esp{energías renovables} sont par exemple le soleil (\esp{energía solar}) et le vent (\esp{energía eólica}). Le pétrole et le charbon sont des énergies fossiles (\esp{energías fósiles}).
\end{enumerate}
\end{corrige}

\begin{corrige}[Exercice 2 — Vrai ou faux ?]
\begin{enumerate}
\item \textbf{Faux.} \esp{La deforestación es mala para los animales porque pierden su casa y su comida.} (La déforestation est mauvaise pour les animaux car ils perdent leur habitat et leur nourriture.)
\item \textbf{Vrai.} \esp{Las inundaciones destruyen las casas y los cultivos, por ejemplo en Douala después de las grandes lluvias.}
\item \textbf{Faux.} \esp{El lago Chad es cada vez más pequeño porque llueve poco y hace mucho calor.} (Le lac Tchad devient de plus en plus petit, il ne grandit pas.)
\item \textbf{Vrai.} \esp{Reciclar los residuos reduce la contaminación y protege el medioambiente.}
\end{enumerate}
\emph{Point de méthode :} la justification doit être une phrase complète en espagnol, avec sujet et verbe conjugué ; une seule phrase suffit.
\end{corrige}

\begin{corrige}[Exercice 3 — Le bon mot : phrases complètes en espagnol]
\begin{enumerate}
\item \esp{En Douala, después de las grandes lluvias, hay a menudo una \textbf{inundación} en algunos barrios.}
\item \esp{La \textbf{deforestación} de la cuenca del Congo amenaza a los gorilas y a los elefantes.}
\item \esp{En el Extremo Norte de Camerún, la \textbf{sequía} hace la agricultura muy difícil.}
\item \esp{La \textbf{contaminación} del aire viene de los coches, las fábricas y los incendios de sabana.}
\item \esp{Para proteger el planeta, debemos \textbf{reciclar} el plástico y el papel.} (\esp{reciclar} reste à l'infinitif après \esp{debemos}.)
\end{enumerate}
\end{corrige}

\begin{corrige}[Exercice 4 — Conjugaison : le futur simple]
\begin{enumerate}
\item \esp{Mañana, nosotros \textbf{plantaremos} diez árboles en el patio del liceo.} — futur régulier : infinitif + \esp{-emos}.
\item \esp{Si no llueve, los campesinos no \textbf{podrán} cultivar el maíz.} — \esp{poder} est irrégulier au futur : radical \esp{podr-} (podré, podrás, podrá, podremos, podréis, podrán).
\item \esp{¿Tú \textbf{vendrás} a la manifestación por el clima el sábado?} — \esp{venir} est irrégulier : radical \esp{vendr-}. N'oublie pas le point d'interrogation ouvrant ¿.
\item \esp{El próximo año \textbf{habrá} menos plástico en nuestra ciudad.} — \esp{haber} est irrégulier : \esp{habrá} (il y aura), accentué sur la dernière syllabe.
\item \esp{Los científicos \textbf{dirán} la verdad sobre el clima.} — \esp{decir} est irrégulier : radical \esp{dir-} (diré, dirás, dirá, diremos, diréis, dirán).
\end{enumerate}
\emph{Rappel :} au futur simple, on ajoute les terminaisons \esp{-é, -ás, -á, -emos, -éis, -án} à l'infinitif entier pour les verbes réguliers ; les verbes irréguliers changent de radical mais gardent les mêmes terminaisons.
\end{corrige}

\begin{corrige}[Exercice 5 — Thème : traduis en espagnol]
\begin{enumerate}
\item \esp{Si no protegemos el bosque, desaparecerá.} — « protégeons » = présent de l'indicatif en français (après \og si \fg{}), et c'est aussi le présent de l'indicatif en espagnol : \esp{protegemos}. « Disparaîtra » = futur : \esp{desaparecerá}.
\item \esp{Habrá más sequías y los animales sufrirán.} — « il y aura » = \esp{habrá} (irrégulier) ; « souffriront » = \esp{sufrirán}.
\item \esp{Plantaremos árboles y reciclaremos los residuos.} — deux futurs réguliers à la première personne du pluriel : \esp{-emos}.
\item \esp{Si apagas la luz, ahorrarás energía.} — « tu éteins » = \esp{apagas} (présent après \esp{si}) ; « tu économiseras » = \esp{ahorrarás} (futur). On peut aussi dire \esp{consumirás menos energía}.
\end{enumerate}
\emph{Point de vigilance :} ne jamais mettre le futur après \esp{si} : on dit \esp{si protegemos} et non \esp{si protegeremos}.
\end{corrige}

\begin{corrige}[Exercice 6 — Transforme les phrases (condition + conséquence)]
\begin{enumerate}
\item \esp{Si cortamos los árboles, los animales perderán su casa.} — \esp{perderán} : futur régulier de \esp{perder}.
\item \esp{Si usas el coche todos los días, contaminarás el aire.} — \esp{contaminarás} : futur, deuxième personne du singulier.
\item \esp{Si reciclamos el plástico, protegeremos el medioambiente.} — \esp{protegeremos} : futur régulier, première personne du pluriel.
\item \esp{Si no llueve durante meses, habrá una sequía terrible.} — « il y aura » = \esp{habrá}, futur irrégulier de \esp{haber}.
\end{enumerate}
\emph{Méthode :} la première phrase devient la condition (verbe au présent après \esp{si}) ; la deuxième devient la conséquence (verbe au futur simple).
\end{corrige}

\begin{corrige}[Exercice 7 — Texte à trous]
Texte complété :
\begin{quote}
\esp{El \textbf{calentamiento global} es un problema muy grave para nuestro planeta: las temperaturas suben cada año. En el norte de Camerún, la \textbf{sequía} destruye las cosechas porque no llueve durante muchos meses. En cambio, en las ciudades del sur, una \textbf{inundación} puede cubrir las calles de agua después de las lluvias. La \textbf{deforestación} de los bosques mata a muchos animales. Para luchar contra estos problemas, debemos usar \textbf{energías renovables} como el sol y el viento, y también debemos \textbf{reciclar} nuestros residuos todos los días.}
\end{quote}
\emph{Indices :} « las temperaturas suben » appelle \esp{calentamiento global} ; « no llueve durante muchos meses » appelle \esp{sequía} ; « las calles de agua » appelle \esp{inundación} ; « el sol y el viento » sont des \esp{energías renovables} ; après \esp{debemos}, il faut l'infinitif \esp{reciclar}.
\end{corrige}

\begin{corrige}[Exercice 8 — Appariements]
\begin{center}
\begin{tabularx}{\linewidth}{|X|X|}
\hline
\rowcolor{popblueL} \textbf{Español} & \textbf{Français} \\
\hline
1. \esp{la sequía} & c. la sécheresse \\
\hline
2. \esp{la inundación} & d. l'inondation \\
\hline
3. \esp{la deforestación} & e. la déforestation \\
\hline
4. \esp{las energías renovables} & a. les énergies renouvelables \\
\hline
5. \esp{reciclar} & b. recycler \\
\hline
\end{tabularx}
\end{center}
\emph{Attention :} \esp{sequía} prend un accent sur le \esp{i} (\esp{se-quí-a}) ; \esp{reciclar} est un verbe du premier groupe, régulier.
\end{corrige}

\begin{corrige}[Exercice 9 — Compréhension écrite : la forêt amazonienne]
\textbf{I. Vrai ou faux ?}
\begin{enumerate}
\item \textbf{Faux.} Citation : \esp{« La Amazonía es el bosque más grande del mundo. »} C'est le plus grand bosque du monde, pas un petit bosque d'Europe.
\item \textbf{Faux.} Citation : \esp{« Los hombres cortan miles de árboles para cultivar la tierra, criar vacas o buscar oro. »} Ils coupent les arbres pour cultiver, élever des vaches ou chercher de l'or, pas pour construire des écoles.
\item \textbf{Vrai.} Citation : \esp{« Algunos gobiernos y organizaciones ya trabajan para salvar el bosque. »}
\end{enumerate}
\textbf{II. Questions ouvertes (réponses modèles).}
\begin{enumerate}
\item \esp{La Amazonía está en peligro porque los hombres cortan miles de árboles cada año para cultivar la tierra, criar vacas y buscar oro. Además, la deforestación contamina el aire y cambia el clima.}
\item \esp{Si no protegemos la Amazonía, desaparecerán muchas especies de plantas y de animales, y habrá más sequías y más incendios en la región.}
\item \esp{Todos podemos consumir menos carne y menos papel, reciclar y usar energías renovables.}
\end{enumerate}
\textbf{III. Questions de langue.}
\begin{enumerate}
\item Deux verbes au futur simple : \esp{« desaparecerán muchas especies de plantas y de animales »} (\esp{desaparecerán}, futur de \esp{desaparecer}) et \esp{« habrá más sequías y más incendios »} (\esp{habrá}, futur irrégulier de \esp{haber}). On trouve aussi \esp{reduciremos} et \esp{protegeremos}.
\item Exemple : \esp{« Si consumimos menos carne y menos papel, reduciremos la presión sobre el bosque. »} Construction : \esp{si} + verbe au \textbf{présent de l'indicatif} (\esp{consumimos}) dans la proposition de condition, et verbe au \textbf{futur simple} (\esp{reduciremos}) dans la proposition principale qui exprime la conséquence.
\end{enumerate}
\emph{Barème indicatif (sur 20) :} I. 6 pts (1 pt par réponse + 1 pt par citation exacte) ; II. 6 pts (2 pts par question) ; III. 8 pts (4 pts par question : 2 pts pour la phrase recopiée, 2 pts pour l'explication).
\emph{Points de vigilance :} recopier la citation exactement, avec les accents ; répondre aux questions ouvertes par des phrases complètes, pas par un seul mot ; ne pas confondre \esp{criar} (élever) et \esp{crear} (créer).
\end{corrige}

\begin{corrige}[Exercice 10 — Thème et version]
\textbf{A. Thème — proposition de traduction :}\par
\vspace{0.3em}
\esp{« Si no protegemos el bosque, desaparecerá. Habrá más sequías y los animales sufrirán. Plantaremos árboles y reciclaremos los residuos. »}\par
\vspace{0.3em}
\emph{Justifications :} après \esp{si}, présent de l'indicatif \esp{protegemos} (jamais de futur) ; « elle disparaîtra » = \esp{desaparecerá} ; « il y aura » = \esp{habrá} (futur irrégulier de \esp{haber}) ; « souffriront » = \esp{sufrirán} ; « nous planterons » = \esp{plantaremos} ; « nous recyclerons » = \esp{reciclaremos}. « Les déchets » se traduit par \esp{los residuos} ou \esp{los desechos}.\par
\vspace{0.5em}
\textbf{B. Version — texte original et traduction :}\par
\vspace{0.3em}
\begin{quote}
\esp{El lago Chad, en el norte de Camerún, se vuelve cada vez más pequeño. Los científicos dicen que, si no llueve más, el lago desaparecerá casi totalmente. Los pescadores no podrán trabajar y las familias sufrirán. Por eso, los países de la región plantarán árboles y protegerán el agua del lago.}
\end{quote}
\vspace{0.3em}
\emph{Traduction :} « Le lac Tchad, au nord du Cameroun, devient de plus en plus petit. Les scientifiques disent que, s'il ne pleut pas davantage, le lac disparaîtra presque totalement. Les pêcheurs ne pourront pas travailler et les familles souffriront. C'est pourquoi les pays de la région planteront des arbres et protégeront l'eau du lac. »\par
\vspace{0.3em}
\emph{Justifications :} \esp{se vuelve cada vez más pequeño} = « devient de plus en plus petit » ; \esp{si no llueve más} = « s'il ne pleut pas davantage » (présent après \esp{si}) ; \esp{desaparecerá} = « disparaîtra » ; \esp{no podrán trabajar} = « ne pourront pas travailler » (futur irrégulier de \esp{poder}) ; \esp{por eso} = « c'est pourquoi » ; \esp{los pescadores} = « les pêcheurs ».\par
\vspace{0.3em}
\emph{Barème indicatif (sur 20) :} thème 10 pts (2 pts par phrase : temps corrects 1 pt, vocabulaire et orthographe 1 pt) ; version 10 pts (fidélité du sens 6 pts, qualité du français 4 pts).
\emph{Points de vigilance :} ne pas écrire \esp{si no protegeremos} (futur interdit après \esp{si}) ; ne pas oublier l'accent de \esp{habrá}, \esp{desaparecerá}, \esp{podrán} ; \esp{el lago Chad} se traduit « le lac Tchad ».
\end{corrige}

\begin{corrige}[Exercice 11 — Production écrite type Bac]
\textbf{Proposition de rédaction modèle (environ 100 mots) :}
\begin{quote}
\esp{El medioambiente está en peligro, pero mi familia y yo queremos actuar. Primero, reciclaremos el plástico, el papel y el vidrio. Si reciclamos nuestros residuos, reduciremos la contaminación de nuestra ciudad. También plantaremos árboles en nuestro barrio para luchar contra la deforestación. Mi padre usará menos el coche y caminaremos más. Si apagamos la luz y los aparatos eléctricos, ahorraremos mucha energía. En el futuro, mi familia utilizará energías renovables como el sol. Estoy seguro de que, si todos hacemos un esfuerzo, protegeremos nuestro planeta para los jóvenes de mañana.}
\end{quote}
\emph{Traduction du modèle :} « L'environnement est en danger, mais ma famille et moi voulons agir. D'abord, nous recyclerons le plastique, le papier et le verre. Si nous recyclons nos déchets, nous réduirons la pollution de notre ville. Nous planterons aussi des arbres dans notre quartier pour lutter contre la déforestation. Mon père utilisera moins la voiture et nous marcherons davantage. Si nous éteignons la lumière et les appareils électriques, nous économiserons beaucoup d'énergie. À l'avenir, ma famille utilisera des énergies renouvelables comme le soleil. Je suis sûr que, si nous faisons tous un effort, nous protégerons notre planète pour les jeunes de demain. »\par
\vspace{0.3em}
\textbf{Vérification des consignes dans le modèle :}
\begin{itemize}
\item lexique de la leçon : \esp{reciclar}, \esp{contaminación}, \esp{deforestación}, \esp{energías renovables}, \esp{medioambiente} (5 mots, minimum 3 exigés) ;
\item futurs simples : \esp{reciclaremos}, \esp{reduciremos}, \esp{plantaremos}, \esp{usará}, \esp{caminaremos}, \esp{ahorraremos}, \esp{utilizará}, \esp{protegeremos} ;
\item phrases avec \esp{si} : \esp{Si reciclamos\ldots reduciremos} ; \esp{Si apagamos\ldots ahorraremos} ; \esp{si todos hacemos\ldots protegeremos}.
\end{itemize}
\emph{Barème indicatif (sur 20) :} respect des consignes et de la longueur 4 pts ; correction grammaticale (futurs, construction avec \esp{si}) 6 pts ; richesse du lexique 4 pts ; orthographe et accents 3 pts ; organisation et cohérence du texte 3 pts.
\emph{Points de vigilance :} compter ses mots (80 à 100) ; vérifier l'accent de chaque futur (\esp{-é}, \esp{-ás}, \esp{-á}, \esp{-emos}, \esp{-án}) ; ne jamais mettre le futur après \esp{si} ; accorder les adjectifs (\esp{energías renovables}, \esp{mucha energía}) ; éviter le calque « \esp{realizar un esfuerzo} » (\esp{realizar} = accomplir, pas « se rendre compte ») ; préférer \esp{hacer un esfuerzo}.
\end{corrige}

\newpage

\coursec{Fiche bilan}

\begin{popfigure}
\begin{center}
\begin{tikzpicture}[mindmap, grow cyclic, every node/.style={concept, align=center, font=\small},
  root concept/.append style={concept color=popgreen, font=\small\bfseries, text=white},
  level 1/.append style={level distance=3.4cm, sibling angle=51.4},
  level 1 concept/.append style={font=\footnotesize\bfseries, text=white},
  level 2/.append style={level distance=2.1cm, sibling angle=45},
  level 2 concept/.append style={font=\tiny, text=popdark}]
\node[root concept]{El medioambiente y el cambio climático}
  child[concept color=popblue]{ node{Lexique du climat}
    child[concept color=popblue!50]{ node{calentamiento global, sequía} }
    child[concept color=popblue!50]{ node{contaminación, inundación} }
  }
  child[concept color=poporange]{ node{Causes}
    child[concept color=poporange!50]{ node{deforestación, gases} }
  }
  child[concept color=poppurple]{ node{Conséquences}
    child[concept color=poppurple!50]{ node{lac Tchad, forêt du Congo} }
  }
  child[concept color=popgold]{ node{Futur simple}
    child[concept color=popgold!50]{ node{infinitif + é, ás, á, emos, éis, án} }
  }
  child[concept color=poppink]{ node{Si + présent, futur}
    child[concept color=poppink!50]{ node{Si reciclamos, salvaremos...} }
  }
  child[concept color=popgreen!70!black]{ node{Solutions}
    child[concept color=popgreen!50]{ node{reciclar, energías renovables} }
  }
  child[concept color=popmuted]{ node{Méthode}
    child[concept color=popmuted!50]{ node{commentaire : lire, relever, analyser} }
  };
\end{tikzpicture}
\end{center}
\legende{Carte mentale de la leçon : environnement et changement climatique}
\end{popfigure}

\begin{bilan}
\textbf{1. Lexique clé à connaître par cœur}
\begin{itemize}
  \item \esp{el medioambiente} : l'environnement ; \esp{el cambio climático} : le changement climatique
  \item \esp{el calentamiento global} : le réchauffement climatique ; \esp{la contaminación} : la pollution
  \item \esp{la sequía} : la sécheresse ; \esp{la inundación} : l'inondation
  \item \esp{la deforestación} : la déforestation ; \esp{los incendios} : les incendies
  \item \esp{las energías renovables} : les énergies renouvelables ; \esp{reciclar} : recycler
  \item \esp{proteger} : protéger ; \esp{contaminar} : polluer ; \esp{ahorrar agua/energía} : économiser l'eau/l'énergie
  \item \esp{el planeta está en peligro} : la planète est en danger
\end{itemize}

\textbf{2. Grammaire : le futur simple (futuro simple)}

\begin{center}
\begin{tabularx}{\linewidth}{|l|X|X|}\hline
\rowcolor{popblueL} Personne & Terminaison & Exemple (hablar) \\ \hline
yo & -é & hablar\cle{é} \\ \hline
tú & -ás & hablar\cle{ás} \\ \hline
él/ella/usted & -á & hablar\cle{á} \\ \hline
nosotros & -emos & hablar\cle{emos} \\ \hline
vosotros & -éis & hablar\cle{éis} \\ \hline
ellos/ustedes & -án & hablar\cle{án} \\ \hline
\end{tabularx}
\end{center}

\begin{itemize}
  \item Formation : \cle{infinitif entier} + terminaisons (identiques pour -ar, -er, -ir).
  \item Radicaux irréguliers fréquents : \esp{tener $\to$ tendr-}, \esp{hacer $\to$ har-}, \esp{poder $\to$ podr-}, \esp{decir $\to$ dir-}, \esp{salir $\to$ saldr-}, \esp{haber $\to$ habr-} (il y aura : \esp{habrá}).
  \item Emploi : actions futures, prévisions (\esp{El clima cambiará.} « Le climat changera. »).
\end{itemize}

\textbf{3. La conséquence : si + présent, futur}

\begin{center}
\begin{tabularx}{\linewidth}{|X|X|}\hline
\rowcolor{popgreenL} Estructura & Ejemplo \\ \hline
Si + présent, futur & \esp{Si reciclamos, salvaremos el planeta.} « Si nous recyclons, nous sauverons la planète. » \\ \hline
Futur + si + présent & \esp{Habrá sequías si no llueve.} « Il y aura des sécheresses s'il ne pleut pas. » \\ \hline
\end{tabularx}
\end{center}

\begin{itemize}
  \item Jamais de futur après \esp{si} : on dit \esp{si llueve}, jamais « si lloverá ».
  \item Équivalent français : « si + présent, futur » — structure identique, aucun piège majeur.
\end{itemize}

\textbf{4. Méthode express : commenter un texte sur l'environnement}
\begin{enumerate}
  \item Lire deux fois ; repérer le thème, la source, la thèse de l'auteur.
  \item Souligner le lexique du climat et les chiffres ou exemples (lac Tchad, forêt du Congo).
  \item Relever les procédés : futur de prévision, phrases avec \esp{si}, champs lexicaux.
  \item Rédiger : présentation du texte, résumé des idées, analyse, avis personnel argumenté (\esp{En mi opinión...}, \esp{Pienso que...}).
\end{enumerate}

\textbf{5. Gestes pour protéger l'environnement (production)}
\begin{itemize}
  \item \esp{Debemos reciclar la basura.} « Nous devons recycler les déchets. »
  \item \esp{Hay que plantar árboles y ahorrar agua.} « Il faut planter des arbres et économiser l'eau. »
  \item \esp{Si usamos energías renovables, reduciremos la contaminación.} « Si nous utilisons les énergies renouvelables, nous réduirons la pollution. »
\end{itemize}
\end{bilan}

\begin{aretenir}[Checklist avant le Bac]
\begin{itemize}
  \item $\square$ Je connais le lexique du climat : \esp{calentamiento global, sequía, inundación, deforestación, contaminación}.
  \item $\square$ Je sais traduire \esp{energías renovables} et \esp{reciclar} sans hésiter.
  \item $\square$ Je forme le futur simple : infinitif + \esp{é, ás, á, emos, éis, án}.
  \item $\square$ Je connais les radicaux irréguliers : \esp{tendré, haré, podré, diré, saldré, habrá}.
  \item $\square$ Je n'oublie pas les accents du futur (\esp{hablaré}, \esp{hablará}, \esp{hablarán}).
  \item $\square$ Je construis la conséquence : \esp{si} + présent, futur — jamais de futur après \esp{si}.
  \item $\square$ Je sais citer des exemples africains : le lac Tchad, la forêt du bassin du Congo.
  \item $\square$ Je sais proposer des solutions : \esp{reciclar, plantar árboles, ahorrar agua, usar energías renovables}.
  \item $\square$ Je respecte l'accord de l'adjectif : \esp{la sequía larga}, \esp{las inundaciones graves}.
  \item $\square$ Je distingue \esp{tiempo} (temps/météo) et \esp{clima} (climat) — faux ami avec « temps ».
  \item $\square$ Je sais structurer un commentaire : présentation, résumé, analyse, avis personnel.
  \item $\square$ Je vérifie la ponctuation espagnole : ¿ ... ? et ¡ ... ! dans mes productions.
\end{itemize}
\end{aretenir}

\findecours

\end{document}
